%% Template for the submission to:
%%   The Annals of Statistics [AOS]
%%
%%%%%%%%%%%%%%%%%%%%%%%%%%%%%%%%%%%%%%%%%%%%%%
%% In this template, the places where you   %%
%% need to fill in your information are     %%
%% indicated by '???'.                      %%
%%                                          %%
%% Please do not use \input{...} to include %%
%% other tex files. Submit your LaTeX       %%
%% manuscript as one .tex document.         %%
%%%%%%%%%%%%%%%%%%%%%%%%%%%%%%%%%%%%%%%%%%%%%%

\documentclass[aos]{imsart}

%% Packages
\RequirePackage{amsthm,amsmath,amsfonts,amssymb}
\RequirePackage[numbers]{natbib}
%\RequirePackage[authoryear]{natbib}%% uncomment this for author-year citations
\RequirePackage[colorlinks,citecolor=blue,urlcolor=blue]{hyperref}%% uncomment this for coloring bibliography citations and linked URLs
\RequirePackage{graphicx}%% uncomment this for including figures

\startlocaldefs
%%%%%%%%%%%%%%%%%%%%%%%%%%%%%%%%%%%%%%%%%%%%%%
%%                                          %%
%% Uncomment next line to change            %%
%% the type of equation numbering           %%
%%                                          %%
%%%%%%%%%%%%%%%%%%%%%%%%%%%%%%%%%%%%%%%%%%%%%%
%\numberwithin{equation}{section}
%%%%%%%%%%%%%%%%%%%%%%%%%%%%%%%%%%%%%%%%%%%%%%
%%                                          %%
%% For Axiom, Claim, Corollary, Hypothesis, %%
%% Lemma, Theorem, Proposition              %%
%% use \theoremstyle{plain}                 %%
%%                                          %%
%%%%%%%%%%%%%%%%%%%%%%%%%%%%%%%%%%%%%%%%%%%%%%
\theoremstyle{plain}
\newtheorem{thm}{\textsc{Theorem}}[section]
\newtheorem{prop}[thm]{\textsc{Proposition}}
\newtheorem{cor}[thm]{\textsc{Corollary}}
\newtheorem{lem}[thm]{\textsc{Lemma}}
%%%%%%%%%%%%%%%%%%%%%%%%%%%%%%%%%%%%%%%%%%%%%%
%%                                          %%
%% For Assumption, Definition, Example,     %%
%% Notation, Property, Remark, Fact         %%
%% use \theoremstyle{remark}                %%
%%                                          %%
%%%%%%%%%%%%%%%%%%%%%%%%%%%%%%%%%%%%%%%%%%%%%%
%\theoremstyle{remark}
%\newtheorem{???}{???}
%\newtheorem*{???}{???}
%\newtheorem{???}{???}[???]
%\newtheorem{???}[???]{???}
%%%%%%%%%%%%%%%%%%%%%%%%%%%%%%%%%%%%%%%%%%%%%%
%% Please put your definitions here:        %%
%%%%%%%%%%%%%%%%%%%%%%%%%%%%%%%%%%%%%%%%%%%%%%
\newcommand{\david}[1]{{\small{\color{blue} \sf $\heartsuit$ David: #1}}}
\newcommand{\piotr}[1]{{\small{\color{red} \sf $\clubsuit$ Piotr: #1}}}
\newcommand{\paul}[1]{{\small{\color{green} \sf $\spadesuit$ Paul: #1}}}
\newcommand{\E}{\mathbb{E}}
\newcommand{\N}{\mathbb{N}}
\newcommand{\R}{\mathbb{R}}
\newcommand{\I}{\mathbb{I}}
\newcommand{\M}{\mathcal{M}}
\newcommand{\T}{\mathcal{T}}
\newcommand{\bs}{\boldsymbol}
\endlocaldefs

\begin{document}

\begin{frontmatter}
%%%%%%%%%%%%%%%%%%%%%%%%%%%%%%%%%%%%%%%%%%%%%%
%%                                          %%
%% Enter the title of your article here     %%
%%                                          %%
%%%%%%%%%%%%%%%%%%%%%%%%%%%%%%%%%%%%%%%%%%%%%%
\title{Limits on variable selection and external information}
%\title{A sample article title with some additional note\thanksref{T1}}
\runtitle{???}
%\thankstext{T1}{A sample of additional note to the title.}

\begin{aug}
%%%%%%%%%%%%%%%%%%%%%%%%%%%%%%%%%%%%%%%%%%%%%%%
%% Only one address is permitted per author. %%
%% Only division, organization and e-mail is %%
%% included in the address.                  %%
%% Additional information can be included in %%
%% the Acknowledgments section if necessary. %%
%% ORCID can be inserted by command:         %%
%% \orcid{0000-0000-0000-0000}               %%
%%%%%%%%%%%%%%%%%%%%%%%%%%%%%%%%%%%%%%%%%%%%%%%
\author[A]{\fnms{???}~\snm{???}\ead[label=e1]{???@???}},
\author[B]{\fnms{???}~\snm{???}\ead[label=e2]{???@???}}
\and
\author[B]{\fnms{???}~\snm{???}\ead[label=e3]{???@???}}
%%%%%%%%%%%%%%%%%%%%%%%%%%%%%%%%%%%%%%%%%%%%%%
%% Addresses                                %%
%%%%%%%%%%%%%%%%%%%%%%%%%%%%%%%%%%%%%%%%%%%%%%
\address[A]{???\printead[presep={,\ }]{e1}}

\address[B]{???\printead[presep={,\ }]{e2,e3}}
\end{aug}

\begin{abstract}
Lorem ipsum dolor sit amet, consectetur adipiscing elit. Sed auctor enim metus, vel commodo neque laoreet ut. Praesent ac maximus mauris. Aliquam at lorem vitae nisi porta viverra. Integer dictum at ante at lobortis. Nulla placerat tortor eros, non condimentum elit molestie a. Aliquam orci est, congue vel hendrerit sed, scelerisque in ante. Duis mattis id urna nec tempus. Pellentesque efficitur auctor ex a egestas. Cras euismod tincidunt nisi eget egestas. Vivamus purus augue, semper vitae iaculis eu, egestas sit amet odio. Interdum et malesuada fames ac ante ipsum primis in faucibus.

Pellentesque nisi nunc, commodo id orci ac, imperdiet ullamcorper ante. Donec iaculis consequat interdum. Vivamus at efficitur est. Phasellus vitae nisi ac ipsum pellentesque faucibus in quis odio. Nunc vehicula nulla ac egestas posuere. Morbi sit amet neque feugiat, semper neque quis, finibus quam. Sed eget hendrerit mi. Nam fermentum tempus vulputate. Pellentesque convallis risus tellus, eget viverra elit egestas et. Ut felis turpis, pharetra ac convallis quis, facilisis quis urna. Quisque quis ipsum sit amet odio imperdiet blandit ut non ex.

Vestibulum rhoncus pretium ex, egestas facilisis neque efficitur quis. Morbi malesuada volutpat erat, sed aliquet lectus dictum ac. Aenean ligula massa, consequat a blandit id, placerat.
\end{abstract}

\begin{keyword}[class=MSC]
\kwd[Primary ]{???}
\kwd{???}
\kwd[; secondary ]{???}
\end{keyword}

\begin{keyword}
\kwd{???}
\kwd{???}
\end{keyword}

\end{frontmatter}
%%%%%%%%%%%%%%%%%%%%%%%%%%%%%%%%%%%%%%%%%%%%%%
%% Please use \tableofcontents for articles %%
%% with 50 pages and more                   %%
%%%%%%%%%%%%%%%%%%%%%%%%%%%%%%%%%%%%%%%%%%%%%%
%\tableofcontents

%%%%%%%%%%%%%%%%%%%%%%%%%%%%%%%%%%%%%%%%%%%%%%
%%%% Main text entry area:

\section{Introduction} 

In many modern applications, the sample size is insufficient to reliably do inference. This limitation has motivated the development of high-dimensional techniques that frequently rely on strict assumptions on sparsity and signal strength. Yet, often, external information is available and can be leveraged to enhance inference.
%Our motivation stems from the observation that, often, 
In particular, external information can guide our decisions regarding which parameters should be included in a given model. In regression, this information may originate from the inherent nature of each variable (e.g., clinical history versus genomic markers, sociodemographics versus job history) or external data that distinguishes relevant variables from irrelevant ones (e.g., functional annotations on genes). 
Our paper investigates this concept in the the Gaussian Sequence Model (GSM) and linear regression, where external information helps partition variables into groups with potentially distinct characteristics.

Incorporating external information is often advocated within the framework of data integration, as joint learning can lead to more accurate inference compared to analyzing datasets separately. Indeed, numerous applied works have employed external information to guide inference. For instance, \cite{stingo:2011, quintana, cassese:2014} proposed Bayesian variable selection methods for gene expression, where prior probabilities for non-zero coefficients depend on biological knowledge and meta-covariates. Additionally, \cite{chen_tinghuei:2021} predicted disease outcomes by allowing LASSO penalties to depend on functional annotation categories. Beyond regression, node and edge covariates have been incorporated in \cite{peterson} and \cite{jewson2023graphical} to drive edge inclusion in Gaussian graphical models, while \cite{schiavon2022generalized} used meta-covariates to determine non-zero loadings in factor models. In causal analysis, covariates of interest (referred to as treatments) guide the inclusion of other covariates based on their degree of association with the treatments \cite{belloni2014inference,antonelli2021bayesian}. Collectively, empirical evidence consistently demonstrates improved structural learning when external data is suitably integrated. However, despite this empirical success, a theoretical framework explaining precisely why this occurs remains elusive.

Here we show that incorporating external information allows pushing the theoretical limits under which consistent variable selection is possible, understood as recovering the true model with probability going to 1 as $n$ and the number of variables $p$ grow. Previous literature has extensively explored how high-dimensional variable selection is constrained by inherent characteristics of the dataset, such as $n$, $p$, the signal strength, the correlation between variables and the number of variables truly associated to the outcome \cite{buhlmannvandegeerbook,overviewvarsel,handbookbayesvarsel,wainwrighthighdimbook}. In this paper we analyse how conditions for consistency can be relaxed in the presence of external information. We focus on variable selection but using external information is also expected to benefit parameter estimation and prediction problems. 

To make our ideas concrete, consider the variable selection problem in the linear model with Gaussian noise:
\begin{equation}\label{eq:linearmodel}
	\bs y \;=\; \bs X \bs \beta^{\star} + \bs\epsilon,\qquad \bs \epsilon \sim N(0,\sigma^2 I_n),
\end{equation}
where $\bs X\in \R^{n\times p}$, $\bs y\in \R^n$, and $\sigma >0$. A large class of methods operate by penalizing the size of an estimated $\hat{\bs  \beta}$. For instance, penalized likelihood methods are based on optimizing the log-likelihood plus a penalty term driven by the $\ell_q$ "norm" of $\hat{\bs  \beta}$ for $q\in [0,1]$ or folded concave penalties \citep{lasso,overviewvarsel,bertsimas}. 
In Bayesian settings, variable selection is often based on posterior model probabilities that are directly connected to $\ell_0$ penalties \citep{BIC,EBIC,rossell2022concentration}. 

We study the case where the external information partitions the $p$ variables into $b$ blocks denoted $B_j$, $j=1,\ldots,b$ across which key characteristics such as the level of sparsity or signal strength may vary. We introduce block $\ell_0$ penalties that allow modulating the strength of the penalty in each block. Unlike standard $\ell_0$ penalties, like BIC \cite{BIC}, RIC \cite{RIC} or EBIC \cite{EBIC}, block penalties are non-exchangeable in the sense that the complexity penalty for adding a variable may depend on its block $B_j$.

\paul{Reorganized the intro to have this contribution paragraph.}
Our main contributions are as follows. We show that block-based non-exchangeable $\ell_0$ penalties soften the theoretical conditions for consistent model recovery in the Gaussian sequence model and in linear regression.
Specifically, that it is possible for an oracle to set block-based penalties so that variable selection consistency is either attained where otherwise not possible or at a faster rate. Another important contribution is that we introduce data-based variable selection procedures, based on empirical Bayes, that realize the theoretical benefits without requiring an oracle. Other more minor contributions include a tight necessary conditions on signal strength in the GSM and linear regression as well as consistency results for the EBIC in a wider class of asymptotic regimes than proved in \cite{EBIC} and \cite{ebicdiverging}.

Our work has connections with ideas from multiple hypothesis testing. In this line of research, \cite{genovese2006} proposed a false discovery rate (FDR) procedure in which p-values are weighted based on prior information. They show that if the weights are positively associated to the null hypotheses being false, their procedure improves power. Subsequent works discussed oracle choices of weights either based on external information or derived from the data \cite{roederwasserman,basu}. 
Adaptive FDR methods \cite{blanchard} estimate the
proportion of truly null hypotheses to improve statistical power.
Recently, \cite{ramdas2019prior} combined external information and adaptivity in a group-based FDR control. 
Relative to this work, we study variable selection consistency in high-dimensional regression as well as the associated conditions on sparsity and signal strength.

Also related to our work, \cite{scarlett2012compressed} studied compressed sensing when prior information allows splitting parameters into blocks, where one knows the true proportion of non-zero parameters in each block. They show that sparse signal recovery requires smaller $n$ than with exchangeable penalties. Key differences with our work are their assuming independence across covariates, which renders the results inapplicable to regression. Also, they do not consider the sequence model, nor that the proportions of non-zero parameters are unknown in practice.

The paper is organized as follows. Section~\ref{sec:intro_blockpen} introduces block $\ell_0$ penalization. Section~\ref{sec:properties_orth} presents its
variable selection properties and its benefits in the Gaussian sequence model, in an oracle setting where the sparsity and betamin conditions are known for all groups.
Those results apply to essentially any penalized likelihood or Bayesian method, including $\ell_0$ and $\ell_1$ penalties, because in that setting selection operates by thresholding \cite{davidomiros,rigollethighdimensional}. Section~\ref{sec:properties_legression} studies block $\ell_0$ penalties in linear regression, and shows analogous benefits to those in Section~\ref{sec:properties_orth}. These results can be extended to a wide class of Bayesian variable selection methods. Section~\ref{sec:informed_l0} presents data-based 
procedures that achieve faster convergence without requiring an oracle. Section~\ref{sec:simulation} shows empirical examples and Section~\ref{sec:discussion} concludes. Proofs are gathered at the end of the paper.
\medskip

\noindent \textbf{Notation:}  
We denote by $\bs \beta\in \R^p$ the parameters of interest and by $\bs \beta^*$ their true values. Let  $V=\{1,\ldots,p\}$. For any $A\subseteq V$ and any vector $\bs x\in \R^p$, $\bs x_A$ denotes the subvector of $\bs x$ with entries corresponding to indices in $A$. For any matrix $\bs X \in \R^{n\times p}$, $\bs X_A$ denotes the submatrix of $\bs X$ obtained by selecting the \emph{columns} with indices in $A$. We denote by
$S=\left\{i\in V:\,\,  \beta^{\star}_i \neq 0 \right\}$ the true support of $\bs \beta^{\star}$, its size is $s=|S|$, and hence $p-s$ is the number of truly inactive parameters.
Denote $\beta_{\min}^{\star}=\min_{i \in S}|\beta^{\star}_i|$ the smallest true signal.
We assume that $V$ is partitioned into a fixed number $b$ of disjoint blocks $B_j \subseteq V$ for $j=1,\ldots,b$.
We denote by $S_j=\left\{i\in B_j:\,\,  \beta^{\star}_i \neq 0\right\}$ the set of active parameters in block $B_j$,  $s_j=|S_j|$ its size, $p_j-s_j$ the number of inactive parameters in the block,
and $\beta_{\min,j}^{\star}=\min_{i \in S_j}|\beta^{\star}_i|$.
Given sequences $f(n) > 0$ and $g(n) >0$, $f(n)=O(g(n))$ means that there exists a constant $c<\infty$ such that $f(n) \leq c g(n)$ for all $n \geq n_0$ and some fixed $n_0$, and $f(n)=\Theta(g(n))$ means that $f(n) = O(g(n))$ and $g(n) = O(f(n))$. 
$f(n)=o(g(n))$ means that $f(n) / g(n) \to 0$, and $f(n)=\Omega(g(n))$ means that $\lim_{n\to\infty}f(n) / g(n) > 0$. 
\medskip

\noindent\textbf{Our asymptotic regime:} Although not denoted explicitly, $s_j$, $p_j$ and $\beta_{\min,j}^{\star}$ are  functions of $n$, and so are $s$, $p$ and $\beta_{\min}^{\star}$. We study asymptotic regimes where \\[5pt]
     (A1)\quad $n \,\to \, \infty$, and the number of blocks $b$ is constant. \\[5pt]
     (A2)\quad For all $j$, $p_j-s_j \,\to \, \infty$.\\[5pt]
Assumption A2 implies that $p-s$ goes to infinity but we do not make any assumption on the asymptotic regime linking $n$ and $p$. We provide results for both cases $s_j\geq1$ fixed or diverging.
\david{But surely we will end up needing some conditions involving the $s_j$'s, at least in regression. And on how $p$ can grow with $n$. To be cautious, I wouldn't write "we do not impose any restriction...", it may raise eyebrows and possibly disappoint later if we end up setting conditions (even if indirectly) on $(p,n)$ or $(s,s_j)$. For example, you could just say "Assumption A2 is a mild assumption requiring that one has a growing number of inactive covariates in each block."}\paul{I changed it}

\section{Variable selection via informed block penalization}\label{sec:intro_blockpen}

Consider a set of candidates models $\M$ given by subsets $M\subseteq V$ and their corresponding coordinate subspaces
$$
\mathcal{L}_M\;:=\;\{\beta\in \R^p: \beta_i=0\mbox{ if }i\notin M\}.
$$ \paul{deleted a sentence in which we were redefining S as the smallest subset such that $\bs\beta^\star\in L_S$}
Consider a standard $\ell_0$ selection procedure with penalty $\eta(M)$ that depends on $M$ linearly through its cardinality, $\eta(M)=\kappa |M|$ for some $\kappa>0$. The selected model is
\begin{equation}\label{eq:stdl0}
\hat{S}= \arg\max_{M\in \M} \left\{\max_{\beta\in \mathcal{L}_M}\ell(\bs y; \bs \beta)\;-\;\kappa |M|\right\},
\end{equation}
where $\ell(\bs y; \bs \beta)$ is the log-likelihood function. 
We study a block-informed procedure that penalizes differently each block $B_1,\ldots,B_b$,
\begin{equation}\label{eq:Mhat}
  \hat S^b\;\in\;\arg\max_{M\in \M} \left\{\max_{\beta\in \mathcal{L}_M}\ell(\bs y; \bs \beta)-\sum_{j=1}^b\kappa_j |M_j| \right\},
\end{equation}
where $\kappa_1>0,\ldots,\kappa_b>0$. We denote a model by $M=M_1\cup\cdots\cup M_b$, and $M_j\subseteq B_j$ represents the  selected parameters in block $j$.
Note that $\hat{S}^1$ is the standard $\ell_0$ selector in \eqref{eq:stdl0}.
\piotr{Section 2 is half-page. Section 3 takes 10 pages. Is there a way to reorganize sectioning so that things are lighter and more balanced? (this is not critical)}\paul{I don't see how we could reorganize that unless we merge section 2 and the introduction}
\section{Gaussian sequence model}\label{sec:properties_orth}

We now introduce the Gaussian sequence model, its connection to orthogonal linear regression and exhibit how $\hat S^b$ performs block-wise hard thresholding in that setting. We also provide its variable selection properties, discuss its benefits in terms of lessening the conditions for consistency variable selection and convergence rates, and illustrate said advantages under different asymptotic regimes.

In this section, in addition to A1--A2, we assume that \\[5pt]
(A3)\quad $\text{for all } j, \;\; s_j \,\to \, \infty$.\\[5pt]
Assumption A3 is made to simplify the exposition, and because we derive a new necessary condition on signal strength under A3, but we obtained analogous results when the $s_j$'s are finite (Supplement XXXX).

\subsection{Gaussian sequence model and thresholding}\label{sec:GSMgen}

The Gaussian sequence model assumes:
\begin{equation}\label{eq:GSM}
	\bs y \;=\; \sqrt{n} \bs \beta^{\star} + \bs\epsilon,\qquad \bs \epsilon \sim N(0,\sigma^2 I_p),
\end{equation}
where $\bs y \in \mathbb{R}^p$ and without loss of generality we set $\sigma^2=1$ to streamline notation. It is a popular simplified model for high-dimensional inference, and is also central to non-parametric statistics. See \cite{johnstone}, Chapter 3 for a complete introduction.

Let $\tilde{\bs \beta}=\frac{1}{\sqrt{n}}\bs y \sim N(\bs \beta^*,\tfrac{1}{n}I_p)$ be the MLE under \eqref{eq:GSM}. The next proposition shows that, by orthogonality, the block-wise $\ell_0$ penalty reduces to hard-thresholding $\tilde{\bs \beta}$, with block-dependent thresholds.
\begin{prop}\label{prop:isthresh}
	In the Gaussian sequence model \eqref{eq:GSM}, $\hat{S}^b$ defined in \eqref{eq:Mhat} satisfies $\hat S^b=\hat S^b_1\cup \cdots \cup \hat S^b_b$, where, for each $j=1,\ldots,b$,
	$$
	\hat S_j^b\;:=\;\left\{i\in B_j:\; |\tilde{\beta}_i|> \sqrt{\frac{2\kappa_j}{n}}\right\}.
	$$ 
\end{prop}
Most popular variable selection procedures also take this block-based thresholding form in \eqref{eq:GSM} with non-exchangeable penalization. It is the case of the LASSO and adaptive LASSO when letting the penalty vary by block: see Lemma \ref{lem:l1isthreshold} in Auxiliary results in the Supplement. It is also the case of Bayesian procedures under most standard priors \citep{davidomiros} where one sets a different prior inclusion probability in each block. Thus, equivalently, we study generic hard thresholding model selectors of the form
\begin{equation}\label{eq:thresholdS}
\hat S^{b}\;:=\;\{i\in B_j: |\tilde{\beta}_i|>\tau_j\}, \qquad\bs \tau=(\tau_1,\ldots,\tau_b)\in \R^b_{>0}.	
\end{equation}

Results in this section also generalize to orthogonal linear regression with normalized columns, i.e. when $\bs X$ in \eqref{eq:linearmodel} satisfies $\bs X^\top \bs X=n I_p$. In that setting the MLE is also distributed $N(\bs \beta^*,\tfrac{1}{n}I_p)$ and by orthogonality, variable selection with $\hat{S}^b$ also amounts to block thresholding.

\subsection{Selection based on block thresholds}
\label{ssec:block_thresholds}


We now study the statistical performance of the block-based thresholding operator in \eqref{eq:thresholdS}. Block thresholding was previously studied in the context of parameter estimation for wavelet-based models with equally sized blocks; see \cite{johnstone}, Chapters 7 to 9.
We focus here on its variable selection properties. We successively analyze properties relative to false positives, false negatives, exact recovery and rates of convergence. 
\medskip

\noindent\textbf{False positives.} 
A false positive
occurs when there exists $j=1,\ldots,b$ and $i\in B_j$ such that $\beta^\star_i=0$
but $|\tilde \beta_i|>\tau_j$. In other words, $\hat S^{b}\not\subseteq S$. A sufficient assumption to guarantee that asymptotically $\hat S^{b}\subseteq S$ is:\\[5pt]
(A4)\quad for all $j=1,\ldots,b$ and all sufficiently large $n$,\;\; $\tau_j \;\geq\;  \sqrt{\frac{{2\ln(p_j-s_j)}}{{n}}}$.
\\[5pt]
The next proposition shows that
A4 is also essentially necessary.
\begin{prop}\label{prop:suffnecIblockthresh}In the
sequence model \eqref{eq:GSM}, assume A1 and A2.
\begin{longlist}
    \item If A4 holds, then
$P(\hat{S}^{b}\subseteq S) \to 1$.\\
\item If for some $j\in\{1,\ldots,b\}$\; $
\lim_{n\to \infty} \frac{\sqrt{n}\tau_j}{\sqrt{2\ln(p_j-s_j)}} \;<\;1$,
then $P(\hat{S}^{b}\subseteq S ) \not \to 1$.
\end{longlist}
\end{prop}
By Proposition~\ref{prop:suffnecIblockthresh},
a necessary assumption to have, asymptotically, no false positives
is: 
\begin{equation}\label{cond:necIallblockthresh}
\lim_{n\to \infty} \frac{\tau_j}{\sqrt{\tfrac{2\ln(p_j-s_j)}{n}}}\;\geq\;1\qquad\mbox{for all }j=1,\ldots,b.	
\end{equation}
Proposition~\ref{prop:suffnecIblockthresh} extends to $\hat{S}^{b}$ results known for $\hat{S}$ in orthogonal linear regression \cite{ slope, buhlmannvandegeerbook, wainwrighthighdimbook} and in the Gaussian sequence model, \cite{ariascastro} and \cite{johnstone}, Chapter 3.
\bigskip

\noindent\textbf{False negatives.} 
False negatives occur when $S\not\subseteq \hat S^{b}$, i.e. for some $j \in \{1,\ldots,b\}$ and $i\in B_j$,  we have that $\beta^\star_i\neq 0$ and $|\tilde\beta_i|\leq \tau_j$.
Control of false negatives in a block $B_j$ requires bounding the minimum of $|S_j|=s_j$ independent Gaussians with non-zero means $\beta^{\star}_i$, $i\in S_j$. This can be achieved by setting a betamin condition
on $\beta_{\min,j}^{\star}\;:=\;\min\{|\beta^\star_i|: i\in S_j\}$,
and a assumption on the thresholds $\tau_j$ for $j=1,\ldots,b$.  
Intuitively, to have no false positives (also known as screening property), each $\tau_j$ needs to be smaller than $\beta^*_{\min,j}$.
A closely-related \paul{I don't get this "closely-related"} sufficient assumption is: \\[5pt]
(A5)\quad For all $j=1,\ldots,b$ and sufficiently large $n$, \;
	 $\beta_{\min,j}^{\star} - \tau_j \;\geq\; \sqrt{\frac{2\ln(s_j)}{n}}$.

\begin{prop}\label{prop:suffnecIIblockthresh}
 In the sequence model \eqref{eq:GSM}, assume A1--A3 and A5, then
$P(S\subseteq \hat{S}^{b} ) \to 1$.
\end{prop}
Assumption A5 for standard thresholding was proved to be sufficient for selection consistency in orthogonal linear regression in \cite{wainwrighthighdimbook}. Proposition~\ref{prop:suffnecIIblockthresh} generalizes it to block thresholds. The study of {\it necessary} assumptions
is more complicated. We provide below a first general assumption valid for any vector $\bs \beta^*$ and a second tight assumption for a worst case where $\beta_i^*=\beta_{\min}^*$ for all $i\in S$.
\begin{prop}\label{prop:necIIblockthresh}In the
sequence model \eqref{eq:GSM}, assume A1 and A2.
\begin{longlist}
    \item Suppose that $\tau_j<\beta^\star_{\min,j}$ satisfies \eqref{cond:necIallblockthresh} for some $j \in \{1,\ldots,b\}$. If \begin{equation}\label{cond:necIIallblockthresh2}
    \lim_{n\to \infty} \sqrt{n}(\beta_{\min,j}^{\star}-\tau_j)\; <\; \infty
\end{equation}
then $P(S\subseteq \hat{S}^{b} ) \not \to 1$.
\item Suppose there exists $j\in \{1,\ldots,b\}$ such that  $\beta^{\star}_i=\beta_{\min,j}^{\star}$ for all $i \in S_j$, $\tau_j<\beta^\star_{\min,j}$ satisfies \eqref{cond:necIallblockthresh}, and  $s_j/p_j<1$. If A3 holds and, in addition, 
\begin{equation}\label{cond:necIIallblockthresh}
\lim_{n\to \infty} 	\frac{\beta_{\min,j}^{\star}-\tau_j}{\sqrt{\frac{\pi}{2}\frac{\ln(s_j)}{n}}} \;\leq\;  1\end{equation}
then $P(S\subseteq \hat{S}^{b} ) \not \to 1$.
\end{longlist}\end{prop}
\paul{The AOS style requires using "longlist" instead of "enumerate" which makes it a bit ugly and not very readable. Maybe we could have the equations inline? As I formated it in Proposition \ref{prop:suffnecIIblockthresh}}
\cite{abraham2023sharp} showed the necessity of \eqref{cond:necIIallblockthresh2} for consistency in the single-block case
in more restrictive asymptotic regimes satisfying $p/s \to \infty$.
Proposition~\ref{prop:necIIblockthresh} (ii) gives a new precise necessary  rate of growth of $\sqrt{n}\left(\beta_{\min}^{\star}-\tau\right)$ for a worst case in which all signals are equal. It shows that in that case sufficient assumption~A5 is essentially necessary up to a constant factor, and that A5 cannot be relaxed in general.
\bigskip

\noindent\textbf{Exact model recovery guarantees.} 
Combining assumption A4 (Proposition~\ref{prop:suffnecIblockthresh}) and A5 (Proposition~\ref{prop:suffnecIIblockthresh}),
we asymptotically recover the correct support if for all $j=1,\ldots,b$,
\begin{equation*}
  \sqrt{\frac{2\ln(p_j-s_j)}{n}}\;\;\leq\;\;\tau_j\;\;\leq\;\;\beta_{\min,j}^{\star}-\sqrt{\frac{2\ln(s_j)}{n}}
\end{equation*}
for sufficiently large $n$. This implies
that for all $j=1,\ldots,b$,
\begin{equation}\label{eq:betamin}
    {\beta_{\min}^{\star}}_j \geq \sqrt{\frac{2\ln(p_j-s_j)}{n}} + \sqrt{\frac{2\ln(s_j)}{n}}
\end{equation}
By Propositions~\ref{prop:suffnecIblockthresh} and~\ref{prop:suffnecIIblockthresh} signals as small as the right-hand side in \eqref{eq:betamin} can be recovered in block $B_j$. Moreover, Propositions~\ref{prop:suffnecIblockthresh} (ii) and \ref{prop:necIIblockthresh} show that assumptions A4--A5 are essentially necessary. It follows that betamin condition \eqref{eq:betamin} is near-necessary for selection consistency in the following sense.
\begin{lem}\label{lem:necbetaminblockthresh}
    If there exists $j\in\{1,\ldots,b\}$ such that $\beta^{\star}_i=\beta_{\min,j}^{\star}$ for all $i \in S_j$, $s_j/p_j<1$ and
    \begin{equation}\label{eq:necbetaminblockthresh}
      \lim_{n\to \infty}\frac{\beta_{\min,j}^*}{\sqrt{\tfrac{2\ln(p_j-s_j)}{n}}+\sqrt{\tfrac{\pi}{2}\tfrac{\ln(s_j)}{n}}}\;<\; 1  
    \end{equation}
then under assumptions A1--A3,\;  $P(\hat{S}^{b} = S)\not \to 1$
\end{lem}



\bigskip

\noindent\textbf{Rate of convergence.} From the
preceding results, we can bound the rate of convergence of $P(\hat S^{b}\neq S)$.\piotr{Here we do not assume neither A1 nor A2? also I slightly edited this result so check if you agree}\paul{We don't need A1 and A2 to state the bounds no.}
\begin{thm}\label{theo:fullrecovGS}
Assume A4 and A5. Then, for sufficiently large $n$,
\begin{equation} \label{eq:rateconvgsm}
P(\hat{S}^{b} \neq S)\;\leq\;\sum_{j=1}^b\,\,\frac{e^{-\tfrac{n}{2}\big[\tau_j^2-\tfrac{2\ln(p_j-s_j)}{n}\big]}}{\sqrt{\pi\ln(p_j-s_j)}}+\frac{e^{-\tfrac{n}{2}\big[(\beta^\star_{\min,j}-\tau_j)^2-\tfrac{2\ln(s_j)}{n}\big]}}{\sqrt{\pi\ln(s_j)}}.
\end{equation}
On the other hand, if $\beta^\star_{\min,j}$ satisfy \eqref{eq:betamin} for every block $j$, and the thresholds $\tau_j$ take oracle values
	\begin{equation} \label{eq:optimtauj}
\tau^\star_j \;=\; \frac{1}{2}\beta^\star_{\min,j}+\frac{1}{n\beta^\star_{\min,j}}(\ln({p_j-s_j})-\ln({s_j})),
\end{equation}
then A4 and A5 hold and, for sufficiently large $n$,
\begin{equation}\label{eq:optimrateconvgsm}
    P(\hat{S}^{b} \neq S) \; \leq \; 2 \sum_{j=1}^be^{-[\frac{n}{8}{\beta^\star_{\min,j}}^2-\max\{\ln(p_j-s_j),\ln(s_j)\}]}
\end{equation}
\end{thm}
In \eqref{eq:rateconvgsm} the first component of each term in the sum decreases exponentially in $n\tau^2/2$, while the second component decreases exponentially in $n(\beta_{\min,j}^*-\tau)^2/2$. The choice $\tau^*_j$ ensures that both terms are equal and hence approximates the values minimizing \eqref{eq:rateconvgsm}. We call this ideal $\tau^*_j$ oracle threshold because it depends on the quantities $s_j$ and $\beta_{\min,j}$ unknown in practice.  
The resulting bound in \eqref{eq:optimrateconvgsm} is tightest in the worst-case scenario where $\beta_i^*=\beta_{\min}^*$ for all $i \in S$, and approximates the fastest rate achievable in that worst case.
\subsection{Benefits of block penalties}\label{sec:prosconsblock}


We now examine the benefits brought by block penalization. We start with a simple example to build intuition and then discuss
two types of benefit: those relative to conditions for variable selection consistency and those relative to rates of convergence.
\bigskip

Intuitively, block penalties allow adapting 
to varying sparsity $s_j$  or signal strengths $\beta^\star_{\min,j}$ across blocks. 
When $p_j-s_j$ is small one may use smaller thresholds $\tau_j$ and still control false positives, and when $\beta^\star_{\min,j}$ is large one may use larger $\tau_j$ and still control false negatives.
As the blocks become more discriminative, that is, more distinct $p_j-s_j$ and $\beta^\star_{\min,j}$ across blocks, the larger the potential benefits.
Consider a somewhat extreme example where $n = 1$, $p = 100$, $s = 10$ and $\beta_i^\star=4$ for all $i\in S$. There are $b=2$ blocks: a large sparse block with $s_1=0$, $p_1-s_1=80$ and a smaller non-sparse block with $s_2=10$, $p_2-s_2=10$. In the left panel of Figure~\ref{fig:starter}, a standard threshold selector is used with $\tau=\sqrt{2 \ln(p-s)/ n}$, the minimal threshold such that the FWER vanishes (dashed line). The large number of truly zero parameters $p-s$ forces a high threshold, which leads to a false negative. In the right panel, block thresholds are set to $\tau_j=\sqrt{2 \ln(p_j-s_j)/ n}$.
Adjusting $\tau_2$ to the lower number of truly zero means in block 2 ($p_2-s_2$) allows recovering the true model. 
\begin{figure}[t]
\includegraphics[scale=.6]{images/starter.example.pdf}
\centering
\caption{In both panels, the 90 white dots are generated from $N(0, 1)$ and the 10 black dots are generated from $N(4, 1)$. Left panel: the dashed line is $\sqrt{2 \ln (p-s)/n}$, the minimal threshold such that the FWER vanishes. Right panel: observations are grouped in two blocks. The dashed lines are the minimal block thresholds such the FWER vanishes}
\david{In the figure x-axis, I would delete "observation" and "observations grouped". Then may add "global thresholding" and "block-based thresholding" on top of the figure. Also, journals usually require that each panel is uploaded as a separate pdf. Then use tabular to create the two panels, as you would with a table.}
\label{fig:starter}
\end{figure}
\bigskip

\noindent\textbf{Conditions for variable selection consistency.} In the example, the possibility to lower thresholds in the less sparse blocks with $\hat S^2$ allowed the recovery of the smallest active signal and the true model. More precisely, assumptions A4--A5 give ranges of threshold values that are sufficient and essentially necessary for asymptotic support recovery. For the standard threshold selector $\hat{S}$, the range for the single threshold $\tau$ is
\begin{equation}\label{eq:rangethreshex}
  \sqrt{\frac{2\ln(p-s)}{n}}\;\;\leq\;\;\tau\;\;\leq\;\;\beta^{\star}_{\min}-\sqrt{\frac{2\ln(s)}{n}}.	
\end{equation}
For a block threshold selector $\hat{S}^{b}$, assuming without loss of generality that the global minimum $\beta_{\min}^*$ is in block $b$, the ranges for the $\tau_j$'s are
    \begin{eqnarray}\label{eq:rangeblockthreshex}
  \sqrt{\frac{2\ln(p_1-s_1)}{n}}\;\;&\leq\;\;\tau_1\;\;\leq&\;\;\beta^{\star}_{\min,1}-\sqrt{\frac{2\ln(s_1)}{n}}\nonumber\\
  &\;\;\vdots \\
  \sqrt{\frac{2\ln(p_b-s_b)}{n}}\;\;&\leq\;\;\tau_b\;\;\leq&\;\;\beta^{\star}_{\min}-\sqrt{\frac{2\ln(s_b)}{n}} \nonumber.       
    \end{eqnarray}

By definition, the ranges in \eqref{eq:rangeblockthreshex} are wider than the range in \eqref{eq:rangethreshex}, making conditions to get selection consistency milder with $\hat S^b$. If for some $\tau_k$ and $\tau_l$ the ranges do not overlap, which, for example, occurs if $\beta_{\min}^{\star} - \sqrt{(2/n) \ln(s_b)} < \sqrt{(2/n) \ln(p_1-s_1)}$,
then a global threshold $\tau$ cannot possibly satisfy conditions in 
\eqref{eq:rangeblockthreshex} and consistent model recovery is essentially not possible.
We next give precise conditions in which consistent model recovery is not possible with $\hat{S}$ but is with $\hat{S}^b$, that follow from Propositions~\ref{prop:suffnecIblockthresh} and \ref{prop:necIIblockthresh}.

\begin{cor}\label{cor:selconstimpossible}
Assume A1--A5, that $s/p<1$, and $\beta^{\star}_i=\beta_{\min}^{\star}$ for all $i \in S$. If
\begin{eqnarray*}
    &\lim_{n\to \infty}\frac{\beta_{\min}^*}{\sqrt{\tfrac{2\ln(p-s)}{n}}+\sqrt{\tfrac{\pi}{2}\tfrac{\ln(s)}{n}}}\;<\; 1 
\end{eqnarray*}
then 
$P(\hat{S} = S)\not \to 1$ and $P(\hat{S}^{b} = S)\to 1$.
\end{cor}

% for $\hat S^b$ and $\hat S$ depends on the chosen thresholds and decrease as thresholds shrink. However, thresholds also need to be large enough to avoid false positives. To make betamin conditions comparable and understand the complete potential gains with $\hat S^b$,  $\tau_j=\sqrt{2\ln(p_j-s_j)/n}$ for all $j$ for $\hat S^b$, $\tau=\sqrt{2\ln(p-s)/n}$ for $\hat S$. This choice allows to precisely assess what signals can or cannot be recovered with the two selectors. 

% For $\hat S$, the betamin condition A5 constrains only the global $\beta_{\min}^{\star}$ which solely determines whether signal strength is large enough for variable selection consistency. To under , we compare next the constraints that selection consistency with $\hat S$ and $\hat S^b$ impose on that $\beta_{\min}^{\star}=\beta_{\min,b}^{\star}$. For simplicity we assume A5 is met for $\hat S^b$ in blocks other than $b$. 



% \paul{You had asked the question below. I have added more explanation below. The idea of this section is to analyze the smallest signal recoverable, the constraints on the overall $\beta_{\min}^{\star}$}
% \david{Why is it enough to look only at block $b$ ? I mean, even if $\beta_{\min}^\star$ is in block $b$, $(p_j,s_j)$ from other blocks could be such that \eqref{eq:betamincon.suffblockthresh} holds in block $b$, but not for some other block $j$ no?}\piotr{Yes, this should be clarified.}

% From \eqref{eq:betamin}, we obtain the sufficient betamin condition for $\hat{S}$:
% \begin{equation}\label{eq:betamincon.suffthresh}
%     \beta_{\min}^{\star} \;>\; \sqrt{\frac{2\ln(p-s)}{n}}+\sqrt{\frac{2\ln(s)}{n}}
% \end{equation} 
% For the block threshold selector $\hat{S}^{b}$, since $\beta_{\min}^{\star}$ is located in block $b$, assumption \eqref{eq:betamin} requires:

% \begin{equation}\label{eq:betamincon.suffblockthresh}
%     \beta_{\min}^{\star} \;>\; \sqrt{\frac{2\ln(p_b-s_b)}{n}}+\sqrt{\frac{2\ln(s_b)}{n}}
% \end{equation} 
% We then set: 
% \begin{equation}\label{eq:betaminbardef}
%     \underline{\beta}_{\min}^{*} := \sqrt{\frac{2\ln(p-s)}{n}} + \sqrt{\frac{2\ln(s)}{n}},\qquad
%     \underline{\beta}^{*}_{\min,b} := \sqrt{\frac{2\ln(p_b-s_b)}{n}} + \sqrt{\frac{2\ln(s_b)}{n}}
% \end{equation} Since $p_b-s_b \leq p - s$ and $s_b \leq s$, $\underline{\beta}^{*}_{\min,b} \leq \underline{\beta}_{\min}^{*}$ and assumption \eqref{eq:betamincon.suffblockthresh} is always milder than \eqref{eq:betamincon.suffthresh}. Selector $\hat{S}^{b}$ is then able to recover smaller signals than $\hat{S}$. The magnitude of the difference depends on how $p-s$ and $s$ compare to $p_b-s_b$ and $s_b$. The gain thus depends on whether the block structure induced by the external information is indicative of varying sparsity. \david{I didn't find sub-section "Minimal betamin conditions" super-informative (given what we already said earlier). This is however used later for Figure Figure~\ref{fig:ratiobetamincond}, so maybe we can introduce $\underline{\beta}_{\min}^{*}$ later, when presenting that example.}\piotr{agreed}

\bigskip

\noindent\textbf{Rates of convergence.}   We just saw that block penalties can attain consistent variable selection where standard penalties cannot. We now discuss differences in the probability of correct variable selection, when consistent selection is possible for both procedures. We assume the block thresholds 
are set to their oracle values $\tau_j^\star$ 
and similarly the single threshold
takes its oracle value $\tau^\star$.
Let ${OR}^b_{orth}$ be the oracle convergence rate for $\tau_j^\star$
and ${OR}_{orth}$ that for $\tau^\star$, as given in \eqref{eq:optimrateconvgsm}.
Then
\begin{equation}\label{eq:ratioCR}
    \frac{{OR}^b_{orth}}{{OR}_{orth}}\;=\;\sum_{j=1}^be^{-\frac{n}{8}\left({\beta^\star_{\min,j}}^2-{\beta_{\min}^\star}^2\right)-(\max\{\ln(p-s),\ln(s)\}-\max\{\ln(p_j-s_j),\ln(s_j)\}\big)}
\end{equation}
Since $\beta_{\min,j}^\star \geq \beta_{\min}^\star$ and $\max\{\ln(p-s),\ln(s)\}\geq\max\{\ln(p_j-s_j),\ln(s_j)\}\big)$, equation \eqref{eq:ratioCR} shows that for $n$ large enough ${OR}^b_{orth} \leq {OR}_{orth}$ for any partition in blocks, i.e. the convergence rate for $\tau_j^\star$ is never worse than for $\tau^\star$. The magnitude of the gain depends on the discriminative power of the blocks.
For any sparse setting where $s_j<p_j-s_j$ for every $j$\paul{You added "midly" sparse, I removed it. I does not matter if it is highly or midly sparse. By sparse I just mean $s_j<p_j-s_j$ in every block to know what is $\max\{\ln(p_j-s_j),\ln(s_j)\}$ and simplify notation}, assuming without loss of generality that $\beta_{\min,b}^\star=\beta_{\min}^\star$, we have
\begin{equation*}
    \frac{{OR}^b_{orth}}{{OR}_{orth}}\;=\; \frac{p_b-s_b}{p-s} \,+\,\sum_{j<b}\frac{p_j-s_j}{p-s}e^{-\frac{n}{2}\left({\beta^\star_{\min,j}}^2-{\beta_{\min}^\star}^2\right)}
\end{equation*}

There are then two sources of improvement in convergence rates: in any block $j$, a smaller number of truly zero means, that is $p_j-s_j< p-s$, or a larger smallest signal, $\beta^\star_{\min,j} > \beta_{\min}^\star$. Depending on the sparsity of the block, one of the two effects usually prevails. In blocks where $\beta^\star_{\min,j} \approx \beta_{\min}^\star$, improvement comes mainly from having a smaller number of truly zero means. In those blocks
one sets $\tau^\star_j < \tau^\star$, so that one may detect smaller signals. 
If
$\beta^\star_{\min,j} \gg \beta_{\min}^\star$,
one sets $\tau^\star_j > \tau^\star$
and the probability of false positives is reduced.


On the contrary, there is no improvement 
(${OR}^b_{orth}={OR}_{orth}$), 
when for example, as before $p_j-s_j>s_j$, but $p_j-s_j=(p-s)/b$ 
 and $\beta^\star_{\min,j}=\beta_{\min}^\star$ for all $j$. In such case, the external information fails to identify blocks that are substantially differently sparse. 

\subsection{Examples with two blocks}\label{sec:examples}

We illustrate our results with concrete examples contemplating different sparsity regimes and informativeness of the blocks, as summarized in Table~\ref{tab:examples}. We focus on a sparse setting where $\ln(s)=o(\ln(p-s))$, with two blocks ($b=2$). We assume that $\beta_{\min}^{\star}$, the smallest non-zero $|\beta_i^*|$ is located in block $2$. 

\begin{table*}[ht]
\caption{Sparsity assumptions in examples}
\label{tab:examples}
\begin{tabular}{@{}cccccccc@{}}
\hline \\[-1em]
Example & $p-s$ & $p_1-s_1$ & $p_2-s_2$ & $s_1$ 
& $s_2$\\[3pt]
\hline\\[-1em]
$1$    & $\frac{3}{2}\,n$ & $\frac{3}{2}\,n-\sqrt{n}$ & $\sqrt{n}$  & $e^{\sqrt{\ln(n)}}/2$  & $e^{\sqrt{\ln(n)}}/2$ \\[3pt]
$2$    & $e^{n/10}$ & $e^{n/10}-n^2$ & $n^2$  & $e^{\sqrt{\ln(n)}}/2$  & $e^{\sqrt{\ln(n)}}/2$   \\[3pt]
$3$    & $n$ & $n-\ln(n)$ & $\ln(n)$  & $e^{\sqrt{\ln(n)}}/2$  & $e^{\sqrt{\ln(n)}}/2$   \\[3pt]
$4$    & $n$ & $n/2$ & $n/2$  & $e^{\sqrt{\ln(n)}}/2$  & $e^{\sqrt{\ln(n)}}/2$ \\[3pt]
\hline
\end{tabular}
\end{table*}

In Examples 1 to 3, external information is discriminative as it singles out a block $B_2$ such that $\bs\beta^*_{B_2}$ is significantly less sparse than the global $\bs\beta^*$. In Example 1, external information is moderately discriminative --- $p-s/p_2-s_2$ is a power of $n$, in Example 2 it is highly discriminative --- $p-s/p_2-s_2$ is exponential in $n$ and in Example 3 is highly discriminative in the sense that $\bs\beta^*$ is sparse but $\bs\beta^*_{B_2}$ is non-sparse ($\ln(p_2-s_2)=o(\ln(s_2))$). In Example 4, external information is non-discriminative, the block partition is an average random guess where each block has half of the inactive variables. Examples 3 and 4 differ from 1 and 2 in that they are set in only very midly sparse settings. Note also that in Example 2, by Corollary~\ref{cor:selconstimpossible}, recovery is asymptotically not possible with $\hat{S}$ since $\sqrt{\frac{2\ln(p-s)}{n}}+\sqrt{\frac{\ln(s)}{n}} \to \sqrt{2/10} \approx 0.44$ and $\beta_{\min}=1/10$, but it is with $\hat S^b$. 
\paul{It is not super important but I think that what matters in terms of usefulness of the external information is $p-s/p_2-s_2$ but David has been stressing $p_1-s_1/p_2-s_2$ in his edits. My point of view is that the main difference between $\hat S$ and $\hat S^b$ is that you can recover smaller signals in the blocks that are less sparse than the global set of variables. If you have a three block situation such: $p-s=e^n$, $p_1-s_1=e^n-n^2-n$, $p_2-s_2=n^2$ and $p_3-s_3=n$, $p_2-s_2/p_3-s_3$ is not telling of the gain with $\hat S^b$.}

To showcase the benefit in terms of conditions for selection consistency, we start by plotting the range of valid thresholds for $\hat{S}$ and $\hat{S}^{b}$ as given in assumptions A4 and A5. We then compare the \textit{smallest signal recoverable} by $\hat S$ and $\hat S^b$ while still being selection consistent. We denote those $\underline{\beta}_{\min}^{*}$ and $\underline{\beta}_{\min,b}^{*}$ respectively. In betamin condition A5, we set the threshold(s) to the lowest value allowing the vanishing of the FWER as per Proposition~\ref{prop:suffnecIblockthresh}. Since we assumed the the global minimum $\beta_{\min}^*$ is in block $b$, we get:
\begin{equation}\label{eq:betaminbardef}
    \underline{\beta}_{\min}^{*} := \sqrt{\frac{2\ln(p-s)}{n}} + \sqrt{\frac{2\ln(s)}{n}},\qquad
    \underline{\beta}^{*}_{\min,b} := \sqrt{\frac{2\ln(p_b-s_b)}{n}} + \sqrt{\frac{2\ln(s_b)}{n}}
\end{equation} 
and discuss the ratio $\underline{\beta}^{*}_{\min,b}/\underline{\beta}_{\min}^{*}$ in our examples. We then illustrate the benefit of block penalties in terms of rate of convergence with a plot of the ratio ${OR}^b_{orth}/{OR}_{orth}$.

\bigskip
\noindent\textbf{Valid thresholds for $\hat S$ and $\hat S^b$.}
Figure \ref{fig:boundthresh} shows the upper and lower bounds for $\tau$ as in \eqref{eq:rangethreshex} and $\tau_2$ as in \eqref{eq:rangeblockthreshex} in the four examples assuming $\beta^*_{\min,1}=1/2$  and $\beta^*_{\min,2}=\beta^*_{\min}=1/10$. Observe that, for $\tau$, the sparser $\bs \beta^*$ is, the larger the $n$ needed to make the upper bound to exceed the lower bound, i.e. to essentially make support recovery possible. 
The block threshold $\tau_2$ requires smaller $n$, 
with the gap with respect to $\tau$ increasing as the difference in sparsity between block 2 and $V$
increases. The gap is close to null in Example 4 when external information is non-discriminative. In Example 2 on the contrary, recovery with $\hat{S}$ is not possible in the range of values of $n$ considered while it is with $\hat{S}^{b}$. Example 3 shows that the gain with $\hat{S}^{b}$ can be large even in setting where $V$ is only somewhat sparse.
\begin{figure}[t]
\includegraphics[scale=1]{images/boundthresh.pdf}
\centering
\caption{Bounds on $\tau$, $\tau_1$ and $\tau_2$ as per assumptions A4 and A5 in Examples 1 to 4}
 \label{fig:boundthresh}
\end{figure}


\bigskip
\noindent\textbf{Smallest signal recoverable.} Figure~\ref{fig:ratiobetamincond} plots the ratio of smallest signal recoverable $\underline{\beta}^{*}_{\min,b}/\underline{\beta}^{*}_{\min}$ in our four examples. In Example 1, with discriminative blocks, the
smallest signal recoverable with $\hat{S}^{b}$ is asymptotically about 30\% smaller than with the standard selector $\hat{S}$. 
In Example 2 where blocks are even more discriminative, the $\underline{\beta}^{*}_{\min,b}/\underline{\beta}^{*}_{\min}$ converges to 0: the smallest signal recoverable with $\hat S^b$ is infinitely smaller.
On the contrary, in Example 4,
the ratio converges to 1. With non-discriminative blocks, the benefits in terms of
signal recoverable fade away as $n$ grows. 
Example 3 illustrates how highly discriminative blocks can also bring significant benefits in terms of signal recoverable in a regime that is very mildly sparse.
\begin{figure}[t]
\includegraphics[scale=0.7]{images/betamin.bound.ratio.pdf}
\centering
\caption{Ratio of minimal betamin bounds $\underline{\beta}^{*}_{\min,b}$ and $ \underline{\beta}_{\min}^{*} $ for Example 1 to 4}
\label{fig:ratiobetamincond}
\end{figure}

\bigskip
\noindent\textbf{Oracle convergence rates.} We now plot the ratio ${OR}^b_{orth}/{OR}_{orth}$ in Examples 2 to 4 assuming $\beta_{\min }^{\star}=O\Big(\sqrt{\frac{\ln(p-s)}{n}}\Big)$ for illustrative purposes. \david{Is there a reason for choosing this betamin? (before we used a fixed betamin)} \paul{If betamin is fixed the ratio drops to 0 almost immediately for sparser examples. It is a rate of increase that give visually good results for the 3 examples in a same plot.} Since the oracle rate for $\hat{S}^{b}$ depends not only on the $p_j-s_j$'s but also the $\beta_{\min,j}^*$'s, we consider two cases within each example:
one case in which signals $\beta^*_{\min,1}=\beta_{\min,2}^*$ and one case in which $\beta^*_{\min,1}=1.1 \, \beta_{\min,2}^*$. Figure~\ref{fig:ratioconvrate} shows that, when the blocks are discriminative (Examples 2-3), the oracle rate for $\hat S^{b}$ is much
than for $\hat S^{T,1}$. Particularly so when blocks are
discriminative not only in in terms of the number of truly 0 parameters, but also in terms of signal strength. In the non-discriminative Example 1 the improvement stagnates as $n$ grows.
\david{In the left panel, the ratio doesn't seem to stagnate. Also, can we use the same line types in Fig 4 as we did in Fig 3? (e.g. the solid line is Example 1 in Fig 1, Example 2 in Fig 2.). Btw, in example one would we get a ratio=1 in Fig 4? Even if obvious, maybe worth adding that line. Finally, little things: in the figure top legend I'd just write $\beta_{\min}^1=1.1\beta_{\min,2}$. And can we show normal labels for n? (100 instead of 1e+02, etc)}


\begin{figure}[t]
\includegraphics[scale=0.9]{images/comp.bound.prob.rec.optim.pdf}
\centering
\caption{Ratio of bounds on convergence rates ${OR}^b$ and ${OR}^1$ in Examples 1 to 3. In the left panels  smallest signals are equal across blocks while they differ in the right panel}
\label{fig:ratioconvrate}
\end{figure}

\david{Devoting 3 figures to these examples seemed a little too much. How about merging Figs 3-4 into a single figure, with 3 panels?}


\section{High-dimensional linear regression}
\label{sec:properties_legression}

In Section~\ref{sec:properties_orth} we showed that, in the sequence model, using block penalties allows consistent model recovery in settings where it is otherwise not possible (e.g. smaller signals), and improves oracle consistency rates. We now extend the results to linear regression. 
We present sufficient and necessary assumptions, 
betamin conditions and rates for the probability of correct selection for the block informed $\ell_0$-penalized selector $\hat S^b$ in \eqref{eq:Mhat}. We also compare the properties of $\hat S^b$ to those of the standard $\ell_0$ selector $\hat S$. 
Whereas the framework is more involved than in the sequence model and the required proof techniques are different, the results remain conceptually the same. 

%Here $V$ is the set of all variables and $p_j-s_j$ and $s_j$ are respectively the numbers of inactive and active variables in block $B_j$.
\paul{I removed a sentence that said here "Here $V$ is the set of all variables and $p_j-s_j$ and $s_j$ are respectively the numbers of inactive and active variables in block $B_j$." I don't see why we need to repeat it.} We require Assumptions A1--A2, but the number of active variables $s_j$ may either be fixed or go to infinity. 
That is, we do not assume A3 anymore. For simplicity, we assume that $\hat{S}^b$ is unique and that one constrains attention to the set of models $\M$,
such that
$\bs X_M$ has full column rank for any $M \subset \M$, and that the data-generating true support $S$ lies in $\M$. Non-uniqueness of $\hat{S}^b$ and non-full rank models can be accommodated though, at the expense of a slightly more involved treatment.

Our analysis relies on a classical connection
between $\ell_0$ penalties and Bayesian variable selection. 
Section \ref{ssec:regression_proof_strategy} reviews this connection and the proof strategy for our results.
The technical nature of the proof precludes a detailed exposition, we instead only present the most important ideas.
In Section \ref{ssec:regression_mainresult} we state our main theorem
giving variable selection consistency of $\hat{S}^{b}$ and oracle convergence rates as we did for the sequence model. In Section~\ref{ssec:neccond_linreg}, we  give necessary assumptions and betamin conditions variable selection consistency with $\hat{S}^{b}$. Section \ref{ssec:regression_gains} discusses the gains of block penalization in further detail. Section \ref{ssec:nobetaminhyp} gives a general convergence result for $\hat{S}^{b}$ under no betamin condition which is used in Section \ref{sec:informed_l0}.


\subsection{Proof strategy}
\label{ssec:regression_proof_strategy}


We state a well-known
reformulation of $\hat S^b$ in linear regression \eqref{eq:linearmodel}. 
\begin{lem}\label{lem:useful-reformulation} Let $\tilde{\bs \beta}^{(M)}= (\bs X_{\!\!M}^\top \bs X_{\!\!M})^{-1}\bs X_{\!\!M}^\top\bs y \in \R^{|M|}$ be the MLE under model $M\in \M$. Then $\hat S^b\;\in\;\arg\max_{M\in \M} {\rm C}(M)$, where $M=M_1\cup \ldots\cup M_b$ and
\begin{equation}\label{eq:CM}
    {\rm C}(M)\;: =\;\tfrac{1}{2}\|\bs X_M\tilde{\bs\beta}^{(M)}\|^2-\sum_{j=1}^b\kappa_j |M_j|,
\end{equation}
\end{lem}
By Lemma~\ref{lem:useful-reformulation}, $\hat S^b$ selects
$M\in \M$ with the largest score $C(M)$, or equivalently the largest normalized score
\begin{equation}\label{eq:normalizedscore}
    NC(M)\;:=\;\frac{e^{{\rm C}(M)}}{\underset{M' \in \M}{\sum} e^{{\rm C}(M')}}.
\end{equation}
This normalized score can be understood as a pseudo-posterior probability for model $M$ in a Bayesian variable selection framework \citep{BIC}.
Proposition 1 in \cite{rossell2022concentration}, reproduced below as Lemma \ref{lem:L0toL1convergence}, proves that the expectation of such posterior probabilities bounds the probability of an incorrect variable selection $(\hat S^b\neq S)$. 

\begin{lem}\label{lem:L0toL1convergence}
\begin{longlist}
\item 
$P(\hat{S}^{b} \neq S) \,\,\leq \,\, 2 \sum_{M \in \M\setminus \{S\}} \E\left(NC(M)\right).$
\item  For any $M,M'\in\M$, such that $M\neq M'$, 
$NC(M)\;\leq \;\big(1+e^{C(M')-C(M)}\big)^{-1}$.
\end{longlist}
\end{lem}
We use  Lemma~\ref{lem:L0toL1convergence} (i) to show the variable selection consistency of $\hat S^b$ by guaranteeing the vanishing, in expectation, of the sum of the normalized scores. Such vanishing bears resemblance to what is known as \textit{strong selection consistency} property \citep{narisettyhe2014}: the concentration of pseudo-posterior model probabilities on the truth $S$. We prove however $L_1$ convergence instead of convergence in probability.
To bound $\E(NC(M))$, by Lemma~\ref{lem:L0toL1convergence} (ii) with $M'=S$, we may instead bound a simple function of $C(S)-C(M)$: the pairwise comparison of criteria values of $S$ and $M$. 
This is achieved by noting that $C(S)-C(M)=\tfrac12 L_{SM}+\Delta_{MS}$, where for any two models $M,T\subseteq V$, we denote

\begin{equation}\label{eq:DeltaM}
\Delta_{MT}\;:=\;\sum_{j=1}^b \kappa_j(|M_j|-|T_j|). 
\end{equation}
and for any two models $Q,M\subseteq V$, we denote 
$$L_{QM}\;:=\;\|\bs X_{Q}\tilde{\bs\beta}^{(Q)}\|^2-\|\bs X_{M}\tilde{\bs\beta}^{(M)}\|^2.$$
To lower bound $C(S)-C(M)$ we use that $L_{SM}$ can be expressed in terms of chi-squared variables.
For any over-fitted model $M \supset S$, we have $\Delta_{MS}>0$ and it is well-known that $-L_{SM} \sim \chi_{|M \setminus S|}^2$. Hence, if one sets large enough $\kappa_j$ (and thus $\Delta_{MS}$), then $NC(M)$ vanishes. 
For non-overfitted  models, the idea is to take the union model $Q_S=S\cup M$, and to note that $L_{SM}=L_{Q_SM}-L_{Q_SS}$. Since $M \subset Q_S$, we may use the next lemma from \cite{rossell2022concentration} to bound $L_{Q_SM}$ and hence also $L_{SM}$.

\begin{lem}\label{lemma:s7}
    Let $M,Q$ be any two nested models such that $M \subseteq Q$. Then
$$
L_{QM}\;=\;\|\bs X_{Q}\tilde{\bs\beta}^{(Q)}\|^2-\|\bs X_{M}\tilde{\bs\beta}^{(M)}\|^2\;\sim\; \chi_{|Q\setminus M|}^2\left(\mu_{Q M}\right), 
$$
where $\chi_{k}^2(\mu)$ denotes the noncentral chi-squared distribution with $k$ degrees of freedom and noncentrality parameter $\mu$. The noncentrality parameter $\mu_{QM}$ is given by \begin{equation}\label{eq:muQM}
    \mu_{QM}\;\;:=\;\;\|\big(I_n-P_M\big) \bs X_{Q\setminus M} \bs \beta^{\star}_{Q\setminus M}\|^2
\end{equation}
where $P_M=\bs X_{M}\left(\bs X_{M}^{\top} \bs X_{M}\right)^{-1} \bs X_{M}^{\top}$.
\end{lem}

For example, by Lemma~\ref{lemma:s7}, for a under-overfitted $M\subset S$ we have $L_{SM}\sim \chi^2_{|S|-|M|}(\mu_{SM})$, which is centered around $|S|-|M|+\mu_{SM}$.
If the noncentrality parameter $\mu_{SM}$ is large enough, then $L_{SM}$ and $C(S)-C(M)$ are also large, and $NC(M)$ vanishes. Lemma \ref{lem:boundrho} shows that large $\mu_{Q_SM}$ can be achieved by setting a betamin condition and an eigenvalue condition on $\bs X$ involving the following quantity.  
 \begin{equation}\label{eq:rhoX}
        \rho(\bs X)\,=\,\min_{M \in \M:\, S \not \subseteq M} \lambda_{\min}\big(\tfrac{1}{n} \boldsymbol{X}_{S \setminus M}^{\top}\left(I_n- P_M\right) \boldsymbol{X}_{S \setminus M}\big) \geq 0.
    \end{equation}
The quantity $\rho(\bs X)$ is nonnegative and relates to how distinguishable the other models $M\in \M$ are from $S$ \citep{infotheowainwright}.  More specifically, $\tfrac{1}{n} \boldsymbol{X}_{S \setminus M}^{\top}\left(I_n- P_M\right) \boldsymbol{X}_{S \setminus M}$ is the sample covariance matrix of the residuals when regressing $\boldsymbol{X}_{S \setminus M}$ on $\boldsymbol{X}_M$.
That is, $\rho(\bs X)$ measures how well one can approximate the 
$\boldsymbol{X}_{S \setminus M}$ with $\boldsymbol{X}_M$.
In an orthonormal case where $\bs X^{\top}\bs X=n\I_p$, then $\rho(\bs X)=1$\piotr{Not $1/n$?}\paul{no, it's 1}.  

\begin{lem}\label{lem:boundrho}
 For any $M\in\M$, let $Q_S=M \cup S$. The non-centrality parameter defined in \eqref{eq:muQM} satisfies: 
    \begin{equation}\label{eq:musmbound}
    \mu_{Q_SM}\; \geq \;n\,\rho(\bs X)\,\sum_{j=1}^b|S_j\setminus M_j|\,{\beta_{\min,\,j}^{\star}}^2
\end{equation}
\end{lem}



\subsection{Selection based on block $\ell_0$ penalty}
\label{ssec:regression_mainresult}

We state two sufficient assumptions for asymptotic support recovery. 
Building on our previous discussion, the first requires the block penalties $\kappa_j$ to be large enough,  
and the second is a betamin condition:\\[5pt] 
(A6) \quad For each block $j$, there exists $f_j\to \infty$ (as $n\to \infty$) such that for $n$ sufficiently large,
\begin{equation*}
    \kappa_j \;=\;\ln(p_j-s_j) + f_j
\end{equation*} 
(A7) \quad For each block $j$, there exists $h_j\to +\infty$ such that for $n$ sufficiently large,
    \begin{equation*}
        \sqrt{\frac{(1-\gamma) n\rho(\bs X)}{6}}{{\beta_{\min}^{\star}}_j} \,- \,\sqrt{\kappa_j} \;=\; \sqrt{\ln(s_j)} + h_j .
    \end{equation*}
where $\gamma:=\frac12(1+\max_j \alpha_j)\in (\tfrac12,1)$, and
$\alpha_j:=\frac{\ln(p_j-s_j)}{\kappa_j}$.\\[5pt] 
These assumptions are similar to A4--A5 formulated for the sequence model.
 By Proposition~\ref{prop:isthresh}, in the orthonormal 
 case $\bs X^{\top}\bs X=n\I_p$, setting block penalties $\kappa_j$ 
 is equivalent to hard thresholding with block thresholds $\tau_j=\sqrt{2 \kappa_j/n}$.
 Assumptions A4--A5 can then be directly translated into assumptions on $\kappa_j$, taking $\rho(\bs X)=1$. In that case, A6 is of the same order as A4, up to a term $f_j$ that can grow at an arbitrarily slow rate. An oracle choice for $f_j$ is presented in Theorem~\ref{theo:suffcondlinearmodel}. A natural practical choice stemming from a Bayesian perspective discussed later in Section~\ref{sec:informed_l0}, is $f_j=0.5\ln(n)$ for all $j$. The betamin condition~A7 is the same as A5 up to, again, a term $h_j$ increasing at arbitrarily slow rate, and the factor $\sqrt{(1-\gamma)/3}$. This factor is smaller than $1/2$ and depends on the choice of $f_j$'s in A6. It is bounded away from 0 as long as $f_j= \Omega(\ln(p_j-s_j))$ for all $j$. This is the case for the earlier natural choice whenever $n^k=\Omega(p)$ for some $k>0$ \paul{essentially NOT ultra high dimension}. When $f_j= o(\ln(p_j-s_j))$ for some $j$, A7 requires $\sqrt{\kappa_j}+\sqrt{\ln(s_j)} = o(n\rho(\bs X)\beta_{min,j})$.
\david{Earlier we said that $f_j$ can grow at an arbitrarily slow rate, but we now see that $f_j$ has to be at least of order $p_j-s_j$. How about saying: This factor depends on the choice of $f_j$ in A6. For example, by default we take $f_j=...$, in which case $\gamma$ is bounded below by ... Perhaps also relate the choice of $(f_j,g_j)$ to the oracle $\kappa_j^\star$'s discussed later in the theorem, like "an ideal choice of $(f_j,g_j)$ leads to oracle penalties $\kappa_j^\star$ given in Theorem 4.2 below" (and, if it helps, briefly explain how the oracle $(f_j,g_j)$ depend on $s_j,p_j,\beta_{\min}^\star$)} \paul{$f_j$ can grow at any rate.  For any growth rate you will discard asymptotically the overfitting models. When $f_j= o(\ln(p_j-s_j))$ for some $j$, the betamin requirement is stronger than $\Theta(\sqrt{\kappa_j} +\; \sqrt{\ln(s_j)} + h_j)$. To see it, the betamin requirement is essentially: $\sqrt{\frac{ n\rho(\bs X)}{2}}{{\beta_{\min}^{\star}}_j}  \;=2(\max_j\frac{\ln(p_j-s_j)}{f_j}+1)(\sqrt{\kappa_j} +\; \sqrt{\ln(s_j)}) + h_j$. It is not very intuitive to have such change in betamin condition i know. It is linked to a technicality of the proof. I have tried hard to get rid of it but could not. I have explained more.\\}
Finally, we remark that in our proof one could take a betamin condition that is slightly less strict than A7, but we present A7 here to compare it more directly to A5 in the sequence model.  
\david{The rest of the paragraph appeared after Theorem 4.1, but felt more natural here, as it discussed A6--A7. I'm not sure that it's strictly needed though, it basically repeats the notion that A6--A7 are very similar to the sequence model no?} In fact, from A6--A7 we can infer the smallest signals that can be recovered with $\hat{S}^{b}$.
We get the minimal betamin condition:
\begin{equation}\label{betamincond:minimalreg} \mbox{for all }j=1,\ldots,b \qquad
    \lim_{n\to \infty}\frac{\sqrt{(1-\gamma) n\rho(\bs X)/3}\;{\beta_{\min}^{\star}}_j}{\sqrt{2\ln(p_j-s_j)} + \sqrt{2\ln(s_j)}} >1
\end{equation}
Note that, up to the factor $\sqrt{(1-\gamma)\rho(\bs X)/3}$, Expression \eqref{betamincond:minimalreg} is very similar to betamin conditions for the sequence model (Corollary \ref{cor:selconstimpossible}, Expression \eqref{eq:betamincon.suffthresh}).




We can now state our main theorem on the strong variable selection consistency of $\hat{S}^{b}$:
\begin{thm}\label{theo:suffcondlinearmodel}

Under assumptions A1, A2, A6 and, A7, we have
$$\sum_{M \in \M\setminus \{S\}}\E\left(NC(M)\right)\to 0 \qquad \mbox{and} \qquad P(\hat{S}^b = S) \to 1$$
\end{thm}


Theorem~\ref{theo:suffcondlinearmodel} considers an additive penalty across blocks
$\sum_{j=1}^b \kappa_j|M_j|$. This additivity assumption can be easily relaxed. In Supplement XXXX we give a consistency results for non-additive penalties. Those apply for example
to the EBIC penalty \citep{EBIC} and improve on known requirements on asymptotic regimes and the signal strength to have variable selection consistency \cite{ebicdiverging}.

From the proof of Theorem~\ref{theo:suffcondlinearmodel}, we can bound the convergence rate of $P(\hat S^{b}\neq S)$, given in Theorem \ref{prop:fullrecovreg}. We also give oracle block penalties that approximately optimize the bound, and provide the resulting oracle rate of convergence. 
In the statement, $\delta<1$ and $r_b>1$ should be understood as being arbitrarily close to 1.
\begin{thm}\label{prop:fullrecovreg}
Assume A6--A7. Then, for all sufficiently large $n$ and any $\delta \in (0,1)$,
\begin{equation}\label{eq:rateconvreg}
 P(\hat{S}^{b} \neq S)\;\leq\; 6(2^{2b}-2b)r_b\,\sum_{j=1}^b\,\, e^{-\tfrac{\delta}{2}\big[\kappa_j-\ln(p_j-s_j)\big]} + 
    e^{-\tfrac{\delta}{2}\big[(\sqrt{\frac{(1-\gamma) n\rho(\bs X)}{6}}{{\beta_{\min}^{\star}}_j} - \sqrt{\kappa_j})^2 -\ln(s_j) \big]}
\end{equation}
where $r_b\to 1$ as $n\to +\infty$.

Moreover, suppose that the $\beta^\star_{\min,j}$ satisfy \eqref{betamincond:minimalreg} and the $\kappa_j$ are set at the oracle values:
	\begin{equation} \label{eq:oraclepenreg}
\sqrt{\kappa^\star_j\,} \;=\; \frac12\sqrt{\frac{(1-\gamma) n\rho(\bs X)}{6}}{{\beta_{\min}^{\star}}_j}+\frac12\sqrt{\frac{6}{(1-\gamma) n\rho(\bs X)}}\frac{1}{{\beta_{\min}^{\star}}_j}\big(\ln({p_j-s_j})-\ln({s_j})\big).
\end{equation}
Then for all sufficiently large $n$, A6 and A7 hold and for any $\delta \in (0,1)$: 
\begin{equation}\label{eq:oraclerateconvreg}
    P(\hat{S}^{b} \neq S) \; \leq \; 12(2^{2b}-2b)r_b \sum_{j=1}^b e^{-\tfrac{\delta}{2}\big[\frac{(1-\gamma) n\rho(\bs X)}{24}{{\beta_{\min}^{\star}}_j}^2-\max\{\ln(p_j-s_j),\ln(s_j)\}\big]},
\end{equation}
where $r_b\to 1$ as $n\to+\infty$.
\end{thm}

\subsection{Necessary assumptions for consistent variable selection with block $\ell_0$ penalty}\label{ssec:neccond_linreg}

We derive assumptions on the $\kappa_j$ and $\beta_{min,j}$ that are necessary for consistent variable selection with $\hat{S}^b$. Recall that the strength of the block penalties $\kappa_j$ drives the vanishing of the pseudo-posterior probabilities of over-fitted models. The block penalties need then to be large enough for the pseudo-posterior probabilities of even the least penalized overfitted models to vanish. For every $j=1,\ldots,b$, denote the set of models that over-fit by only one variable from block $B_j$:
$$O_j \:=\,\{M\in \M\;|\; M=S\cup\{i\} \text{ where } i\in B_j\setminus S_j\}$$
The asymptotic recovery of $S$ implies that, $\max_{M\in O_j}NC(M)/NC(S)<1$ with probability going to 1 as $n$ grows, for every $j=1,\ldots,b$, . For every $M_k\in O_j$, the ratio $NC(M_k)/NC(S)$ grows with $L_{M_kS}$ which, by Lemma~\ref{lemma:s7}, is $\chi^2_1$ distributed. There exists then, for every $M_k\in O_j$, $Z_{M_k} \sim N(0,1)$ such that $L_{M_kS}=Z_{M_k}^2$. A necessary assumption on $\kappa_j$ follows from the properties of the maximum of the correlated $Z_M$. Denote the smallest eigenvalue of their correlation matrix:
\begin{equation}
    \underline{\lambda}_j\;:=\;\lambda_{min}(R^j)\;\;\text{where} \;\;R^j_{k,l}=\operatorname{corr}\big(Z_{M_k},Z_{M_l}\big),\;\;\forall M_k,M_l \in O_j.
\end{equation}
\paul{note it is not the correlation matrix of the $\sqrt{L_{MS}}=|Z_M|$. This result is based on \cite{hartigan} that gives a lower bound for maximum of correlated Gaussians and not maximum of absolute value of correlated Gaussians. The lower bound is probably the same but proving it would require additional lemmas.}
We denote $\underline{\lambda}$, the smallest eigenvalue of the correlation matrix of the $Z_{M_k}$, $M_k\in\cup_{j=1}^b O_j$. \paul{This quantity appears later in describing settings in which $\hat S$ is not consistent but $\hat S^b$ is.}
A lower bound on the maximum of correlated Gaussians in \cite{hartigan} yields the following result.
\begin{prop}\label{prop:nectypeIreg}
     If for some $j=1,\ldots,b$,  $\lim_{n\to\infty}\frac{\kappa_j}{\underline{\lambda}_j^2\ln(p_j-s_j)} <1$, then under assumptions A1 and A2, $$P\Big(\max_{M\in O_j}\frac{NC(M)}{NC(S)}<1\Big)\to 0 \quad\text{and}\quad P(\hat{S}^b=S) \to 0.$$
\end{prop}

\paul{Note in passing: \cite{hartigan} nicely complements the result Piotr found on independence being the "worst case" for the maximum of Gaussians. Interestingly it suggests that we could set lower penalty when inactive variables are highly correlated between each other. I remember very well that someone in Toronto remarked that it was somewhat strange that sufficient assumptions for penalty strength did not depend on the level of correlation between variables. I could not answer well back then!}

Proposition~\ref{prop:nectypeIreg} shows that 
a necessary assumption to recover $S$ asymptotically
is: 
\begin{equation}\label{cond:necIallblockpen}
\lim_{n\to\infty}\frac{\kappa_j}{\underline{\lambda}_j^2\ln(p_j-s_j)}\;\geq\;1\qquad\mbox{for all }j=1,\ldots,b.
\end{equation}
When the inactive variables are orthogonal, for every $j=1,\ldots,b$ and $M\in O_j$, the $Z_M$ are independent and $\underline{\lambda}_j=1$. Under orthogonal design, assumption~\eqref{cond:necIallblockpen} then recovers assumption~\eqref{cond:necIallblockthresh}.

Consider now an under-fitted model $M\subset S$. The recovery of $S$ implies that with probability going to 1 as $n$ grows, $NC(M)/NC(S)<1$. By definition of $NC$ the ratio grows with block penalties $\kappa_j$ and shrinks with $L_{SM}$ which is distributed $\chi^2_{|S\setminus M|}(\mu_{SM})$ by Lemma~\ref{lemma:s7}. Small values of the ratio are attained only if $\mu_{SM}$ grows fast enough compared to the $\kappa_j$. Lemma~\ref{lem:upperboundnoncentral} shows that the growth $\mu_{SM}$ is bounded by the growth of the largest active signals in $S$ that is not in $M$.
\begin{lem}\label{lem:upperboundnoncentral}
    For any $T\subseteq S$, $$\mu_{ST} \leq n\,\bar{\lambda}\,\sum_{j=1}^b|S_j\setminus T_j|\,\max_{i\in S_j\setminus T_j}{\beta_{i}^{\star}}^2 \quad \text{where}\quad \bar{\lambda}\,:=\,\lambda_{\max}\Big(\frac1n \bs X_{S}^\top\bs X_{S}\Big)$$
\end{lem}

It follows that a necessary assumption is that the missing active signals in any underfitted $M$ are not too "small". We next introduce a formal distinction between small and large signals in $S$. We also introduce a subset of signals that are not small or large that will be used in Section~\ref{ssec:nobetaminhyp}.  For a fixed set of penalties $\kappa_j$'s and every $j=1,\ldots,b$, $\gamma$ and $h_j$ defined as in A7, let: 
\begin{eqnarray}
    &S_j^{S}(\kappa)\;&:=\;\Bigg\{\beta_i^{\star} \in S_j \,\Big|\,\sqrt{n\bar{\lambda}}|\beta_{i}^{\star}|=o\big(\sqrt{\kappa_j}\big) \Bigg\}\nonumber\\
    &S_j^{L}(\kappa)\;&:=\;\Bigg\{\beta_i^{\star} \in S_j \,\Big|\,\sqrt{\frac{(1-\gamma) n\rho(\bs X)}{6}}|\beta_{i}^{\star}| \,- \,\sqrt{\kappa_j} \;=\; \sqrt{\ln(s_j)} + h_j\Bigg\}\label{eq:defslsi}\\
    &S_j^{I}(\kappa)\;&:=\;S_j \setminus \big(S_j^{L}(\kappa) \cup S_j^{S}(\kappa)\big)\nonumber 
\end{eqnarray} \\
The susbet $S_j^{S}(\kappa)$ gathers signals in $S_j$ that are small with respect to the block penalty $\kappa_j$. The subset $S_j^{L}(\kappa)$ gathers the signals in $S_j$ that large in the sense that they satisfy assumption~A7. The subset $S_j^{I}(\kappa)$ gathers the signals in $S_j$ that are neither large nor small. Take $M=S\setminus S^{S}(\kappa)$ where $S^{S}(\kappa)=\cup_{j=1}^b S_j^{S}(\kappa)$. If $S^{S}(\kappa)\neq \emptyset$, the missing active signals in $M$ and thus $\mu_{SM}$ are all small compared to the $\kappa_j$'s. Using the properties of the noncental $\chi^2$ distribution, we then establish the following convergence of the probabilities of $(NC(M)/NC(S)<1)$ and thus of $(\hat{S}^{b}=S)$.

\begin{prop}\label{prop:nectypeIIreg}
    If $S^{S}(\kappa)\neq\emptyset$ and for all $j=1,\ldots,b$, $\kappa_j\to \infty$, then, under assumptions A1,
    $$P\Big(\frac{NC(S\setminus S^{S}(\kappa))}{NC(S)}<1\Big) \to 0 \quad\text{and}\quad P(\hat{S}^b=S)\to 0$$
\end{prop}
Property~\ref{prop:nectypeIIreg} shows that for any increasing penalty, when $S$ contains signals that are too small with respect to that penalty, the smaller model $S\setminus S^{S}(\kappa)$ is favored at the expense of $S$ and recovery fails. It follows that a necessary assumption for variable selection consistency is:
\begin{equation}\label{cond:necIIallblockpen}
\lim_{n\to\infty}\frac{\sqrt{n\bar{\lambda}}\beta_{\min,j}^{\star}}{\sqrt{\kappa_j}}\;>\;0\qquad\mbox{for all }j=1,\ldots,b.
\end{equation}
When active variables are orthogonal and normalized, that is $\bs X_S^{\top}\bs X_S = nI$, $\bar{\lambda}=1$. Under orthogonal design, the assumption~\eqref{cond:necIIallblockpen} is then milder and less tight than the previously derived assumption~\eqref{cond:necIallblockthresh} for the sequence model. 

An immediate consequence of Proposition~\ref{prop:nectypeIreg} and Proposition~\ref{prop:nectypeIIreg} is the next result.
\begin{cor}\label{cor:necrecovreg}
    If for some $j\in\{1,\ldots,b\}$   $\lim_{n\to\infty}\frac{\sqrt{n\bar{\lambda}}\beta_{\min,j}^{\star}}{\underline{\lambda}_j\sqrt{\ln(p_j-s_j)}} = 0$, then, under assumption A1 and A2,  $P(\hat{S}^b=S)\not\to 1$
\end{cor}
Corollary~\ref{cor:necrecovreg} yields a necessary assumption for consistent variables selection under A1 and A2:
\begin{equation}\label{eq:betaminnecreg}
\lim_{n\to\infty}\frac{\sqrt{n\bar{\lambda}}\beta_{\min,j}^{\star}}{\sqrt{\underline{\lambda}_j\ln(p_j-s_j)}} > 0\qquad\mbox{for all }j=1,\ldots,b.
\end{equation}
\david{In the denominator above it should be $\underline{\lambda}_j \sqrt{p_j-s_j}$ no? That's what you have in Prop 4.1.}\paul{No, $\beta_{min}$ has to be compared to $\sqrt{\kappa_j}$.}
In an orthogonal and normalized design such that $\bs X^{\top} \bs X=nI$, for all $j=1,\ldots,b$, $\underline{\lambda}_j=1$ and $\bar{\lambda}=1$. In the particular case of orthogonality and diverging $s$, the previously derived necessary assumption \eqref{eq:betamin} is then stricter and tighter that assumption \eqref{eq:betaminnecreg}. Assumption \eqref{eq:betaminnecreg} is however a more general necessary betamin condition that applies in all correlated, orthogonal, low and high dimensional setting, for $s$ fixed or diverging. \paul{I am happy about this result because I believe there was no such tight necessary assumption for variable selection consistency in regression in the literature. The only ones I found were for LASSO (far less efficient/tight for selection in correlated settings) or compressed sensing}.


\subsection{Benefits of block penalties}
\label{ssec:regression_gains}

We discuss separately the benefits in conditions for variable selection consistency and those in the convergence rate. The main message
is that benefits discussed in the sequence model in Section~\ref{sec:properties_orth} extend to linear regression. 
\bigskip

\noindent Assumptions A6--A7 give ranges of penalties values that are sufficient (but not necessary) for asymptotic support recovery. For simplicity, we restrict our discussion to choices of penalty $\kappa_j$ such that, in A6, $f_j=\Omega(\ln(p_j-s_j))$ and $\gamma$ is bounded in $(0,1/2)$. For the standard threshold selector $\hat{S}$, the single penalty $\kappa$ is essentially required to satisfy, up to constant and bounded factors:
\begin{equation}\label{eq:rangepenex}
  \sqrt{\ln(p-s)}\;\;<\;\;\sqrt{\kappa}\;\;<\;\;\sqrt{n\rho(\bs X)}\beta^{\star}_{\min}+\sqrt{\ln(s)}.	
\end{equation}
\paul{This is not completely accurate. Is that fine as approximation? Alternatively we could define a $\gtrsim$ such that $g\gtrsim h$ means $g=h+$ any increasing function.}
For a block threshold selector $\hat{S}^{b}$, assuming without loss of generality that the global minimum $\beta_{\min}^*$ is in block $b$, the ranges for the $\kappa_j$'s essentially are, up to constant and bounded factors:
    \begin{eqnarray}\label{eq:rangeblockpenex}
  \sqrt{\ln(p_1-s_1)}\;\;&<\;\;\sqrt{\kappa_1}\;\;<&\;\;\sqrt{n\rho(\bs X)}\beta^{\star}_{\min,1}-\sqrt{\ln(s_1)}\nonumber\\
  &\;\;\vdots \\
  \sqrt{\ln(p_b-s_b)}\;\;&<\;\;\sqrt{\kappa_b}\;\;<&\;\;\sqrt{n\rho(\bs X)}\beta^{\star}_{\min}-\sqrt{\ln(s_b)} \nonumber.       
    \end{eqnarray}
As for the sequence model, by definition, the ranges for $\hat S^b$ are wider than that for $\hat S$. Sufficient conditions to have variable selection consistency are then milder with block penalties and results of Section~\ref{ssec:neccond_linreg} allow to precisely describe settings in which $\hat S^b$ is consistent and $\hat S$ is not.
\begin{cor}\label{cor:selconstimpossiblereg}
Assume A1, A2, A5, A6. If
\begin{eqnarray*}
&\lim_{n\to\infty}\frac{\sqrt{n\bar{\lambda}}\beta_{\min}^{\star}}{\underline{\lambda}\sqrt{\ln(p-s)}} = 0
\end{eqnarray*}
then 
$P(\hat{S} = S)\not \to 1$ and $P(\hat{S}^{b} = S)\to 1$.
\end{cor}

Observe that equations in \eqref{eq:rangepenex} and \eqref{eq:rangeblockpenex} and the same as for the sequence model, up to rescaling by $\sqrt{2/n}$ and factor $\rho(\bs X)$. Similarly, assuming $\beta_{min}$ in block $b$, the smallest signal recoverable by $$


Assume without loss of generality that $\beta_{\min}^{\star}$ is located in block $b$.  By \eqref{betamincond:minimalreg}, the minimal betamin conditions for the block informed $\ell_0$ penalty $\hat{S}^b$ and standard $\ell_0$ penalty $\hat{S}$ are respectively: 
$$\lim_{n\to \infty}\frac{\sqrt{(1-\gamma) n\rho(\bs X)}\,{\beta_{\min}^{\star}}}{\sqrt{2\ln(p_b-s_b)} + \sqrt{2\ln(s_b)}} >1 \qquad \mbox{and} \qquad \lim_{n\to \infty}\frac{\sqrt{(1-\gamma) n\rho(\bs X)}\,{\beta_{\min}^{\star}}}{\sqrt{2\ln(p-s)} + \sqrt{2\ln(s)}} >1$$
The requirements on $\beta_{\min,j}^\star$ and $\beta_{\min}^\star$ again relate to
%The relevant quantities to compare the minimal requirement on the growth rate of betamin are thus again 
$\sqrt{\ln(p_b-s_b)} + \sqrt{\ln(s_b)}$ and $\sqrt{\ln(p-s)} + \sqrt{\ln(s)}$, hence the discussion in Section \ref{sec:prosconsblock} on the benefits in betamin conditions extend directly to linear regression.
In particular, there are linear regression settings where variable selection consistency is not possible with $\hat{S}$, but is with $\hat{S}^b$. For example, in a simplest case where $\bs X^\top \bs X=n I_p$, consistency results for the sequence model apply here
and Corollary \ref{cor:selconstimpossible} outlines such instances.

Let ${OR}^b_{reg}$ be the oracle convergence rate for $\hat{S}^{b}$ in Theorem~\ref{prop:fullrecovreg}, and ${OR}_{reg}$ that for $\hat{S}^{1}$ . Then
\begin{equation*}
    \frac{{OR}^b_{reg}}{{OR}_{reg}}=(2^{2b-1}-b)\frac{r_b}{r_1}\sum_{j=1}^be^{-\frac{\delta}{2}\left[\frac{(1-\gamma) n\rho(\bs X)}{32}\left({\beta^\star_{\min,j}}^2-{\beta_{\min}^\star}^2\right)+\max\{\ln(p-s),\ln(s)\}-\max\{\ln(p_j-s_j),\ln(s_j)\}\right]}
\end{equation*}
which is essentially the same as in the sequence model
in \eqref{eq:ratioCR}, up to certain
factors. 
Specifically, $r_b/r_1$ and $\delta$ are both close to 1, whereas $(1-\gamma) \rho(\bs X)/32$ slows the convergence rate gains
but does not alter the essence of the ratio
as long as $\rho(\bs X)$ remains bounded above 0. The factor $2^{2b-1}-b$ highlights a potential limitation: guarantees of gains with block $\ell_0$ penalties deteriorate when one considers a large number of blocks $b$.
We remark that such deterioration may be a consequence of our proof strategy, rather than an inherent limitation of block penalties. Studying cases with large $b$, including $b \to \infty$, is left as future research.

\subsection{Convergence with no betamin condition}\label{ssec:nobetaminhyp}

We now derive a convergence result for pseudo-posterior probabilities $NC(M)$, $M\in\mathcal{M}$ under no assumption on the minimal signal strength. The result generalizes Theorem~\ref{theo:suffcondlinearmodel} and is key to the proofs of Section~\ref{sec:informed_l0}.

In \eqref{eq:defslsi}, we defined, for a choice of penalties $\kappa_j$, the set of large and intermediate signals in $S$. Consider the set of models that contain $S^L(\kappa)=\cup_{j=1}^b S^{L}_j(\kappa)$ and some or none of the signals in $S^{I}(\kappa)=\cup_{j=1}^b S^{I}(\kappa)_j$.
\begin{equation}
    \T \;:=\;\Big\{M \in \M \,\big|\, M\,=\, S^L(\kappa)\,\cup\,R,\; R\in \mathcal{P}\big(S^{I}(\kappa)\big)  \Big\}
\end{equation}
where for any set $A$, $ \mathcal{P}(A)$ denotes the power set of $A$.
Following a strategy similar to that of the proof of Theorem~\ref{theo:suffcondlinearmodel}, we demonstrate, assuming only sufficiently large penalties (A6), the vanishing of the sum of pseudo-probabilities of all models but those in $\T$. The main difference with respect to the proof of Theorem~\ref{theo:suffcondlinearmodel} is that, for any $M\not\in \T$, we bound $\E(NC(M))$ by a function of $C(M)-C(T)$ for a well chosen model $T$ in $\T$. 
\begin{thm}\label{thm:convtoT}
    Assume A1, A2, A6 and $|S^{S}_j(\kappa)|=O(p_j-s_j)$ for every $j=1,\ldots,b$:
    \begin{longlist}
        \item $\sum_{M \in \M\setminus \T} \E\left(NC(M)\right)\to 0$,
        \item if $|S^{I}(\kappa)|=O(1)$, then $P(\hat{S}^b \in \T)\to 1$.
    \end{longlist}
\end{thm}
Theorem~\ref{thm:convtoT} shows that, asymptotically, with only assuming sufficiently large penalties, all inactive and small active signals are discarded. On the contrary, all large signals are retained. Intermediate signals in $S^{I}(\kappa)$ may or may not be retained. Note that assumption $|S^{I}(\kappa)|=O(1)$ in $(ii)$ was made for simplicity and $|S^{I}(\kappa)|$ can be allowed to increase moderately. 

\section{Data analysis with block $\ell_0$ penalization}\label{sec:informed_l0}

\david{I made a number of edits to the section, to clarify the connection between block $\ell_0$ and the Bayesian paradigm, particularly for non-Bayesian readers. Piotr in particular please chime in, if things don't make sense.}

\subsection{A Bayesian-inspired estimator of sparsity}

In Sections~\ref{sec:properties_orth} and~\ref{sec:properties_legression}, we derived oracle properties for block penalties where one sets the penalties to oracle values.
We now
propose two data analysis methods anchored in a Bayesian perspective
that do not require oracle values, and adapt to the unknown sparsity in the data-generating truth.
A Bayesian approach provides a natural estimator of the number of truly active variables in each block and thus, by using empirical Bayes, a straightforward way to adapt penalties to sparsity in each block.
\\
\\
Consider a Bayesian framework for \eqref{eq:linearmodel} where
the model $M$ is a $p$-dimensional vector of variable inclusion indicators $m_i=I(\beta_i \neq 0)$. Consider a joint prior on parameters and models 
\begin{equation}\label{eq:bayesianframework}
    p(\boldsymbol{\beta},M \mid \boldsymbol{\theta})=p( \boldsymbol{\beta} \mid  M) p(M \mid \boldsymbol{\theta})
\end{equation}
where $p(M \mid \boldsymbol{\theta})$ is the prior probability of model $M$, which depends on hyperparameter $\bs \theta$,
and $p( \boldsymbol{\beta} \mid  M)$ is a prior on regression coefficients given the model. 

Posterior model probabilities are $p(M \mid \bs y, \bs \theta) \propto p(\bs y \mid M) p(M \mid \bs \theta)$, where $p(\bs y \mid M)$ is the so-called marginal likelihood of model $M$. The BIC approximation \citep{BIC} gives
\begin{eqnarray*}
 p( \bs y  \mid M)\approx p(\bs y \mid \tilde{\bs \beta}^{(M)})n^{-|M|/2}.
\end{eqnarray*}
It follows that the posterior model probabilities satisfy
\begin{eqnarray*}
\ln p(M \mid\bs y) \approx \ln p(\bs y \mid \tilde{\bs \beta}^{(M)}) - \frac{1}{2}\ln(n)|M| + \ln p(M\mid\bs \theta) + \text{constant}
\end{eqnarray*}
and any block or standard $\ell_0$ penalty selector discussed in earlier sections can be written as $\hat{M}= \arg\max_M p(M \mid \bs y)$. 
In particular, suppose that variable inclusions are independent a priori, and that inclusion probabilities are constant within each block. That is,
\begin{eqnarray} \label{eq:block_modelprior}
p( M \mid \boldsymbol{\theta})=
        \prod_{i=1}^p \operatorname{Bern}\left(m_i ; \theta_i\right) I(M \in \M) \mbox{, where } \forall \, i\in B_j, \: \theta_i=\theta^{(j)}
\end{eqnarray}
then $\ln p(M \mid \bs \theta)$ defines the additive block penalties studied in Section \ref{sec:properties_legression} with:
\begin{equation}\label{eq:blockpenpriorinclu}
    \kappa_j=\tfrac{1}{2}\ln(n)+\ln\big({\theta^{(j)}}^{-1}-1\big).
\end{equation}
This connection motivates an estimator of the number of truly active variables in any block $B_j$, for any block $\ell_0$ procedure:
\begin{equation}\label{eq:shat}
    \widehat{s}_j \,:=\, \sum_{i\in B_j}\,\,\sum_{M\in \M | i \in M}NC(M)
\end{equation}
The normalized score $NC(M)$ can be seen as the posterior probability for model $M$ in Bayesian framework \eqref{eq:bayesianframework}. Replacing normalized scores  by posterior model probabilities in \eqref{eq:shat} gives:
\begin{equation*}
    \sum_{i\in B_j}\,\,\sum_{M\in \M | i \in M}P(M\mid \bs y) = \sum_{i\in B_j}P(\beta_i \neq 0 \mid \bs y) = E\big(s_j \mid \bs y\big)
\end{equation*}
From a Bayesian perspective, $\widehat{s}_j$ is then the sum of posterior inclusion probabilities and the posterior mean of $s_j$. The estimator of the proportion of truly active variables $\widehat{s}_j/p_j$ can also be formally justified as an empirical Bayes procedure. The empirical Bayes estimator of $\bs \theta$ (cf \eqref{eq:block_modelprior}) maximizes the marginal likelihood of $\bs y$ given $\bs \theta$ and is defined as:
\begin{eqnarray}
\hat{\bs \theta}^{EB}\,:=\, \arg\max_{\bs \theta} p(\bs y \mid \bs \theta)=
\arg\max_{\bs \theta} \int p(\bs y \mid \bs \beta, M) d P(\bs \beta, M \mid \bs \theta).
\nonumber
\end{eqnarray}
 In the Bayesian framework \eqref{eq:block_modelprior}, $\widehat{s}_j/p_j$ corresponds to the empirical Bayes estimate of $\bs \theta$.
\begin{lem}\label{lem:empiricalbayesinterpret} For $p(M \mid \bs \theta)$ in \eqref{eq:block_modelprior}, the empirical Bayes estimator satisfies $$\hat{\theta}^{EB}_i = \frac{1}{p_j}\sum_{i\in B_j}P(\beta_i \neq 0 \mid \bs y)$$
\end{lem}
Estimator $\widehat{s}_j/p_j$ is well founded from a Bayesian perspective and also has attractive frequentist properties under mild assumptions. In Section~\ref{sec:properties_legression} we defined $S^{S}_j(\kappa)$ to be the set of small signals in block $j$ and  $S^{L}_j(\kappa)$ the set of larger signals (\textit{cf} \eqref{eq:defslsi}). In Theorem~\ref{thm:convtoT} we also derived a convergence result on pseudo posterior probabilities that yields the following asymptotic bounds on the expectation of $\widehat{s}_j/p_j$.
\begin{prop}
    \label{prop:shatconsistence}
Assume A1, A2, A6 and $|S^{S}_j(\kappa)|=O(p_j-s_j)$ for every $j=1,\ldots,b$. We have:$$ \,\, \frac{|S^{L}_j(\kappa)|}{p_j}\leq \lim_{n\to\infty}\E\bigg(\frac{\widehat{s}_j}{p_j}\bigg)\leq \frac{s_j-|S^{S}_j(\kappa)|}{p_j}\qquad\text{for all }j=1,\ldots,b$$
\end{prop}
An immediate consequence is that, if A7 also holds, then $\widehat{s}_j/p_j \stackrel{L_1}{\longrightarrow} s_j/p_j$
because $S_j=S_j^{L}(\kappa)$ and $S_j^{S}(\kappa)=\emptyset$ for all $j$. The proposition also shows that, when A7 is not met, $\widehat{s}_j/p_j$ is asymptotically downward biased by at least $|S^{S}_j(\kappa)|/p_j$ but larger than the proportion of signals satisfying A7. 

\subsection{Data-based block selection method}

We propose a selection method that builds on the connection in \eqref{eq:blockpenpriorinclu} and relies on $\widehat{s}_j$ to adapt the $\ell_0$ penalty in each block. It has two variants. The first one follows the empirical Bayes approach of Lemma~\ref{lem:empiricalbayesinterpret} to set block penalties. The second variant is motivated by assumption A6 on the sufficient penalty strength for variable selection consistency. The BIC approximation in \eqref{eq:blockpenpriorinclu} provides a natural choice for $f_j$ in A6, $f_j=\frac{1}{2}\ln(n)$ for all $j$. The choice $\theta^{(j)}=(p-s+1)^{-1}$ then yields block penalties that immediately satisfy A6. The second variant approximates this choice.\\

\textbf{Algorithm}
\begin{enumerate}
    \item Set
    $\kappa_j=\kappa^{\circ}=\ln(p)+\frac{1}{2}\ln(n)$ for $j=1,\ldots,b$. Compute $\widehat{s}_j/p_j$ in \eqref{eq:shat} for $j=1,\ldots,b$.

    \item Obtain $\hat{S}^{EB,b}$ solving \eqref{eq:Mhat} with $\kappa_j^{EB}= \ln(p_j/\widehat{s}_j -1)+ \tfrac{1}{2}\ln(n)$.
    Alternatively, obtain $\hat{S}^{IA,b}$ solving \eqref{eq:Mhat} with
    $\kappa_j^{IB} = \ln(p_j-\widehat{s}_j)+ \tfrac{1}{2}\ln(n)$.
\end{enumerate}

We refer to $\hat{S}^{EB,b}$ as the block empirical Bayes selector, and to $\hat{S}^{IA,b}$ as the block adaptive selector.
Step 1 approximates posterior model probabilities under prior inclusion probabilities $\theta^{(j)}=(p+1)^{-1}$ in \eqref{eq:block_modelprior} that are equal across blocks, and estimates the proportion of truly active coefficients with $\hat{s}_j/p_j$. Step 2 selects a model using the block penalties $\kappa_j^{EB}$ induced by $\theta^{(j)}=\hat{s}_j/p_j$,
or alternatively setting $\kappa_j^{IB}= \kappa_j^{EB} + \ln(\hat{s}_j)$.
Under mild assumptions both $\kappa_j^{EB}$ and $\kappa_j^{IB}$ lead to variable selection consistency.
As shown below, the main difference is that in non-sparse settings where $s_j$ is fairly large, the penalty $\kappa_j^{EB}$ can be insufficient and there $\kappa_j^{IB}$ might be preferred.

To show the consistency of $\hat{S}^{EB,b}$ and $\hat{S}^{IA,b}$ we rely on Theorem~\ref{theo:suffcondlinearmodel}. 
Consider assumptions:\\[5pt]
(A8)\quad $s=o(\sqrt{n})$\\[5pt]
(A9)\quad For each block $j$, there exists $h'_j\to +\infty$such that for $n$ sufficiently large, 
\begin{equation*}\label{cond:suffIIreg2}
        \sqrt{\frac{ (1-\psi)n\rho(\bs X)}{6}}{{\beta_{\min}^{\star}}_j} \,- \,\sqrt{\ln\bigg(\frac{p_j}{|S^{L}_j(\kappa^{\circ})|}-1\bigg)+\frac12\ln(n)} \;=\; \sqrt{\ln(s_j)} + h'_j .
    \end{equation*}
where $\gamma'=\frac12\big(1+\max_j \frac{\ln(p_j-s_j)}{\ln(p_j/s_j-1)+0.5\ln(n)}\big)$.
Using Theorem~\ref{thm:convtoT} and Proposition~\ref{prop:shatconsistence}, we show that under A8 and A9, the $\kappa^{EB}_j$'s satisfy A6 and A7 with probability going to 1 and conclude on the selection consistency  of $\hat{S}^{EB,b}$.

\begin{thm}\label{theo:empbayesselconsist}
    Assume A1, A2, A8, A9, and $|S^{S}_j(\kappa)|=O(p_j-s_j)$ for every $j=1,\ldots,b$. Then
$P(\hat{S}^{EB,b} = S) \to 1$.
\end{thm}
For the consistency of $\hat{S}^{IA,b}$, we no longer require A8, and A9 is replaced by the slightly more stringent A10 below. \\[5pt]
(A10)\quad For each block $j$, there exists $h''_j\to +\infty$ such that for $n$ sufficiently large,
\begin{equation*}
        \sqrt{\frac{ (1-\psi)n\rho(\bs X)}{2}}{\beta_{\min}^{\star}}_j \,- \,\sqrt{\ln\big(p_j-|S^{L}_j(\kappa^{\circ})|\big)+\frac12\ln(n)} \;=\; \sqrt{\ln(s_j)} + h''_j .
    \end{equation*}
where  $\psi'=\frac12(1+\max_j \frac{\ln(p_j-s_j)}{\ln(p_j-s_j)+0.5\ln(n)})$.

\begin{thm}\label{theo:informedl0consist}
    Assume A1, A2, A10, and $|S^{S}_j(\kappa)|=O(p_j-s_j)$ for every $j=1,\ldots,b$. Then 
$P(\hat{S}^{IA,b} = S) \to 1$. \david{as $n$ and $p$ grow? Or as $n$ grows only? Same clarification for Theorem 5.1.} \paul{A1 and A2 are $n\to+\infty$ and $p-s\to\infty$ so it's assuming $n$ and $p$ grow.}
\end{thm}

\subsection{Benefits of data-based block selection method}

Theorem~\ref{theo:empbayesselconsist} also yields selection consistency of standard empirical Bayes selector that sets a single step 2 penalty $\kappa^{EB}=\ln(p/\widehat{s} -1)+ \tfrac{1}{2}\ln(n)$. For such selector the sufficient betamin assumption is:
\begin{equation}\label{eq:betaminIEB}
        \sqrt{\frac{ (1-\gamma'')n\rho(\bs X)}{6}}{\beta_{\min}^{\star}} \,- \,\sqrt{\ln\bigg(\frac{p}{|S^{L}(\kappa^{\circ})|}-1\bigg)+\frac12\ln(n)} \;=\; \sqrt{\ln(s)} + h''_j .
    \end{equation}
where $\gamma''=\frac12\big(1+ \frac{\ln(p-s)}{\ln(p/s-1)+0.5\ln(n)}\big)$ and $h''_j \to \infty$. Observe that assumption A9 is milder than \eqref{eq:betaminIEB}. The smaller $\frac{p_j}{|S^{L}_j(\kappa^{\circ})|}$ is compared to $\frac{p}{|S^{L}(\kappa^{\circ})|}$, the milder is A9 in block $B_j$. 

\paul{I believe we cannot comment on differences in rates of convergence between block and standard EB. Theorem on rates in Section~\ref{sec:properties_legression} hold when A6 and A7 hold. We can only prove A6 and A7 hold with probability going to 1. To have a rate for the EB we would need to have a rate at which A6 and A7 hold.}

\section{Empirical results}\label{sec:simulation}

\section{Discussion}\label{sec:discussion}

We provided an understanding of how external information can ease variable selection in the Gaussian Sequence model and high-dimensional linear regression. We studied the case where external information helps partition variables into blocks with potentially distinct characteristics and introduced externally-informed $\ell_0$ selectors that let the penalty depend on those blocks. We showed that an oracle externally-informed selector converges faster and has less strict betamin condition than the standard $\ell_0$ oracle. We also provided concrete data analysis methods that incorporate external information to improve variable selection properties without requiring oracle knowledge.

A question that immediately arise for future work is how much the assumption of fixed number of blocks $b$ can be relaxed. Our consistency results are robust to moderate increases in the number of blocks but cannot hold when $b=p$ for example. Another interesting direction of research is understanding the benefit of external information for parameter estimation and estimation error. \paul{Adapting a proof in Rigollet's high-dim statistics notes I could easily give here a bound on estimation error in the Gaussian sequence model. I don't how tight it is though. For linear regresison, there is proof for BIC in Rigollet's notes and I need to check more closely if it can be adapted. Do you know David a proof for the estimation error with $\ell_0$ selection in linear regression? Would it be interesting/relevant to have those bounds here?} \paul{We could also mention externally informed $\ell_1$ in correlated deigns as future direction.}
\david{Castillo, Schmidt-Hieber and van der Vaart have work on estimation error for the sequence model and linear regression, using posterior model probabilities that are analogs to our pseudo-posterior probabilities. I suppose that one could adapt their proof strategy. Miquel Torrens also studied MSE's in the sequence model, his proof strategy was different and more along the lines of what you did here. This said, I'd probably leave this for future work. It might be an easy-ish paper, perhaps.}

We also observe that our results bear similarities with conclusions of works on related problems. In Section \ref{sec:examples}, we discussed how the benefits of externally-informed selectors depend on how discriminative is the external information. \cite{scarlett2012compressed} similarly showed that support recovery in compressed sensing becomes easier as the entropy of the distribution of inclusion probabilities across blocks shrinks. In hypothesis testing, \cite{genovese2006} demonstrated that the power of their informed weighted testing procedure increases when assignment of weights is positively associated with the null hypotheses being false. We expect analogous results to be achievable in other related problems like structural learning in graphical models or variable selection over partially ordered sets as introduced in \cite{taeb}. 


%%%%%%%%%%%%%%%%%%%%%%%%%%%%%%%%%%%%%%%%%%%%%%
%% Single Appendix:                         %%
%%%%%%%%%%%%%%%%%%%%%%%%%%%%%%%%%%%%%%%%%%%%%%
%\begin{appendix}
%\section*{???}%% if no title is needed, leave empty \section*{}.
%\end{appendix}
%%%%%%%%%%%%%%%%%%%%%%%%%%%%%%%%%%%%%%%%%%%%%%
%% Multiple Appendixes:                     %%
%%%%%%%%%%%%%%%%%%%%%%%%%%%%%%%%%%%%%%%%%%%%%%
%\begin{appendix}
%\section{???}
%
%\section{???}
%
%\end{appendix}

%%%%%%%%%%%%%%%%%%%%%%%%%%%%%%%%%%%%%%%%%%%%%%
%% Support information, if any,             %%
%% should be provided in the                %%
%% Acknowledgements section.                %%
%%%%%%%%%%%%%%%%%%%%%%%%%%%%%%%%%%%%%%%%%%%%%%
%\begin{acks}[Acknowledgments]
% The authors would like to thank ...
%\end{acks}
%%%%%%%%%%%%%%%%%%%%%%%%%%%%%%%%%%%%%%%%%%%%%%
%% Funding information, if any,             %%
%% should be provided in the                %%
%% funding section.                         %%
%%%%%%%%%%%%%%%%%%%%%%%%%%%%%%%%%%%%%%%%%%%%%%
%\begin{funding}
% The first author was supported by ...
%
% The second author was supported in part by ...
%\end{funding}

%%%%%%%%%%%%%%%%%%%%%%%%%%%%%%%%%%%%%%%%%%%%%%
%% Supplementary Material, including data   %%
%% sets and code, should be provided in     %%
%% {supplement} environment with title      %%
%% and short description. It cannot be      %%
%% available exclusively as external link.  %%
%% All Supplementary Material must be       %%
%% available to the reader on Project       %%
%% Euclid with the published article.       %%
%%%%%%%%%%%%%%%%%%%%%%%%%%%%%%%%%%%%%%%%%%%%%%
%\begin{supplement}
%\stitle{???}
%\sdescription{???.}
%\end{supplement}

%%%%%%%%%%%%%%%%%%%%%%%%%%%%%%%%%%%%%%%%%%%%%%%%%%%%%%%%%%%%%
%%                  The Bibliography                       %%
%%                                                         %%
%%  imsart-???.bst  will be used to                        %%
%%  create a .BBL file for submission.                     %%
%%                                                         %%
%%  Note that the displayed Bibliography will not          %%
%%  necessarily be rendered by Latex exactly as specified  %%
%%  in the online Instructions for Authors.                %%
%%                                                         %%
%%  MR numbers will be added by VTeX.                      %%
%%                                                         %%
%%  Use \cite{...} to cite references in text.             %%
%%                                                         %%
%%%%%%%%%%%%%%%%%%%%%%%%%%%%%%%%%%%%%%%%%%%%%%%%%%%%%%%%%%%%%

%% if your bibliography is in bibtex format, uncomment commands:
\bibliographystyle{imsart-number} % Style BST file (imsart-number.bst or imsart-nameyear.bst)
\bibliography{bibliography.bib}       % Bibliography file (usually '*.bib')

%% or include bibliography directly:
% \begin{thebibliography}{}
% \bibitem{b1}
% \end{thebibliography}
XXXXXXXXXXXXXXXXXXXXXXXXXXXXXXXXXXXXXXXXXXXXXXXXXXXXXXXXXXXXXXXXX\\
EVERYTHING AFTER WILL GO TO SUPPLEMENT\\
XXXXXXXXXXXXXXXXXXXXXXXXXXXXXXXXXXXXXXXXXXXXXXXXXXXXXXXXXXXXXXXXX\\
\section{Auxiliary results}\label{auxiliary}

In this section we collect some technical results. Their proofs are gathered in section~\ref{sec:proofauxi}. 

\begin{lem}\label{lem:l1isthreshold}
	In the sequence model \eqref{eq:GSM}, selecting the non-zero entries  of the solution to 
\begin{equation}\label{eq:genl1}
	\underset{\bs \beta \in \R^p}{\rm minimize} \;\;\tfrac12\|\bs y-\sqrt{n}\bs\beta\|^2+\sum_{j=1}^b \lambda_j \sum_{i \in B_j}|\beta_{i}|
\end{equation}
for non-negative $\lambda_1,\ldots,\lambda_b$  is equivalent to taking the selector $\hat S^{b}$ with $\bs \tau=(\lambda_1/n,\ldots,\lambda_b/n)$. Let $\bs\beta^\circ$ be either the MLE or the LASSO estimator with penalty $\lambda^\circ$. Selecting the non-zero entries  of the solution to the problem
\begin{equation}\label{eq:genl1ada}
	\underset{\bs \beta \in \R^p}{\rm minimize} \;\;\tfrac12\|\bs y-\sqrt{n}\bs\beta\|^2+\sum_{j=1}^b \lambda_j \sum_{i \in B_j}\tfrac{|\beta_{i}|}{|\beta_i^\circ|}.
\end{equation}
is equivalent to taking the selector $\hat S^{b}$ with $\bs \tau=(\sqrt{\lambda_1/n},\ldots,\sqrt{\lambda_b/n})$ if $\bs\beta^\circ$ is the MLE and with $\bs \tau=(\lambda^\circ/2n+\sqrt{(\lambda^\circ/2n)^2+\lambda_1/n},\ldots,\lambda^\circ/2n+\sqrt{(\lambda^\circ/2n)^2+\lambda_b/n})$ if $\bs\beta^\circ$ is a LASSO estimate.
\end{lem}


\begin{lem}\label{lemma:3tailbounds}\textbf{Tail bounds on the maximum and minimum of folded Gaussians} 
\begin{longlist}
        \item If $\bs y\sim N_p\left(0,\frac{1}{n}I_n \right)$, $p>1$, and $a \geq \sqrt{\frac{2\ln(p)}{n}}$, $$P\Big(\underset{i\in\{1,\ldots,p\}}{\max}|y_i|> a \Big) \leq \frac{e^{-\frac{n}{2}\left(a^2-\frac{2\ln (p)}{n}\right)}}{\sqrt{\pi\ln (p)}}.$$
    \item If $\bs y\sim N_s\left(\bs\beta,\sigma^2\right)$, $s>1$, and $a \leq \underset{i\in\{1,\ldots,s\}}{\min}|\beta_i|$,
        $$P\Big(\underset{i\in\{1,\ldots,s\}}{\min}|y_i| > a\Big) \geq P\Big(\underset{i\in\{1,\ldots,s\}}{\min}|\beta_i|-\underset{i\in\{1,\ldots,s\}}{\max}|y_i-\beta_i| > a \Big).$$
    \item If  $\forall i\in\{1,\ldots,s\}$ $y_i\sim N\left(\mu,\frac{1}{n}\right)$, and $a<|\mu|$
        $$P\Big(\underset{i\in\{1,\ldots,s\}}{\min}|y_i| > a \Big) \leq \exp\left\{-\frac{s}{2}\frac{e^{-2n(|\mu|-a)^2/\pi}-e^{-n(a+|\mu|)^2/2}}{\left(1-e^{-2n(|\mu|-a)^2/\pi}\right)^{\frac{1}{2}}+\left(1-e^{-n(a+|\mu|)^2/2}\right)^{\frac{1}{2}}}\right\}.$$
\end{longlist}
\end{lem}

\begin{lem}\label{lem:L0toL1convergencegen} For $\hat S^b$ as in \eqref{eq:Mhat} and any $k$ models $M_1,\ldots,M_k$
$$P(\hat{S}^b\not\in \{M_1,\ldots,M_k\})\;\leq\; (k+1) \sum_{M \in \M\setminus \{M_1,\ldots,M_k\}}\E\left( NC(M)\right).$$
\end{lem}

\begin{lem}\label{lem:boundrhogen}
 For any $T\subseteq S$ and $M\in\M$ such that $T\not\subseteq M$, let $Q_T=M \cup T$. The non-centrality parameter defined in \eqref{eq:muQM} satisfies: 
\begin{equation}
    \mu_{Q_TM}\; \geq \;n\,\rho(\bs X)\,\sum_{j=1}^b|T_j\setminus M_j|\,\min_{i\in T_j\setminus M_j}{\beta_{i}^{\star}}^2.
\end{equation}
\end{lem}

\begin{lem}\label{lemma:s1s3}
Let $W\sim\chi_{\nu}^2(\mu)$ with $\mu \geq 0$, then for any $w>\mu+\nu$
$$P(W > w) \leq e^{-\big(\frac{w+\mu}{2}  - \sqrt{2w( 2 \mu+ \nu)-2 \mu \nu - \nu^2 } \big)}. $$
Moreover, if $w$ is increasing such that $\nu=o(w)$ and $\mu=o(w)$, for any $\phi\in(0,1)$ and $n$ large enough
$$P(W > w)  \leq  e^{-\phi\frac{w}{2}}.$$
\end{lem}

\begin{lem}\label{lemma:s2} Let $W \sim \chi_\nu^2(\mu)$ with $\mu>0$. For any $w<\mu$, 
$$
P(W<w) \leq \frac{e^{-\frac{1}{2}(\sqrt{\mu}-\sqrt{w})^2}}{(\mu / w)^{\nu / 4}}.$$
\end{lem}

\begin{lem}\label{lemma:s18} 
    Let $W\sim \chi_\nu^2(\mu)$ with $\mu \geq 0$. Assume $g$ positive and increasing such that $\nu=o(\ln(g))$, $\mu=o(\ln(g))$. Let $\bar{u},\underline{u}$ in $(0,1)$ such that $1>\bar{u}>\underline{u}\geq\left(1+g^{\phi}e^{-(\nu+\mu)/2}\right)^{-1}$ where $\phi\in(0,1)$, then for $n$ large enough, we have
$$
\int_{\underline{u}}^{\bar{u}} P\left(W>2\ln \left(\frac{g}{1 / u-1}\right)\right) d u \leq  \frac{1}{g^{\phi}}\Big(\bar{u}-\underline{u}+\ln\Big(\frac{\bar{u}}{\underline{u}}\Big)\Big).
$$

\end{lem}

\begin{lem}\label{lemma:s20}
 Let $W\sim \chi_\nu^2(\mu)$ with $\mu \geq 0$. Assume $g$ positive and increasing such that $\nu=o(\ln(g))$, $\mu=o(\ln(g))$, then for any $\alpha\in(0,1)$ and $n$ large enough, we have
$$
\int_0^1 P\left(W>2 \ln \left(\frac{g}{1 / u-1}\right)\right) d u\;\; \leq\;\;  g^{-\alpha}$$
\end{lem}

\begin{lem}\label{lem:boundexpnc} For any $T\subseteq S$ and $M\in \M\setminus\{T\}$, denote $Q_T= M\cup T$ and $A_T:=\gamma\Delta_{MT}+\tfrac{1-\gamma}{6}\mu_{Q_TM}$ (cf \eqref{eq:DeltaM} and \eqref{eq:muQM}). Suppose that, for some $\gamma \in (1/2,1]$, $A_T>0$, $|M\setminus T|=o(A_T)$, and $\mu_{Q_T T}=o(A_T)$. For any $\psi\in(0,1)$ and $n$ large enough:
$$\E\left(NC(M)\right)  \;\;\leq\;\; e^{-\psi A_T}.$$
\end{lem}

\begin{lem}\label{lem:upperboundnoncentralgen}
    For any $T\subseteq S$ and $M\in \M$, let $Q_T=M\cup T$, $$\mu_{Q_TT} \leq n\,\bar{\lambda}\,\sum_{j=1}^b|(S_j\cap M_j )\setminus T_j|\,\max_{i\in (S_j\cap M_j )\setminus T_j}{\beta_{i}^{\star}}^2 \quad \text{where}\quad \bar{\lambda}\,:=\,\lambda_{\max}\Big(\frac1n \bs X_{S}^\top\bs X_{S}\Big)$$
    and
    $$\mu_{Q_TT} \geq n\,\lambda_{\min}\Big(\frac1n \bs X_{S}^\top\bs X_{S}\Big)\,\sum_{j=1}^b|(S_j\cap M_j )\setminus T_j|\,\min_{i\in (S_j\cap M_j )\setminus T_j}{\beta_{i}^{\star}}^2$$
\end{lem}

\section{Proofs of Section~\ref{sec:properties_orth}}
 
 

\subsection{Proof of Proposition~\ref{prop:isthresh}}
First note that by discarding constant term 
 $$\arg\max_{M\in \M} \left\{\max_{\beta\in \mathcal{L}_M}\ell(\bs y; \bs \beta)-\sum_{j=1}^b\kappa_j |M_j| \right\}=\arg\min_{M\in \M} \left\{\min_{\beta\in \mathcal{L}_M}\bigg\{ \tfrac{1}{2}\|\bs y-\sqrt{n}\bs\beta\|^2\right\}+\sum_{j=1}^b\kappa_j |M_j| \bigg\}$$

 We also have $\min_{\beta\in \mathcal{L}_M} \tfrac{1}{2}\|\bs y-\sqrt{n}\bs\beta\|^2=\tfrac{1}{2}\|\bs y-\sqrt{n}\tilde{\bs \beta}_M\|^2=\tfrac{1}{2}\|\bs y\|^2-\tfrac{1}{2}\|\sqrt{n}\tilde{\bs \beta}_M\|^2$. The maximization in \eqref{eq:Mhat} can be then replaced with:
     \begin{equation}\label{eq:MhatGSM}
  \arg\max_{M\in \M} \bigg\{\tfrac{1}{2}\| \sqrt{n}\tilde{\bs \beta}_M\|^2-\sum_{j=1}^b\kappa_j |M_j|\bigg\}, \qquad M=M_1\cup \ldots\cup M_b.	
\end{equation}
Under the assumptions of \eqref{eq:GSM}, we can write 
	$$
	\tfrac{1}{2}\|\sqrt{n}\tilde{\bs \beta}_M\|^2-\sum_{j=1}^b\kappa_j |M_j|\;=\;\sum_{j=1}^b\sum_{i\in M_j}(\tfrac{n}{2}{\tilde{ \beta}_i}^2-\kappa_j).
	$$
	Then \eqref{eq:MhatGSM} can be maximized with respect to each $M_j$ by including $i\in \hat S_j$ whenever $n\tilde\beta_i^2\geq 2\kappa_j$.

\subsection{Proof of Proposition~\ref{prop:suffnecIblockthresh}} By the union bound and by Lemma~\ref{lemma:3tailbounds} (i),
\begin{equation}\label{eq:type1convrate}
P\left(\hat{S}^{b} \not\subseteq S\right)\;\leq\; \sum_{j=1}^b P\Big(\max_{i\in B_j\setminus S_j}|y_i/\sqrt{n}|>\tau_{j}\Big)\;\leq\;\sum_{j=1}^b\frac{e^{-\tfrac{n}{2}\left(\tau_j^2-\tfrac{2\ln(p_j-s_j)}{n}\right)}}{\sqrt{\pi\ln(p_j-s_j)}} 	.
 \end{equation}
 By A4, the nominator on the right is bounded above by $1$ for all sufficiently large $n$. It follows that 
	$$P\left(\hat{S}^{b} \not\subseteq S\right) \;\leq\; \sum_{j=1}^b (\pi\ln (p_j-s_j))^{-1/2} \to 0\qquad \mbox{as } n\to \infty.$$
For the second statement, by independence, we have 
$$P(\hat{S}^{b} \subseteq S)\;=\;\prod_{j=1}^b P(\max_{i \in B_j\setminus S_j}|y_i/\sqrt{n}| \leq \tau_j).$$ 
Consider $j$ such that $\lim_{n\to \infty} \frac{\sqrt{n}\tau_j}{\sqrt{2\ln(p_j-s_j)}} \;<\;1$. In particular, there exists $c<1$ such that, for sufficiently large $n$, we have that $\tau_j\leq c\sqrt{2\ln(p_j-s_j)/n}$. Then, for any such $n$,
$$
P\Big(\max_{i \in B_j\setminus S_j}|y_i/\sqrt{n}| \leq \tau_j\Big)\;\leq\; P\Big(\max_{i \in B_j\setminus S_j}|y_i| \leq c \sqrt{2\ln(p_j-s_j)}\Big)\;\leq\;P\Big(\tfrac{1}{\sqrt{2\ln(p_j-s_j)}}\max_{i \in B_j\setminus S_j}y_i \leq c \Big).
$$
On the other hand, results from extreme value theory\cite{galambos} Example 4.4.1, show that 
 $$\frac{1}{\sqrt{2 \ln(p_j-s_j)}}\max_{i\in B_j\setminus S_j}y_i \;\overset{P}{\longrightarrow} 1$$
 and so 
 $$
 P\Big(\frac{1}{\sqrt{2\ln(p_j-s_j)}}\max_{i \in B_j\setminus S_j}y_i \leq c \Big)\;\to\;0,
 $$ 
 which implies that $P(\hat{S}^{b} \subseteq S)\to 0$. 

\subsection{Proof of Proposition~\ref{prop:suffnecIIblockthresh}}
	 By the union bound, 
	$P\big(S \not\subseteq \hat{S}^{b}\big)\leq \sum_{j=1}^b  P\left(\min_{i\in S_j}|y_i/\sqrt{n}|\leq\tau_j\right)$. 
	By Lemma~\ref{lemma:3tailbounds} (ii), for each $j$,	
	$$P\left(\min_{i\in S_j}|y_i/\sqrt{n}|\leq\tau_j\right)\;\leq\; P\left(\max_{i\in S_j}|y_i/\sqrt{n}-\beta_i| \geq\beta_{\min,j}^{\star}-\tau_j \right).$$ By Lemma~\ref{lemma:3tailbounds} (i)  
	\begin{equation}\label{eq:type2convrate}
P\left(S \not\subseteq \hat{S}^{b} \right) \;\leq\;\sum_{j=1}^b\frac{e^{-\tfrac{n}{2}\left((\beta^\star_{\min,j}-\tau_j)^2-\tfrac{2\ln(s_j)}{n}\right)}}{\sqrt{\pi\ln(s_j)}}.		
	\end{equation}
 By A5, the nominator on the right is bounded above by $1$ for all sufficiently large $n$. It follows that, by A3 
	$$P\left(S \not\subseteq \hat{S}^{b} \right) \;\leq\;  \sum_{j=1}^b (\pi\ln (s_j))^{-1/2} \to 0\qquad \mbox{as } n\to \infty.$$ 

\subsection{Proof of Proposition~\ref{prop:necIIblockthresh}} 
(i) By independence we have 
$P(S \subseteq \hat{S}^{b})=\prod_{j=1}^b P(\min_{i \in S_j}|y_i/\sqrt{n}| > \tau_j)$. Take any $j$, satisfying \eqref{cond:necIIallblockthresh2}. $P(\min_{i \in S_j}|y_i/\sqrt{n}| > \tau_j)=\prod_{i \in B_j} P(|y_i/\sqrt{n}| > \tau_j)$. Denote $i_{\circ}\in B_j$ the entry, such that $\beta_i^*=\beta^*_{\min,\,j}$. From the proof of Lemma \ref{lemma:3tailbounds} (ii) $$P(|y_{i_{\circ}}/\sqrt{n}|>\tau_j) 
    = P\left(z>\sqrt{n}\left(\tau_j-\beta_{\min,\,j}^{\star}\right)\right) + P\left(z<-\sqrt{n}\left(\tau_j+\beta_{\min,\,j}^{\star}\right)\right)$$ where $z\sim N(0,1)$. Since $\tau_j<\beta^\star_{\min,j}$ satisfies condition \eqref{cond:necIallblockthresh} $P(z<-\sqrt{n}\left(\tau_j+\beta_{\min,\,j}^{\star}\right)) \to 0$. By \eqref{cond:necIIallblockthresh2} $\underset{n \to +\infty}{\lim}P\left(z>\sqrt{n}\left(\tau_j-\beta_{\min,\,j}^{\star}\right)\right)<1$ and $\underset{n \to +\infty}{\lim}P\left(S_j \subseteq \hat{S}^{b} \right)<1$.\\
(ii)	Take any $j=1,\ldots,b$ satisfying \eqref{cond:necIIallblockthresh}. By Lemma~\ref{lemma:3tailbounds} (iii), since $\beta^\star_{\min,j}>\tau_j$:
$$P\left(\min_{i \in S_j}\Big|\frac{y_i}{\sqrt{n}}\Big| > \tau_j \right)\;\leq\;\exp\left\{-\frac{s_j}{2}\frac{e^{-\tfrac{2n}{\pi}(\beta^\star_{\min,j}-\tau_j)^2}-e^{-\tfrac{n}{2}(\tau_j+\beta^\star_{\min,j})^2}}{\left(1-e^{-2n(\beta^\star_{\min,j}-\tau_j)^2/\pi}\right)^{\frac{1}{2}}+\left(1-e^{-n(\tau_j+\beta^\star_{\min,j})^2/2}\right)^{\frac{1}{2}}}\right\}. $$
Let $a_{n}=2\left(1-e^{-2n(\beta^\star_{\min,j}-\tau_j)^2/\pi}\right)^{\frac{1}{2}}+2\left(1-e^{-n(\tau_j+\beta^\star_{\min,j})^2/2}\right)^{\frac{1}{2}}$ and note that $a_n\in (0, 4]$.  Thus, to show that $P\left(\min_{i \in S_j}|\tfrac{y_i}{\sqrt{n}}| > \tau_j \right)$ is bounded away from $1$, it is enough to show that 
$$
\lim_{n\to \infty}{s_j}({e^{-2n(\beta^\star_{\min,j}-\tau_j)^2/\pi}-e^{-n(\tau_j+\beta^\star_{\min,j})^2/2}})\;>\;0.
$$
By \eqref{cond:necIIallblockthresh}, $\lim_{n\to\infty} {s_j} e^{-2n(\beta^\star_{\min,j}-\tau_j)^2/\pi}\geq 1$. To conclude, we show that
$$
\lim_{n\to\infty} {s_j} e^{-n(\tau_j+\beta^\star_{\min,j})^2/2}\;\leq\;c\;<\;1.
$$ 
for any $c\in(0,1)$ such that $\tfrac{s_j}{p_j}\leq c$. Such $c$ exists by our assumption $\tfrac{s_j}{p_j}<1$.
We equivalently need
$\sqrt{n}(\tau_j+\beta^\star_{\min,j})\geq\sqrt{{2\ln (s_j/c)}}$ for all $n$ sufficiently large. To show that, note that $\tau_j\geq \beta^\star_{\min,j}$ and $s_j/c\leq p_j$, which gives
$$
\frac{\tau_j+\beta^\star_{\min,j}}{\sqrt{\tfrac{2\ln(s_j/c)}{n}}}\;\geq\;\frac{2\tau_j}{\sqrt{\tfrac{2\ln(p_j)}{n}}}\;=\;\frac{\sqrt{\tfrac{2\ln(p_j-s_j)}{n}}}{\sqrt{\tfrac{2\ln(p_j)}{n}}}\frac{2\tau_j}{\sqrt{\tfrac{2\ln(p_j-s_j)}{n}}}\;=\;\sqrt{\frac{\ln(p_j-s_j)}{\ln(p_j)}}\frac{2\tau_j}{\sqrt{\tfrac{2\ln(p_j-s_j)}{n}}}.
$$
Since the second term on the right converges to something $\geq 2$, it is enough to show that the first term converges to something $>1/2$. Using that $\tfrac{s_j}{p_j}<c$, we get
$$
\sqrt{\frac{\ln(p_j-s_j)}{\ln(p_j)}}\;\geq\;\sqrt{\frac{\ln(p_j)+\ln(1-c)}{\ln(p_j)}}\;\to\;1,
$$
which concludes the proof.

\subsection{Proof of Lemma~\ref{lem:necbetaminblockthresh}}
    We have $\frac{\beta_{\min,j}^*}{\sqrt{\tfrac{2\ln(p_j-s_j)}{n}}+\sqrt{\tfrac{\pi}{2}\tfrac{\ln(s_j)}{n}}}= \frac{\beta_{\min,j}^*-\tau_j+\tau_j}{\sqrt{\tfrac{2\ln(p_j-s_j)}{n}}+\sqrt{\tfrac{\pi}{2}\tfrac{\ln(s_j)}{n}}}$ and \eqref{eq:necbetaminblockthresh} implies that $\lim_{n\to \infty}\frac{\tau_j}{\sqrt{\tfrac{2\ln(p_j-s_j)}{n}}} <1$ or $\lim_{n\to \infty}\frac{\beta_{\min,j}^*-\tau_j}{\sqrt{\tfrac{\pi}{2}\tfrac{\ln(s_j)}{n}}} <1$. 
    If $\lim_{n\to \infty}\frac{\tau_j}{\sqrt{\tfrac{2\ln(p_j-s_j)}{n}}} <1$ by Proposition~\ref{prop:suffnecIblockthresh}, $P(\hat{S}^{b} = S)\not \to 1$. If $\beta_{\min,j}<\tau_j$ recovery is not possible either by definition. 
    The last case to examine is when $\beta_{\min,j}\geq\tau_j$, $\lim_{n\to \infty}\frac{\tau_j}{\sqrt{\tfrac{2\ln(p_j-s_j)}{n}}} \geq 1$ and $\lim_{n\to \infty}\frac{\beta_{\min,j}^*-\tau_j}{\sqrt{\tfrac{\pi}{2}\tfrac{\ln(s_j)}{n}}} <1$. By Proposition~\ref{prop:necIIblockthresh}, in that case $P(S \subseteq \hat{S}^{b} )\not \to 1$ and recovery is not possible asymptotically.

\subsection{Proof of Theorem~\ref{theo:fullrecovGS}}
By the union bound and the bounds in \eqref{eq:type1convrate} and \eqref{eq:type2convrate}, under A4 and A5, we have:
\begin{equation}\label{eq:fhfhfhf}
P(S\neq \hat{S}^b) \; \leq\;\sum_{j=1}^b\frac{e^{-\tfrac{n}{2}\left(\tau_j^2-\tfrac{2\ln(p_j-s_j)}{n}\right)}}{\sqrt{\pi\ln(p_j-s_j)}}+\sum_{j=1}^b\frac{e^{-\tfrac{n}{2}\left((\beta^\star_{\min,j}-\tau_j)^2-\tfrac{2\ln(s_j)}{n}\right)}}{\sqrt{\pi\ln(s_j)}}.
\end{equation}
We continue with showing that if \eqref{eq:betamin} holds then $\tau^\star_j$ satisfies A4 and A5. Under \eqref{eq:betamin}, there exists $c>1$ such that $\beta^\star_{\min,j}=c(\sqrt{\tfrac{2\ln(p_j-s_j)}{n}} + \sqrt{\tfrac{2\ln(s_j)}{n}})$. Denote $a=\sqrt{\tfrac{2\ln(p_j-s_j)}{n}}$ and $b=\sqrt{\tfrac{2\ln(s_j)}{n}}$. We have
$$
\tau_j^\star\;=\;\frac{c}{2} (a+b)+\frac{(a^2-b^2)}{2c(a+b)}\;=\;\frac{c}{2} (a+b)+\frac{a-b}{2c}\;=\;\frac{c^2+1}{2c}a+\frac{c^2-1}{2c}b.
$$
The fact that $\tau^\star_j>a$ (the lower bound in A4) is now clear because $\tfrac{c^2+1}{2c}> 1$  and $\tfrac{c^2-1}{2c}> 0$. To show that the upper bound in A5 hold observe that
$$
\beta_{\min,j}^* - \tau^*_j \;=\; c(a+b)- \frac{c^2+1}{2c}a-\frac{c^2-1}{2c}b \;= \; \frac{c^2-1}{2c}a+\frac{c^2+1}{2c}b\;>\;  b 
$$
Then for $n$ sufficiently large \eqref{eq:rateconvgsm} holds. For convenience, denote $d = \frac{1}{n {\beta_{\min}^*}_j} \ln (p_j / s_j-1)$ such that $\tau^*_j= {\beta_{\min}^*}_j/2 + d$. Then
$$e^{-\frac{n}{2}\left({\tau^*_j}^2-\frac{2\ln(p_j-s_j)}{n}\right)}
=e^{-\left[\frac{n}{8}{\beta_{\min}^{\star}}_j^2+\frac{n}{2}d^2 + \frac{1}{2}\left(\ln (p_j- s_j)-\ln (s_j)\right) -\ln (p_j- s_j)\right]}$$
Then, by considering separately the two posssible maximum in $\max\{\ln(p_j-s_j),\ln(s_j)\}$, we get that:
$$\frac{e^{-\frac{n}{2}\left({\tau^*_j}^2-\frac{2\ln(p_j-s_j)}{n}\right)}}{e^{-[\frac{n}{8}{{\beta_{\min}^*}_j}^2-\max\{\ln(p_j-s_j),\ln(s_j)\}]}}= 
e^{-\left[\frac{n}{2}d^2 + \frac{1}{2}\left|\ln (p_j- s_j)-\ln (s_j)\right|\right]} <1$$
Note that $e^{-\tfrac{n}{2}\big[{\tau_j^*}^2-\tfrac{2\ln(p_j-s_j)}{n}\big]}=e^{-\tfrac{n}{2}\big[(\beta^\star_{\min,j}-{\tau_j^*})^2-\tfrac{2\ln(s_j)}{n}\big]}$, $\sqrt{\pi\ln(p_j-s_j)}>1$, and $\sqrt{\pi\ln(s_j)}>1$. Then from \eqref{eq:fhfhfhf} we have:
$$P(\hat{S}^{b} \neq S) \; \leq \; 2 \sum_{j=1}^be^{-[\frac{n}{8}{\beta^\star_{\min,j}}^2-\max\{\ln(p_j-s_j),\ln(s_j)\}]}.$$

\subsection{Proof of Corollary~\ref{cor:selconstimpossible}}

Since $\hat{S}$ is $\hat{S}^b$ with $b=1$, the assumptions of Lemma~\ref{lem:necbetaminblockthresh} for $\hat{S}$ are met and $P(\hat{S} = S)\not \to 1$. Since A4 and A5 hold, by Proposition~\ref{prop:suffnecIblockthresh} and Proposition~\ref{prop:suffnecIIblockthresh}, $P(\hat{S}^b \not\subseteq S)\to 0$ and $P(S \not\subseteq \hat{S}^b)\to 0$, and then $P(\hat{S}^b = S)\to 1$.

\section{Proofs of Section~\ref{sec:properties_legression}}
\subsection{Proof of Lemma~\ref{lem:useful-reformulation}}
	First note that by discarding constant term 
 $$\arg\max_{M\in \M} \left\{\max_{\beta\in \mathcal{L}_M}\ell(\bs y; \bs \beta)-\sum_{j=1}^b\kappa_j |M_j| \right\}=\arg\min_{M\in \M} \left\{\min_{\beta\in \mathcal{L}_M}\bigg\{ \tfrac{1}{2}\|\bs y-\bs X\bs\beta\|^2\right\}+\sum_{j=1}^b\kappa_j |M_j| \bigg\}$$
 
 We also have $\min_{\beta\in \mathcal{L}_M} \tfrac{1}{2}\|\bs y-\bs X\bs\beta\|^2=\tfrac{1}{2}\|\bs y-\bs X_{\!\!M}\tilde{\bs \beta}^{(M)}\|^2=\tfrac{1}{2}\|\bs y\|^2-\tfrac{1}{2}\|\bs X_{\!\!M}\tilde{\bs\beta}^{(M)}\|^2$. The maximization in \eqref{eq:Mhat} can be then replaced with maximization of $\tfrac{1}{2}
 \|\bs X_{\!\!M}\tilde{\bs\beta}^{(M)}\|^2-\sum_{j=1}^b \kappa_j|M_j|$.

\subsection{Proof of Lemma~\ref{lem:L0toL1convergence}}
The first statement follows directly from Lemma~\ref{lem:L0toL1convergencegen} by taking $\{M_1,\ldots,M_k\}=\{S\}$. For the second statement, note that  $NC(M)=(1+\sum_{N \neq M} e^{C(N)}/ e^{C(M)})^{-1}$. Since for any $M,N\in \M$ and $M\neq N$, $C(N) / C(M)>0$, $NC(M) \leq (1+e^{C(M')-C(M)})^{-1}$.

\subsection{Proof of Lemma~\ref{lemma:s7}}
    This result follows directly from Lemma~S7 in \cite{rossell2022concentration} taking $\phi^*=1$.

\subsection{Proof of Lemma~\ref{lem:boundrho}}
This result follows directly Lemma~\ref{lem:boundrhogen} by taking $T=S$, and observing that $\min_{i\in S_j\setminus M_j}{\beta^*_i}^2\geq {\beta^*_{\min,j}}^2$.

\subsection{Proof of Theorem~\ref{theo:suffcondlinearmodel}} We start by showing that, for any $M\in\mathcal{M}\setminus \{S\}$, $A_S$ as defined in Lemma~\ref{lem:boundexpnc} for $T=S$ satisfies the conditions of the lemma. Observe that $\Delta_{MS}$, defined in \eqref{eq:DeltaM}, can be rewritten as $\Delta_{MS}=\sum_{j=1}^b(|M_j\setminus S_j|-|S_j\setminus M_j|)\kappa_j$. By  Lemma~\ref{lem:boundrho}, for every $n\in \N$ we have 
\begin{equation}\label{eq:nonoverfittingcond}
   A_S\;=\; \gamma \Delta_{MS}+\frac{1-\gamma}{6}\mu_{Q_SM}\;\geq\; \gamma\sum_{j=1}^b|M_j \setminus S_j|\kappa_j + \sum_{j=1}^b |S_j\setminus M_j |\left(\tfrac{1-\gamma}{6} n \rho(\bs X) {\beta^{\star}_{\min,j}}^{2} - \gamma\kappa_j\right) 
\end{equation} 
Since $M\neq S$, $|M\setminus S|\neq0$ or $|S\setminus M|\neq0$, then by assumptions A6 and A7, for $n$ large enough $A_S>0$. We immediately have $\mu_{Q_S S}=o(A_S)$ because $\beta^*_{Q_S\setminus S}=\beta^*_{M\setminus S}=0$ and $\mu_{Q_S S}=0$. If $|M\setminus S|=0$, $|M\setminus S|=o(A_S)$ also immediately. Consider now the case $|M\setminus S|\neq 0$. By assumption A7, the last term in \eqref{eq:nonoverfittingcond} is nonnegative and:
$$\frac{|M\setminus S|}{A_S}\leq \bigg[ \gamma\sum_{j=1}^b\frac{|M_j \setminus S_j|}{|M\setminus S|}\kappa_j \bigg]^{-1} \leq \Big[ \gamma\min_{j=1,\ldots,b}\kappa_j \Big]^{-1} $$
where the last inequality follows from $\sum_{j=1}^b\frac{|M_j \setminus S_j|}{|M\setminus S|}=1$. Then by assumption A6, when $|M\setminus S|=o(A_S)$ when $|M\setminus S|\neq 0$ too.
We conclude that for every $M\in \mathcal M\setminus \{S\}$, by Lemma~\ref{lem:boundexpnc}, for any $\psi\in(0,1)$ and $n$ large enough:
$$
\E(NC(M))\;\leq\;e^{-\psi A_S}.
$$
Let 
$$A_S^*\;:=\;\gamma\sum_{j=1}^b|M_j \setminus S_j|\kappa_j + \sum_{j=1}^b |S_j\setminus M_j |\left(\tfrac{1-\gamma}{6} n \rho(\bs X) {\beta^{\star}_{\min,j}}^{2} - \gamma\kappa_j\right).$$
By \eqref{eq:nonoverfittingcond}, we have for $n$ large enough:
\begin{equation}\label{eq:boundexp}
    \E(NC(M))\;\leq\;e^{-\psi A_S^*}.
\end{equation}
Assumption~A7 implies there exist $g'_j \to +\infty$ such that:
$$\frac{(1-\gamma) n\rho(\bs X)}{6}{{\beta_{\min}^{\star}}_j}^2 \,- \,\kappa_j \;=\; \ln(s_j) + g^{\prime}_{j}.$$
Let $\delta \in (0,1)$ and denote $\bar{m}_j=\max\big\{\frac{2\ln(p_j-s_j)}{f_j},\frac{2\ln(s_j)}{g'_{j}}\big\}$. Take $\psi=\max_{j=1,\ldots,b}\frac{\xi + \delta + \bar{m}_j}{1+\bar{m}_j}$ for some $\xi \in(0, 1- \delta)$ then $\psi \in (0,1)$ and we have for every $j=1,\ldots,b$:
$$\delta \leq \psi-(1-\psi)\bar{m}_j.$$ 
By definition of $\gamma$ we have: 
$$
\gamma\kappa_j\;\geq\;\frac12 (1+\alpha_j)\kappa_j\;=\;\ln(p_j-s_j)+\frac12(\kappa_j-\ln(p_j-s_j)) \;=\; \ln(p_j-s_j)+ \frac12f_j.
$$
From our choice of $\psi$ and $\delta$, we get:
$$\psi\gamma\kappa_j \;\geq\; \psi\big(\ln(p_j-s_j)+ \frac12f_j\big) \;\geq\;  \ln(p_j-s_j)+ \delta\frac12f_j.$$
By A7 and by our choice of $\psi$ and $\delta$,
$$\psi\left(\tfrac{1-\gamma}{6} n \rho(\bs X) {\beta^{\star}_{\min,j}}^{2} - \gamma\kappa_j\right) \;\geq\; \psi\left(\ln(s_j)+g_j'\right) \;\geq\; \ln(s_j)+\delta \frac{1}{2}g'_j.
$$
Using the expression for $A_S^*$, the bound in \eqref{eq:boundexp}, and the above inequalities, we get
$$
\E(NC(M))\;\leq\;\exp\left\{- \sum_{j=1}^b |M_j\setminus S_j|(\ln(p_j-s_j)+\delta \tfrac{f_j}{2})-\sum_{j=1}^b |S_j\setminus M_j|(\ln(s_j)+\delta \tfrac{g'_j}{2})\right\}.
$$

Denote $\mathcal{S}= \sum_{M \in \M\setminus \{S\}} \E\left(NC(M)\right)$ for convenience. We can split $\mathcal{S}$ into sums over the models that have the same number of spurious variables and same number of missing variables in every block, that is the models $M$ such that for all $j$, $|M_j\setminus S_j|=u_j$ and $|S_j\setminus M_j|=w_j$ with $u_j \in \{0,\ldots,p_j-s_j\}$ and $w_j \in \{0,\ldots,s_j\}$, not all simultaneously zero:
$$\mathcal{S} =\sum_{w_1=0}^{s_1}\cdots\sum_{w_b=0}^{s_b}\sum_{u_1=0}^{p_1-s_1}\cdots\sum_{u_b=0}^{p_b-s_b}\sum_{M \neq S:\forall \,j |M_j\setminus S_j|=u_j, |S_j\setminus M_j|=w_j } \E\left(NC(M)\right)$$
The number of models with given collection of $u_j$'s and $w_j$'s is $\prod_{j=1}^b {\binom{p_j-s_j}{u_j}}{\binom{s_j}{w_j}}$. We thus have:
\begin{equation*}
    \mathcal{S} \leq \prod_{j=1}^b\left(1+\sum_{u_j=1}^{p_j-s_j}{\binom{p_j-s_j}{u_j}}e^{- u_j(\ln(p_j-s_j)+\delta \tfrac{f_j}{2})}\right)\left(1+
    \sum_{w_j=1}^{s_j}{\binom{s_j}{w_j}}
    e^{-w_j(\ln(s_j)+\delta \tfrac{g'_j}{2})}\right) - 1
\end{equation*}
where $-1$ stems from discarding $S$ from the product-sum, which corresponds to the case $u_j=0$ and $w_j=0$ for all $j$. A standard bound on binomial coefficient for $1\leq k\leq n$ is:
\begin{equation*}
    \binom{n}{k} \leq \left(\frac{n\,e}{k}\right)^k \leq \left(n\,e\right)^k
\end{equation*}
Then
\begin{equation*}
    \mathcal{S} \;\;\leq\;\; \prod_{j=1}^b\left(1+\sum_{u_j=1}^{p_j-s_j}e^{-u_j\left(\delta \tfrac{f_j}{2}-1\right)}\right)\left(1+
    \sum_{w_j=1}^{s_j}    e^{-w_j\left(\delta \tfrac{g'_j}{2}-1\right)}\right) - 1
\end{equation*}
Denote
$$
d_j\;=\;e^{1-\delta \tfrac{f_j}{2}},\qquad h_j\;=\;e^{1-\delta \tfrac{g'_j}{2}}
$$
where both expressions go to zero as $n$ increases.
For every $j$, we have:
\begin{eqnarray*}
    &1+\sum_{u_j=1}^{p_j-s_j}e^{-u_j\left(\delta \tfrac{f_j}{2}-1\right)}
    =\frac{1-d_j^{p_j-s_j+1}}{1-d_j} \\
    &1+\sum_{w_j=1}^{s_j}    e^{-w_j\left(\delta \tfrac{g'_j}{2}-1\right)} =\frac{1-h_j^{s_j+1}}{1-h_j}.
\end{eqnarray*}
Since both expressions converge to $1$ as $n$ grows, we conclude $\mathcal S\to 0$.

 \subsection{Proof of Theorem~\ref{prop:fullrecovreg}}
From the proof of Theorem~\ref{theo:suffcondlinearmodel} we have:
\begin{equation}\label{eq:prodconstreg}
    \mathcal{S} \;\;\leq\;\; \prod_{j=1}^b\left(1+\sum_{u_j=1}^{p_j-s_j}e^{-u_j\left(\delta \tfrac{f_j}{2}-1\right)}\right)\left(1+
    \sum_{w_j=1}^{s_j}    e^{-w_j\left(\delta \tfrac{g'_j}{2}-1\right)}\right) - 1
\end{equation}
For every $j$, we have:
\begin{eqnarray*}
    &\sum_{u_j=1}^{p_j-s_j}e^{-u_j\left(\delta \tfrac{f_j}{2}-1\right)}
    =g_j\,\frac{1-g_j^{p_j-s_j}}{1-g_j} \\
    &\sum_{w_j=1}^{s_j}    e^{-w_j\left(\delta \tfrac{g'_j}{2}-1\right)} =h_j\,\frac{1-h_j^{s_j}}{1-h_j}.
\end{eqnarray*}
Each of the sum above is smaller than 1 for $n$ sufficiently large. Then every term of degree larger than 1 in the product in \eqref{eq:prodconstreg} is bounded above by the sum of terms of degree 1. For $n$ large enough, we get: 
 \begin{equation*} \mathcal{S} \leq (2^{2b}-2b)\sum_{j=1}^b \bigg[g_j\,\frac{1-g_j^{p_j-s_j}}{1-g_j}+ h_j\,\frac{1-h_j^{s_j}}{1-h_j}\bigg ]. \end{equation*} 
 Take $$r_b=\max_{j=1,\ldots,b}\bigg\{\frac{1-g_j^{p_j-s_j}}{1-g_j} \,,\,\frac{1-h_j^{s_j}}{1-h_j} \bigg\}.$$
 By Lemma~\ref{lem:L0toL1convergence} and observing that $2e \leq 6$, we get: 
 \begin{eqnarray*}
     P(\hat{S}^{b} \neq S)\;&\leq&\;  2 \sum_{M \in \M\setminus \{S\}} \E\left(NC(M)\right) \\\;&\leq&\; (2^{2b}-2b)6r_b\,\,\sum_{j=1}^b\,\, e^{-\tfrac{\delta}{2}\big[\kappa_j-\ln(p_j-s_j)\big]} + e^{-\tfrac{\delta}{2}\big[\frac{(1-\gamma) n\rho(\bs X)}{2}{{\beta_{\min}^{\star}}_j}^2 - \kappa_j -\ln(s_j) \big]}
 \end{eqnarray*}
and 
 $$ P(\hat{S}^{b} \neq S)\;\leq \;(2^{2b}-2b)6r_b\sum_{j=1}^b e^{-\tfrac{\delta}{2}\big[\kappa_j-\ln(p_j-s_j)\big]} + e^{-\tfrac{\delta}{2}\big[(\sqrt{\frac{(1-\gamma) n\rho(\bs X)}{6}}{{\beta_{\min}^{\star}}_j} - \sqrt{\kappa_j})^2 -\ln(s_j) \big]}.$$

We continue with showing that if \eqref{betamincond:minimalreg} holds then $\kappa^\star_j$ satisfies A6 and A7. Under \eqref{betamincond:minimalreg}, there exists $c$ such that $\lim_{n\to\infty} c>1$ and $\sqrt{\frac{(1-\gamma)n}{6}}\beta^\star_{\min,j}=c\big(\sqrt{\ln(p_j-s_j)} + \sqrt{\ln(s_j)}\big)$. Denote $a=\sqrt{\ln(p_j-s_j)}$ and $b=\sqrt{\ln(s_j)}$. Proceeding as in the proof of Theorem~\ref{theo:fullrecovGS} shows that for $n$ large enough, we have $\sqrt{\kappa^*_j} \geq  \big(\frac{(c-1)^2}{2c}+1\big)\sqrt{\ln(p_j-s_j)}$ and $\frac{\sqrt{(1-\gamma) n\rho(\bs X)}}{7}{{\beta_{\min}^{\star}}_j}-\sqrt{\kappa^*} \geq \big(\frac{(c-1)^2}{2c}+1\big)\sqrt{\ln(s_j)}$
    which implies A6 and A7.
Then there exists $r_b\to 1$ such that for $n$ large enough the bound \eqref{eq:rateconvreg} on $P(\hat S^b \neq S)$ holds. Proceeding as in the proof of Theorem~\ref{theo:fullrecovGS} also shows that for any $j$ 
    $$e^{-\frac{\delta}{2}\left[{\kappa^*_j}-\ln(p_j-s_j)\right]}+e^{-\frac{\delta}{2}\big[(\sqrt{\frac{(1-\gamma) n\rho(\bs X)}{2}}{{\beta_{\min}^{\star}}_j} - \sqrt{\kappa_j^*})^2 -\ln(s_j) \big]} <2 e^{-\frac{\delta}{2}[\frac{(1-\gamma) n\rho(\bs X)}{24}{{\beta_{\min}^*}_j}^2-\max\{\ln(p_j-s_j),\ln(s_j)\}]}. $$
from which we conclude \eqref{eq:oraclerateconvreg}.

\subsection{Proof of Proposition~\ref{prop:nectypeIreg}}

Recovery of $S$ requires that $\max_{M\in O_j}\frac{NC(M)}{NC(S)}<1$ with prob-
ability going to 1 as n grows. By the definition of the normalized criterion $NC$:
$$P(\hat{S}^b=S) \leq P\Big(\max_{M\in O_j}\frac{NC(M)}{NC(S)}<1\Big)\leq  P\big(\max_{M\in O_j}\sqrt{L_{MS}}<\sqrt{2\kappa_j}\Big).$$
By Lemma~\ref{lemma:s7}, for every $M\in O_j$, $L_{MS}\in \chi^2_1$ and there exists $Z_M\sim N(0,1)$ such that $\sqrt{L_{MS}}=|Z_M|$. Since for every $M\in O_j$, $|Z_M|\geq Z_M$, we have:
\begin{equation}\label{eq:upperboundrecov}
    P(\hat{S}^b=S) \leq  1-P\big(\max_{M\in O_j}Z_M\geq\sqrt{2\kappa_j}\Big).
\end{equation}

The set $O_j$ has cardinality $p_j-s_j$, then by \cite{hartigan}, Theorem 3.4  and assumption A2, for any $\varepsilon>0$:
$$P\big(\max_{M\in O_j}Z_M \geq \underline{\lambda}_j \sqrt{2 \ln(p_j-s_j)}(1-\varepsilon)\big) \to 1.$$
If $\lim_{n\to\infty}\frac{\kappa_j}{\underline{\lambda}_j^2\ln(p_j-s_j)} <1$, then $P\big(\max_{M\in O_j}Z_M\geq\sqrt{2\kappa_j}\big) \to 1$ and by \eqref{eq:upperboundrecov} $P(\hat{S}^b=S) \to 0$.

\subsection{Proof of Lemma~\ref{lem:upperboundnoncentral}}
The result follows directly from Lemma~\ref{lem:upperboundnoncentralgen}, with $M=S$.

\subsection{Proof of Proposition~\ref{prop:nectypeIIreg}}

Let $M=S\setminus S^{S}(\kappa)$, since $S^{S}(\kappa)\neq \emptyset$, $M\neq S$. The event $(\hat{S}^b=S)$ implies $\big(\frac{NC(M)}{NC(S)}<1\big)$ then $P(\hat{S}^b=S)\leq P\big(\frac{NC(M)}{NC(S)}<1\big)$. By definition of $NC$, we have
$P\big(\frac{NC(M)}{NC(S)}<1\big) = P(L_{SM}> 2\Delta_{SM})$ where by Lemma~\ref{lemma:s7},  $L_{SM}\sim \chi_{|S\setminus M|}^2(\mu_{SM})$. We  bound next $P(L_{SM}> 2\Delta_{SM})$ with Lemma~\ref{lemma:s1s3}. We first show that $|S\setminus M|=o(2\Delta_{SM})$. and $\mu_{SM}=o(2\Delta_{SM})$. We have:
$$\frac{|S\setminus M|}{2\Delta_{SM}}\leq \bigg[\sum_{j=1}^b\frac{2|S_j \setminus M_j|}{|S\setminus M|}\kappa_j \bigg]^{-1} \leq \Big[ 2\min_{j=1,\ldots,b}\kappa_j \Big]^{-1} $$
where the last inequality follows from $\sum_{j=1}^b\frac{|S_j \setminus M_j|}{|S\setminus M|}=1$. By assumption, for all $j=1,\ldots,b$ $\kappa_j\to +\infty$ and the lef-hand side above goes to 0 as $n$ grows. By Lemma~\ref{lem:upperboundnoncentral}, we also have:
$$\frac{\mu_{SM}}{2\Delta_{SM}}\leq \frac{ \sum_{j=1}^b |S_j\setminus M_j|n\bar{\lambda}\max_{i\in S_j\setminus M_j}{\beta_i^*}^2}{\sum_{j=1}^b |S_j\setminus M_j|2 \kappa_j}.$$
Observe that:
\begin{eqnarray*}
    \sum_{j=1}^b\frac{n\bar{\lambda}\max_{i\in S_j\setminus M_j}{\beta_i^{\star}}^2}{2\kappa_j}
    &\geq &\sum_{j=1,|S_j\setminus M_j|\neq 0}^b\frac{|S_j\setminus M_j|n\bar{\lambda}\max_{i\in S_j\setminus M_j}{\beta_i^{\star}}^2}{|S_j\setminus M_j|2\kappa_j} \\
    &=& \frac{\sum_{j=1,|S_j\setminus M_j|\neq 0 }^b|S_j\setminus M_j|n\bar{\lambda}\max_{i\in S_j\setminus M_j}{\beta_i^{\star}}^2\Big(1+\sum_{l\neq j}\frac{|S_l\setminus M_l|2\kappa_l}{|S_j\setminus M_j|2\kappa_j}\Big)}{\sum_{j=1}^b|S_j\setminus M_j|2\kappa_j} \\
    &\geq & \frac{\sum_{j=1}^b|S_j\setminus M_j|n\bar{\lambda}\max_{i\in S_j\setminus M_j}{\beta_i^{\star}}^2}{\sum_{j=1}^b \big|S\setminus M_j\big|2\kappa_j}
\end{eqnarray*}
where last inequality follows from $\Big(1+\sum_{l\neq j}\frac{|S_l^S|\kappa_l}{|S_j^{S}(\kappa)|\kappa_j}\Big)\geq 1$ for all $j$ such that $|S_j^{S}(\kappa)|\neq 0$, and then
$$\frac{\mu_{SM}}{2\Delta_{SM}}\leq \sum_{j=1}^b\frac{n\bar{\lambda}\max_{i\in S_j\setminus M_j}{\beta_i^{\star}}^2}{2\kappa_j}.$$
For every $j=1,\ldots,b$, $S_j\setminus M_j \subset S_j^{S}(\kappa)$ , then by definition of the $S_j^{S}(\kappa)$, the right-hand 
side above goes to 0 and so does $\frac{\mu_{SM}}{2\Delta_{SM}}$. By Lemma~\ref{lemma:s7}, for any $\phi \in (0,1)$, for $n$ large enough:
\begin{equation}
    P(L_{SM}> 2\Delta_{SM}) \leq e^{-\phi2\Delta_{SM}}
\end{equation} 
where the right hand-side goes to 0 since for all $j=1,\ldots,b$, $\kappa_j \to +\infty$. 

\subsection{Proof of Corollary~\ref{cor:necrecovreg}}
If for some $j\in\{1,\ldots,b\}$ $\lim_{n\to\infty}\frac{\sqrt{n\bar{\lambda}}\beta_{\min,j}^{\star}}{\sqrt{\underline{\lambda}_j\ln(p_j-s_j)}} = 0$ then for that $j$ we have:
$$\lim_{n\to\infty}\frac{\sqrt{n\bar{\lambda}}\beta_{\min,j}^{\star}}{\sqrt{\kappa_j}}\sqrt{\frac{\kappa_j}{\underline{\lambda}_j\ln(p_j-s_j)}} = 0$$
If $\lim_{n\to\infty}\frac{\sqrt{n\bar{\lambda}}\beta_{\min,j}^{\star}}{\sqrt{\kappa_j}}>0$, then $\lim_{n\to\infty}\sqrt{\frac{\kappa_j}{\underline{\lambda}_j\ln(p_j-s_j)}} = 0$ and by Proposition~\ref{prop:nectypeIreg} $P(\hat{S}^b=S)\not\to1$. If $\lim_{n\to\infty}\frac{\sqrt{n\bar{\lambda}}\beta_{\min,j}^{\star}}{\sqrt{\kappa_j}}=0$, then by Proposition~\ref{prop:nectypeIIreg} $P(\hat{S}^b=S)\not\to1$.

\subsection{Proof of Corollary~\ref{cor:selconstimpossiblereg}}

Since $\hat{S}$ is $\hat{S}^b$ with $b=1$, the assumptions of XXXX for $\hat{S}$ are met and $P(\hat{S} = S)\not \to 1$. Since A6 and A7 hold, by XXXX $P(\hat{S}^b = S)\to 1$.

\subsection{Proof of Theorem~\ref{thm:convtoT}}
By definition of $\T$, for any $M \not\in\T$, there exists a unique $T^M\in\T$ such that $M\cap S^{I}(\kappa)=T^M\cap S^{I}(\kappa)$ and $M\neq T^M$. We next show that $A_{T^M}$ as defined in Lemma~\ref{lem:boundexpnc} satisfies the conditions of the lemma. Observe that $\Delta_{MT^M}$, defined in \eqref{eq:DeltaM}, can be rewritten as $\Delta_{MT^M}=\sum_{j=1}^b(|M_j\setminus T^M_j|-|T^M_j\setminus M_j|)\kappa_j$. By  Lemma~\ref{lem:boundrhogen}, for every $n\in \N$ we have 
\begin{equation}\label{eq:nonoverfittingcond2}
   \gamma \Delta_{MT^M}+\frac{1-\gamma}{2}\mu_{Q_{T^M}M}\;\geq\; \gamma\sum_{j=1}^b|M_j \setminus T^M_j|\kappa_j + \sum_{j=1}^b |T^M_j\setminus M_j |\Big(\tfrac{1-\gamma}{6} n \rho(\bs X) \min_{i\in |T^M_j\setminus M_j|}{\beta^{\star}_i}^{2} - \gamma\kappa_j\Big) 
\end{equation} 
Since $T^M \subseteq S^L(\kappa) \cup S^{I}(\kappa)$ and $M\cap S^{I}(\kappa)=T^M\cap S^{I}(\kappa)$, $T^M\setminus M \subseteq S^L(\kappa)$ and by definition of $S^L(\kappa)$, the rightmost component in \eqref{eq:nonoverfittingcond2} is nonnegative and positive if $|T^M\setminus M|\neq0$, for $n$ large enough. By assumption A6, component $\gamma\sum_{j=1}^b|M_j \setminus T^M_j|\kappa_j$ is nonnegative and positive if $|M\setminus T^M|\neq0$.  $M\neq T^M$ implies necessarily $|M\setminus T^M|\neq0$ or $|T^M\setminus M|\neq0$ and then, for $n$ large enough, $A_{T^M}>0$. If $|M\setminus T^M|=0$, we immediately have $|M\setminus T^M|=o(A_{T^M})$. If $|M\setminus T^M|\neq 0$, since the rightmost component in \eqref{eq:nonoverfittingcond2} is nonnegative, we have:
$$\frac{|M\setminus T^M|}{A_{T^M}}\leq \bigg[ \gamma\sum_{j=1}^b\frac{|M_j \setminus T^M_j|}{|M\setminus T^M|}\kappa_j \bigg]^{-1} \leq \bigg[ \gamma\min_{j=1,\ldots,b}\kappa_j \bigg]^{-1} $$
where the last inequality follows from $\sum_{j=1}^b\frac{|M_j \setminus T^M_j|}{|M\setminus T^M|}=1$. Then, by assumption A6, we have $|M\setminus T^M|=o(A_{T^M})$ when $|M\setminus T^M|\neq 0$ too. 
We have $M\setminus T^M\subseteq V\setminus (S^{I}(\kappa) \cup S^L(\kappa))$. If $M\cap S^{S}(\kappa)=\emptyset$, we immediately get $\mu_{Q_{T^M} T^M}=o(A_{T^M})$ because $\bs\beta^*_{Q_{T^M} \setminus T^M}=0$ and $\mu_{Q_{T^M} T^M}=0$. If $M\cap S^{S}(\kappa)\neq\emptyset$, by Lemma~\ref{lem:upperboundnoncentralgen}, we then have:
$$\frac{\mu_{Q_{T^M} T^M}}{ A_{T^M}} \leq \frac{n\,\bar{\lambda}\,\sum_{j=1}^b|(S_j\cap M_j )\setminus T^M_j|\,\max_{i\in (S_j\cap M_j )\setminus T^M_j}{\beta_{i}^{\star}}^2}{ \gamma\sum_{j=1}^b|M_j \setminus T^M_j|\kappa_j}.$$
For all $j=1,\ldots,b$, $|M_j \setminus T^M_j| \geq |(S_j\cap M_j )\setminus T^M_j|$. Moreover, since  $M\setminus T^M\subseteq V\setminus (S^{I}(\kappa) \cup S^L(\kappa))$ and $T^M\cap S^{I}(\kappa)=M\cap S^{I}(\kappa)$, $(S_j\cap M_j )\setminus T^M_j \subseteq S^{S}_j(\kappa)$ for all $j=1,\ldots,b$. We then get:
$$\frac{\mu_{Q_{T^M} T^M}}{ A_{T^M}} \leq \frac{n\,\bar{\lambda}\,\sum_{j=1}^b|(S_j\cap M_j )\setminus T^M_j|\,\max_{i\in S^{S}_j(\kappa)}{\beta_{i}^{\star}}^2}{ \gamma\sum_{j=1}^b|(S_j\cap M_j ) \setminus T^M_j|\kappa_j}.$$
Observe that:
\begin{eqnarray*}
    \sum_{j=1}^b\frac{n\bar{\lambda}\max_{i\in S^{S}_j(\kappa)}{\beta_i^{\star}}^2}{\gamma\kappa_j}
    &\geq &\sum_{j=1,|(S_j\cap M_j )\setminus T^M_j|\neq 0}^b\frac{|(S_j\cap M_j )\setminus T^M_j|n\bar{\lambda}\max_{i\in S^{S}_j(\kappa)}{\beta_i^{\star}}^2}{|(S_j\cap M_j )\setminus T^M_j|\gamma\kappa_j} \\
    &=& \frac{\sum_{j=1,|(S_j\cap M_j )\setminus T^M_j|\neq 0 }^b|(S_j\cap M_j )\setminus T^M_j|n\bar{\lambda}\max_{i\in S^{S}_j(\kappa)}{\beta_i^{\star}}^2\Big(1+\sum_{l\neq j}\frac{|(S_l\cap M_l)\setminus T^M_l|\gamma\kappa_l}{|(S_j\cap M_j )\setminus T^M_j|\gamma\kappa_j}\Big)}{\sum_{j=1}^b|(S_j\cap M_j )\setminus T^M_j|\gamma\kappa_j} \\
    &\geq & \frac{\sum_{j=1}^b|(S_j\cap M_j )\setminus T^M_j|n\bar{\lambda}\max_{i\in S^{S}_j(\kappa)}{\beta_i^{\star}}^2}{\sum_{j=1}^b \big|(S\cap M_j)\setminus T^M_j\big|\gamma\kappa_j} 
\end{eqnarray*}
where last inequality follows from $\Big(1+\sum_{l\neq j}\frac{|(S_l\cap M_l)\setminus T^M_l|\gamma\kappa_l}{|(S_j\cap M_j )\setminus T^M_j|\gamma\kappa_j}\Big)\geq 1$ for all $j$ and then,
$$\frac{\mu_{Q_{T^M} T^M}}{ A_{T^M}} \leq \sum_{j=1}^b\frac{n\bar{\lambda}\max_{i\in S^{S}_j(\kappa)}{\beta_i^{\star}}^2}{\gamma\kappa_j}.$$
By definition of the $S_j^{S}(\kappa)$, the right-hand side above goes to 0 as $n$ grows and $\mu_{Q_{T^M} T^M}=o(A_{T^M})$ when $M\cap S^{S}(\kappa)\neq\emptyset$ too. 
We conclude that for every $M\in \mathcal M\setminus \T$, by Lemma~\ref{lem:boundexpnc}, for any $\psi\in(0,1)$ and $n$ large enough:
$$
\E(NC(M))\;\leq\;e^{-\psi A_{T^M}}.
$$
Let 
$$A_{T^M}^*\;:=\;\gamma\sum_{j=1}^b|M_j \setminus T^M_j|\kappa_j + \sum_{j=1}^b |T^M_j\setminus M_j |\left(\tfrac{1-\gamma}{6} n \rho(\bs X) \min_{i\in S^{L}_j(\kappa)}{\beta^{\star}_i}^{2} - \gamma\kappa_j\right) $$
By \eqref{eq:nonoverfittingcond} and because $\min_{i\in |T^M_j\setminus M_j|}{\beta^{\star}_i}^{2}\geq \min_{i\in S^{L}_j(\kappa)}{\beta^{\star}_i}^{2}$, we have for $n$ large enough:
\begin{equation}\label{eq:boundexp2}
    \E(NC(M))\;\leq\;e^{-\psi A_{T^M}^*}.
\end{equation}
The definition of $S^{L}_j(\kappa)$ implies there exist $g'_j \to +\infty$ such that:
$$\frac{(1-\gamma) n\rho(\bs X)}{6}\min_{i\in S^{L}_j(\kappa)}{\beta^{\star}_i}^{2}\,- \,\kappa_j \;=\; \ln(s_j) + g'_{j}.$$
Let $\delta \in (0,1)$ and denote $\bar{m}_j=\max\big\{\frac{2\ln(p_j-s_j)}{f_j},\frac{2\ln(s_j)}{g'_{j}}\big\}$. Take $\psi=\max_{j=1,\ldots,b}\frac{\xi + \delta + \bar{m}_j}{1+\bar{m}_j}$ for some $\xi \in(0, 1- \delta)$ then $\psi \in (0,1)$ and we have for every $j=1,\ldots,b$:
$$\delta \leq \psi-(1-\psi)\bar{m}_j.$$ 
By definition of $\gamma$ we have: 
$$
\gamma\kappa_j\;\geq\;\frac12 (1+\alpha_j)\kappa_j\;=\;\ln(p_j-s_j)+\frac12(\kappa_j-\ln(p_j-s_j)) \;=\; \ln(p_j-s_j)+ \frac12f_j.
$$
From our choice of $\psi,\delta$, we get:
$$\psi\gamma\kappa_j \;\geq\; \psi\big(\ln(p_j-s_j)+ \frac12f_j\big) \;\geq\;  \ln(p_j-s_j)+ \delta\frac12f_j.$$
By A7 and by our choice of $\psi,\delta$,
$$\psi\left(\tfrac{1-\gamma}{6} n \rho(\bs X) \min_{i\in S^{L}_j(\kappa)}{\beta^{\star}_i}^{2} - \gamma\kappa_j\right) \;\geq\; \psi\left(\ln(s_j)+g_j'\right) \;\geq\; \ln(s_j)+\delta \frac{g'_j}{2}.
$$
Using the expression for $A_{T^M}^*$, the bound in \eqref{eq:boundexp2}, and the above inequalities, we get
$$
\E(NC(M))\;\leq\;\exp\left\{- \sum_{j=1}^b |M_j\setminus T^M_j|(\ln(p_j-s_j)+\delta \tfrac{f_j}{2})-\sum_{j=1}^b |T^M_j\setminus M_j|(\ln(s_j)+\delta \tfrac{g'_j}{2})\right\}.
$$

Denote $\mathcal{S}= \sum_{M \in \M\setminus \T} \E\left(NC(M)\right)$ for convenience. Recall that for $M \in \M\setminus \T$ there exists a unique $T^M\in\T$ such that $M\cap S^{I}(\kappa)=T^M\cap S^{I}(\kappa)$.  We can split $\mathcal{S}$ into sums over the models $M$ such that, for all $j$, $|M_j\setminus T^M_j|=u_j$ and $|T^M_j\setminus M_j|=w_j$ with $u_j \in \{0,\ldots,p_j-s_j,\ldots,p_j-s_j+|S_j^{S}(\kappa)|\}$ and $w_j \in \{0,\ldots,|S^{L}_j(\kappa)|\}$, not all simultaneously zero:
$$\mathcal{S} =\sum_{w_1=0}^{|S^L(\kappa)_1|}\cdots\sum_{w_b=0}^{|S^L(\kappa)_b|}\sum_{u_1=0}^{p_1-s_1+|S_1^S|}\cdots\sum_{u_b=0}^{p_b-s_b+|S_b^S|}\sum_{M \neq T^M:\forall \,j |M_j\setminus T^M_j|=u_j, |T^M_j\setminus M_j|=w_j } \E\left(NC(M)\right)$$
The number of models with given collection of $u_j$'s and $w_j$'s is $\prod_{j=1}^b {\binom{p_j-s_j+|S_j^{S}(\kappa)|}{u_j}}{\binom{|S_j^{L}(\kappa)|}{ w_j}}$. We thus have:
\begin{equation*}
    \mathcal{S} \leq \prod_{j=1}^b\left(1+\sum_{u_j=1}^{p_j-s_j+|S_j^{S}(\kappa)|}{\binom{p_j-s_j +|S_j^{S}(\kappa)|}{u_j}}e^{- u_j(\ln(p_j-s_j)+\delta \tfrac{f_j}{2})}\right)\left(1+
    \sum_{w_j=1}^{|S_j^{L}(\kappa)|}{\binom{|S_j^{L}(\kappa)|}{ w_j}}
    e^{-w_j(\ln(s_j)+\delta \tfrac{g'_j}{2})}\right) - 1.
\end{equation*}
where $-1$ stems from discarding the cases $M=T^M$ ($u_j=0$ and $w_j=0$ for all $j$). A standard bound on binomial coefficient for $1\leq k\leq n$ is:
\begin{equation*}
    \binom{n}{k} \leq \left(\frac{n\,e}{k}\right)^k \leq \left(n\,e\right)^k.
\end{equation*}
Then
\begin{equation*}
    \mathcal{S} \;\;\leq\;\; \prod_{j=1}^b\left(1+\sum_{u_j=1}^{p_j-s_j+|S_j^{S}(\kappa)|}e^{-u_j\left(\delta \tfrac{f_j}{2}-\ln\big(1+\frac{|S^{S}_j(\kappa)|}{p_j-s_j}\big)-1\right)}\right)\left(1+
    \sum_{w_j=1}^{|S_j^{L}(\kappa)|}    e^{-w_j\left(\ln\big(\frac{s_j}{|S_j^{L}(\kappa)|}\big)+\delta \tfrac{g'_j}{2}-1\right)}\right) - 1.
\end{equation*}
Denote
$$
d_j\;=\;e^{1+\ln\big(1+\frac{|S^{S}_j(\kappa)|}{p_j-s_j}\big)-\delta \tfrac{f_j}{2}},\qquad h_j\;=\;e^{1-\ln\big(\frac{s_j}{|S_j^{L}(\kappa)|}\big)-\delta \tfrac{g'_j}{2}}.
$$
Since $|S^{S}_j(\kappa)|=O(p_j-s_j)$ and $s_j\geq|S_j^{L}(\kappa)|$, $d_j$ and $h_j$ go to zero as $n$ increases.
For every $j$, we have:
\begin{eqnarray*}
    &1+\sum_{u_j=1}^{p_j-s_j+|S_j^{S}(\kappa)|}e^{-u_j\left(\delta \tfrac{f_j}{2}-1\right)}
    =\frac{1-d_j^{p_j-s_j+|S_j^{S}(\kappa)|+1}}{1-d_j} \\
    &1+\sum_{w_j=1}^{|S_j^{L}(\kappa)|}    e^{-w_j\left(\delta \tfrac{g'_j}{2}-1\right)} =\frac{1-h_j^{|S_j^{L}(\kappa)|+1}}{1-h_j}.
\end{eqnarray*}
Since both expressions converge to $1$ as $n$ grows, we conclude $\mathcal S\to 0$ which proves (i). Finally, (ii) follows from (i) and observing that by Lemma~\ref{lem:L0toL1convergencegen}, 
$P(\hat{S}^{b} \not\in \T ) \,\,\leq \,\, (2^{|S^{I}(\kappa)|}+1) \sum_{M \in \M\setminus \{\T\}} \E\left(NC(M)\right) = (2^{|S^{I}(\kappa)|}+1) \mathcal{S}$.


\section{Proofs of Section~\ref{sec:informed_l0}}

\subsection{Proof of Lemma~\ref{lem:empiricalbayesinterpret}}
Denote $H(\bs \theta)=p(\bs y \mid \bs \theta)$. For every $j=1,\ldots,b$, the partial derivative with respect to $\theta^{(j)}$ is:
    \begin{equation}\label{eq:logderivative}
        \frac{\partial \ln H(\bs \theta)}{\partial \theta^{(j)}}=\frac{\partial H(\bs \theta)}{\partial \theta^{(j)}}\,H(\boldsymbol{\theta})^{-1}
    \end{equation}
    Observe that:
    \begin{equation}\label{eq:partialmarglikelihood}
    H(\bs \theta)=\sum_{M \in \M}p(\bs y \mid M)p(M \mid \bs \theta)
    \quad \text{and}\quad 
    \frac{\partial H(\bs \theta)}{\partial \theta^{(j)}}=\sum_{M\in \M}p(\bs y \mid M)\frac{\partial p(M \mid \bs \theta)}{\partial \theta^{(j)}}
    \end{equation}
    By our choice of prior \eqref{eq:block_modelprior}:
    \begin{equation}\label{eq:partialprobmodel}
    \frac{\partial p(M \mid \bs \theta)}{\partial \theta^{(j)}}= p(M \mid \bs \theta) \bigg(\frac{\sum_{i\in B_j}m_i}{\theta^{(j)}}-\frac{p_j-\sum_{i\in B_j}m_i}{1-\theta^{(j)}} \bigg)        
    \end{equation}
    Replacing \eqref{eq:partialprobmodel} into \eqref{eq:partialmarglikelihood}, we get:
    \begin{equation*}\label{eq:partialmarglikelihood2}
        \frac{\partial H(\bs \theta)}{\partial \theta^{(j)}}=\frac{1}{\theta^{(j)}}\sum_{i \in B_j}\sum_{M\in \M: m_i=1}p(\bs y \mid M)p(M \mid \bs \theta) - \frac{1}{1-\theta^{(j)}}\bigg[p_j H(\bs \theta)-\sum_{i \in B_j}\sum_{M\in \M: m_i=1}p(\bs y \mid M)p(M \mid \bs \theta)
        \bigg]
    \end{equation*}
    Note that:
    $$\frac{\sum_{M\in \M: m_i=1}p(\bs y \mid M)p(M \mid \bs \theta)}{\sum_{M\in \M}p(\bs y \mid M)p(M \mid \bs \theta)} = P(m_i = 1\mid \bs y)$$
    By \eqref{eq:logderivative}, we then get the following expression of the partial derivative for every $j=1,\ldots,b$:
    $$\frac{\partial \ln H(\bs \theta)}{\partial \theta^{(j)}}=\frac{1}{\theta^{(j)}}\sum_{i \in B_j}P(m_i = 1\mid \bs y) - \frac{1}{1-\theta^{(j)}}\bigg[p_j-\sum_{i \in B_j}P(m_i = 1\mid \bs y)
        \bigg]$$
    Setting the partial derivatives to 0 and solving for $\theta^{(j)}$ gives the desired result.

\subsection{Proof of Proposition~\ref{prop:shatconsistence}}
We start by showing the lower bound. For every $j=1,\ldots,b$, we have the decomposition:
\begin{gather}\label{eq:decompshat}
    \frac{\widehat{s}_j}{p_j}
   =A_j\,+\, C_j \\
       \text{where} \quad A_j=\frac{1}{p_j}\sum_{i \in B_j}\sum_{M \in \T|i\in M}NC(M) \quad\text{and}\quad  C_j=\frac{1}{p_j}\sum_{i \in B_j}\sum_{M \in \M\setminus \T|i\in M}NC(M). \nonumber
\end{gather}
For every $j=1,\ldots,b$, $A_j$ can be furthered decomposed:
\begin{eqnarray*}
  A_j
   &=& \sum_{M \in \T}NC(M)\sum_{i \in B_j}\frac{I(i\in M_j)}{p_j}\\
   &=& \sum_{M \in \T}NC(M)\sum_{i \in S^{L}_j(\kappa)}\frac{I(i\in M_j)}{p_j}+\sum_{M \in \T}NC(M)\sum_{i \in B_j\setminus S^{L}_j(\kappa)}\frac{I(i\in M_j)}{p_j}\\
   &=& \frac{|S^{L}_j(\kappa)|}{p_j}\sum_{M \in \T}NC(M)+\sum_{M \in \T}NC(M)\sum_{i \in B_j\setminus S^{L}_j(\kappa)}\frac{I(i\in M_j)}{p_j}.
\end{eqnarray*}
where the last equality follows from $I(i\in M_j)=1$ for all $i\in S^L(\kappa)$ since $M\in\T$. The righmost term above and $C_j$ are nonnegative, we then get:
$$\E\Big(\frac{\widehat{s_j}}{p_j}\Big)\geq\frac{|S^{L}_j(\kappa)|}{p_j}\sum_{M \in \T}\E(NC(M))$$
and by Theorem~\ref{thm:convtoT} (i), $\lim_{n\to\infty}\E\Big(\frac{\widehat{s_j}}{p_j}\Big)\geq \frac{|S^{L}_j(\kappa)|}{p_j}$ for every $j=1,\ldots,b$.

We continue with the upper bound. Since that for all $M\in\T$, $I(i\in M_j)=0$ for all $i\in B_j\setminus(S^{L}_j(\kappa)\cup S^{I}(\kappa)_j)$, $A_j$ can be further written:
$$
  A_j
   = \frac{|S^{L}_j(\kappa)|}{p_j}\sum_{M \in \T}\E(NC(M))+\sum_{M \in \T}\E(NC(M))\sum_{i \in S^{I}(\kappa)_j}\frac{I(i\in M_j)}{p_j}.
$$
Using that $\sum_{i \in S^{I}(\kappa)_j}\frac{I(i\in M_j)}{p_j}\leq \frac{|S^{I}(\kappa)_j|}{p_j}$ for all $M$, we have:
\begin{eqnarray}\label{eq:uuuuu}
    \frac{\widehat{s}_j}{p_j}&\leq \frac{|S^{L}_j(\kappa)|+|S^{I}(\kappa)_j|}{p_j}\sum_{M \in \T}NC(M) + C_j.
\end{eqnarray}
A bound on $C_j$ for every $j=1,\ldots,b$ is:
\begin{equation}\label{eq:uuuuu2}
    C_j = \sum_{M \in \M\setminus \T}NC(M)\sum_{i \in B_j}\frac{I(i\in M_j)}{p_j}\leq \sum_{M \in \M\setminus \T}NC(M)
\end{equation}
where the inequality follows from $\sum_{i \in B_j}\frac{I(i\in M_j)}{p_j}\leq 1$ for all $M$.
Taking expectation in \eqref{eq:uuuuu} and \eqref{eq:uuuuu2} gives:
\begin{eqnarray*}
    \E\bigg(\frac{\widehat{s_j}}{p_j}\bigg) 
    &\leq \frac{|S^{L}_j(\kappa)|+|S^{I}(\kappa)_j|}{p_j}\sum_{M \in \T}\E(NC(M)) + \sum_{M \in \M\setminus \T}\E(NC(M)).
\end{eqnarray*}
 By Theorem~\ref{thm:convtoT} (i), $\sum_{M \in \T}\E(NC(M))\to 1$ and $\sum_{M \in \M\setminus\T}\E(NC(M))\to 0$, then $ \lim_{n\to+\infty}\E\big(\frac{\widehat{s_j}}{p_j}\big) \leq  \frac{|S^{L}_j(\kappa)|+|S^{I}(\kappa)_j|}{p_j}= \frac{s_j-|S^{S}_j(\kappa)|}{p_j}$ for every $j=1,\ldots,b$, which proves the upper bound.


\subsection{Proof of Theorem \ref{theo:empbayesselconsist}}
Denote for any $M\in\M$, $NC^{\circ}(M)$, the normalized criterion value for model $M$ under step 1 penalty $\kappa^{\circ}$. For this choice of penalty and every $j=1,\ldots,b$, we have:
\begin{equation*}
    \frac{\widehat{s_j}}{p_j}
   = \frac{1}{p_j}\sum_{M \in \M}NC^{\circ}(M)\sum_{i \in B_j}I(i\in M) =\sum_{M \in \M | M \neq S}\frac{|M_j|}{p_j}NC^{\circ}(M) +\frac{s_j}{p_j}NC^{\circ}(S) .
\end{equation*}
Using that $NC^{\circ}(S)=1-\sum_{M \in \M | M \neq S}NC(M)$, we get
\begin{equation*}
   \frac{\widehat{s}_j}{p_j}
   =\sum_{M \in \M | M \neq S}\frac{|M_j|-s_j}{p_j}NC^{\circ}(M) \,\,+\,\, \frac{s_j}{p_j}.
\end{equation*}
Splitting the sum in the right-hand side above between models $M$ that are overfitting in $B_j$ in the sense $|M_j|> s_j$ and those that are underfitting in $B_j$, in the sense $|M_j|<s_j$, we can write:
\begin{equation}\label{eq:rewriteshat}
  \frac{\widehat{s}_j}{p_j}  = \frac{s_j}{p_j}\,\,+\,\, O_j^{\circ}-\,\,U_j^{\circ}
\end{equation}
where:
$$
    O_j^{\circ} = \sum_{M \in \M | |M_j| > s_j}\frac{|M_j|-s_j}{p_j}NC^{\circ}(M)\quad\text{and}\quad
    U_j^{\circ} = \sum_{M \in \M | |M_j| < s_j}\frac{s_j-|M_j|}{p_j}NC^{\circ}(M).
$$
By \eqref{eq:rewriteshat}, we have the following decomposition of step 2 penalties :
\begin{eqnarray}\label{eq:decompstep2pen}
    \kappa^{EB}_j
    &=& \ln(p_j-s_j) + \ln\Big(\frac{\sqrt{n}}{s_j}\Big) + \ln\Big(\frac{p_j-\widehat{s}_j}{p_j-s_j}\Big) + \ln\Big(\frac{s_j}{\widehat{s}_j}\Big) \nonumber\\
    &=& \ln(p_j-s_j)+\ln\Big(\frac{\sqrt{n}}{s_j}\Big)+\ln\Big(1-\frac{p_j(O^{\circ}_j-U^{\circ}_j)}{p_j-s_j}\Big)+\ln\Big(\frac{s_j}{\widehat{s}_j}\Big).
\end{eqnarray}
We first show the $\kappa^{EB}_j$'s satisfy assumption~A6 with probability going to 1. Since $U_j^{\circ}$ is nonnegative, a lower bound on $\kappa^{EB}_j$ for every $j=1,\ldots,b$ is:
\begin{equation}\label{eq:decompstep2pen2}
    \kappa^{EB}_j
    \geq  \ln(p_j-s_j)+\ln\Big(\frac{\sqrt{n}}{s_j}\Big)+\ln\Big(1-\frac{p_j O^{\circ}_j}{p_j-s_j}\Big)+\ln\Big(\frac{s_j}{\widehat{s}_j}\Big).
\end{equation}
We have :
$$\frac{p_j O^{\circ}_j}{p_j-s_j} =  \sum_{M \in \M | |M_j| > s_j}\frac{|M_j|-s_j}{p_j-s_j}NC^{\circ}(M)\leq \sum_{M \in \M | |M_j| > s_j}NC^{\circ}(M).$$ 
If $M$ is such that $|M_j|>s_j$, then $M\not\in\T$ and by Theorem~\ref{thm:convtoT}, $O^{\circ}_j$ and $\ln\Big(1-\frac{p_j O^{\circ}_j}{p_j-s_j}\Big)$ vanish in probability. By assumption $\ln(\sqrt{n}s_j^{-1})\to+\infty$. Then, to show A6 holds it is enough to show that with probability going to 1 $\ln\Big(\frac{s_j}{\widehat{s}_j}\Big)$ is nonnegative. By \eqref{eq:uuuuu} and \eqref{eq:uuuuu2} in the proof of Proposition~\ref{prop:shatconsistence}, we have:
\begin{equation*}
    \frac{\widehat{s}_j}{p_j}\leq \frac{|S^{L}_j(\kappa)|+|S^{I}(\kappa)_j|}{p_j}\sum_{M \in \T}NC(M) + \sum_{M \in \M\setminus \T}NC(M).
\end{equation*}
By Theorem~\ref{thm:convtoT} (i), $\sum_{M \in \T}NC(M)$ and $\sum_{M \in \M\setminus\T}NC(M)$ converge in probability to 1 and 0 respectively. We then have with probability going to 1:
\begin{equation*}
    \frac{\widehat{s}_j}{p_j}\leq \frac{|S^{L}_j(\kappa)|+|S^{I}(\kappa)_j|}{p_j} \quad\text{and}\quad\ln\Big(\frac{s_j}{\widehat{s}_j}\Big)\geq \ln\bigg(\frac{s_j}{|S^{L}_j(\kappa)|+|S^{I}(\kappa)_j|} \bigg) \geq 0.
\end{equation*}
We conclude that we probability going to 1:
\begin{equation}\label{eq:eeeeeeeh}
    \kappa^{EB}_j
    \geq  \ln(p_j-s_j)+\ln\Big(\frac{\sqrt{n}}{s_j}\Big)
\end{equation}
and the $\kappa^{EB}_j$ satisfy assumption~A6.\medskip

We go on with showing that assumption~A9 implies A7 with probability going to 1. Observe:
$$\ln\Big(\frac{\widehat{s}_j}{s_j}\Big)=\ln\Big(1+\frac{p_j(O^{\circ}_j-U^{\circ}_j)}{s_j}\Big) .$$
Since $O_j^{\circ}$ is nonnegative, from \eqref{eq:decompstep2pen}, for every $j=1,\ldots,b$, an upper bound on $\kappa^{EB}_j$ is:
\begin{equation}\label{eq:gggggggg}
    \kappa^{EB}_j \leq \ln\big(p_j-s_j\big)+\ln\Big(\frac{\sqrt{n}}{s_j}\Big)+\ln\bigg(\frac{1+\frac{p_j}{p_j-s_j}U^{\circ}_j}{1-\frac{p_j}{s_j}U^{\circ}_j}\bigg)
\end{equation}
We split the sum in $U^{\circ}_j$ between models in $\T$ and those out of $\T$. Using that if $M\in \T$, $s_j-|M_j|\leq s_j-|S^{L}_j(\kappa^{\circ})|$ and if $M\not\in \T$, $s_j-|M_j|\leq s_j$, we have:
\begin{eqnarray*}
    U^{\circ}_j \leq \frac{s_j-|S^{L}_j(\kappa^{\circ})|}{p_j}\sum_{M \in \T | |M_j| < s_j}NC^{\circ}(M)+  \frac{s_j}{p_j}\sum_{M \in \M \setminus \T| |M_j| < s_j}NC^{\circ}(M)
\end{eqnarray*}
By Theorem~\ref{thm:convtoT}, $\sum_{M \in  \T| |M_j| < s_j}NC^{\circ}(M)$ and $\sum_{M \in \M \setminus \T| |M_j| < s_j}NC^{\circ}(M)$ converge in probability to 1 and 0 respectively. We then get than with probability going to 1:
\begin{eqnarray*}
    \frac{p_j}{p_j-s_j}U^{\circ}_j \leq \frac{s_j-|S^{L}_j(\kappa^{\circ})|}{p_j-s_j}\\
    \frac{p_j}{s_j}U^{\circ}_j \leq  \frac{s_j-|S^{L}_j(\kappa^{\circ})|}{s_j}
\end{eqnarray*}
By replacing $\frac{p_j}{p_j-s_j}U^{\circ}_j$ and $\frac{p_j}{s_j}U^{\circ}_j$ in \eqref{eq:gggggggg} by the bounds above, we get that  with probability going to 1:
\begin{equation}\label{eq:finalupperboundkappaj}
  - \,\sqrt{\kappa^{EB}_j} \geq  \,- \,\sqrt{\ln\big(p_j/|S^{L}_j(\kappa^{\circ})|-1\big)+\frac{1}{2}\ln\Big(n\Big) }
\end{equation}
By \eqref{eq:eeeeeeeh}, with probability going to 1:
\begin{equation}\label{eq:rrrrrrrrr}
    1-\frac12\Big(1+\max_{j=1, \ldots,b}\frac{\ln(p_j-s_j)}{\kappa^{EB}_j}\Big) \geq 1-\psi
\end{equation}
By \eqref{eq:finalupperboundkappaj} and \eqref{eq:rrrrrrrrr}, under assumption~A9, assumption~A7 holds with probability going to 1. Since A6 and A7 hold with probability going to 1, by Theorem \ref{theo:suffcondlinearmodel}, $P(\hat{S}^{EB,b}\neq S) \to 0$.

\subsection{Proof of Theorem \ref{theo:informedl0consist}} We have that $\kappa_j^{IB}=\kappa_j^{EB}+\ln(\widehat s_j)$ and by \eqref{eq:decompstep2pen2} in the proof of Theorem \ref{theo:empbayesselconsist}:
\begin{equation*}
    \kappa_j^{IB}
    \geq  \ln(p_j-s_j)+\ln(\sqrt{n})+\ln\Big(1-\frac{p_j O^{\circ}_j}{p_j-s_j}\Big).
\end{equation*}
From the proof of Theorem ~\ref{theo:empbayesselconsist}, by Theorem~\ref{thm:convtoT}, $\ln\Big(1-\frac{p_j O^{\circ}_j}{p_j-s_j}\Big)$ vanishes in probability. With probability going to 1, we then have:
\begin{equation}\label{eq:finallowerboundkappaj2}\kappa^{IB}_j 
    \geq  \ln(p_j-s_j)+\ln(\sqrt{n})
\end{equation}
and the $\kappa^{IB}_j$ satisfy assumption~A6.

Using that $\kappa_j^{IB}=\kappa_j^{EB}+\ln(\widehat s_j)$ and $O_j^{\circ}$ in \eqref{eq:decompstep2pen} is nonnegative, for every $j=1,\ldots,b$, an upper bound on $\kappa^{IB}_j$ is:
\begin{equation*}
    \kappa^{IB}_j \leq \ln\big(p_j-s_j\big)+\ln(\sqrt{n})+\ln\bigg(1+\frac{p_j}{p_j-s_j}U^{\circ}_j\bigg)
\end{equation*}
As in the proof of Theorem ~\ref{theo:empbayesselconsist}, by Theorem~\ref{thm:convtoT}, with probability going to 1:
\begin{equation}\label{eq:finalupperboundkappaj2}
    \frac{p_j}{p_j-s_j}U^{\circ}_j \leq \frac{s_j-|S^{L}_j(\kappa^{\circ})|}{p_j-s_j}\quad\text{and}\quad -\sqrt{\kappa_j} \geq -\sqrt{\ln\big(p_j-|S^{L}_j(\kappa^{\circ})|\big)+\frac{1}{2}\ln(n)}
\end{equation}
Moreover by \eqref{eq:finallowerboundkappaj2}, we have with probability going to 1:
\begin{equation}\label{eq:rrrrrrrrr2}
    1-\frac12\Big(1+\max_{j=1, \ldots,b}\frac{\ln(p_j-s_j)}{\kappa_j}\Big) \geq 1-\psi
\end{equation}
By \eqref{eq:finalupperboundkappaj2} and \eqref{eq:rrrrrrrrr2}  under assumption~A10, assumption~A7 holds with probability going to 1. Since A6 and A7 hold with probability going to 1, by Theorem \ref{theo:suffcondlinearmodel}, $P(\hat{S}^{IA,b}\neq S) \to 0$.

\section{Proofs of auxiliary results}\label{sec:proofauxi}

\subsection{Proof of Lemma~\ref{lem:l1isthreshold}}
    Denote $\hat{\bs \beta}$, the minimizer of the problem \eqref{eq:genl1}. The MLE under the full model is $\tilde{\bs\beta}=\tfrac{1}{\sqrt{n}} \bs y$ and the optimized function in \eqref{eq:genl1} can be then rewritten as
	$$
	\tfrac{1}{2}\bs y^{\top}\bs y +\tfrac{n}{2}\sum_{j=1}^b\sum_{i\in B_j}\left(\beta_i^2-2\tilde\beta_i\beta_i+2\tfrac{\lambda_j}{n}|\beta_i|\right)
	$$
	which we optimize with respect to each $\beta_i$ separately. For any $i\in B_j$ we have $\hat\beta_i=0$ if and only if $|\tilde \beta_i|\leq \frac{\lambda_j}{n}$. Similarly  for \eqref{eq:genl1ada}, denoting $\hat{\bs \beta}$ the minimizer of the problem, for any $i\in B_j$ we have $\hat\beta_i=0$ if and only if $|\tilde \beta_i|\leq \frac{\lambda_j}{|\beta^{\circ}_i|n}$. If $\bs \beta^{\circ}= \tilde{\bs\beta}$, then for any $i\in B_j$ $\hat{\beta}_i = 0$ if and only if $|\tilde \beta_i|\leq \sqrt{\tfrac{\lambda_j}{n}}$. If $\bs\beta^\circ$ is a LASSO estimate with penalization $\lambda^\circ$, then for any $i\in B_j$ 
 $\hat{\beta}_i = 0$ if and only if $\tilde{\beta}^2_i + \operatorname{sign}(\lambda^{\circ}/n-\tilde{\beta}_i) \frac{\lambda^\circ}{n}\tilde{\beta}_i-\frac{\lambda_j}{n} \leq 0$. Equivalently, for any $i\in B_j$ 
 $\hat{\beta}_i = 0$ if and only if $|\tilde{\beta}_ i| \leq  \frac{1}{2}\left(\lambda^\circ/n+\sqrt{(\lambda^\circ/n)^2+4 \lambda_j/n}\right)$.

\subsection{Proof of Lemma~\ref{lemma:3tailbounds}}

	(i) \,\, By union bound and the identical distribution of the $y_i$'s,
\begin{center}
    $P\Big(\underset{i\in\{1,\ldots,p\}}{\max}|y_i|> a \Big)     \leq \sum_{i=1}^{p}{ P\left(|\sqrt{n}y_i|>\sqrt{n}a \right)}
    = p\; P\left(|z|>\sqrt{n}a\right).$
\end{center}
    where $z\sim N(0,1)$. 
    By symmetry and the standard tail bound for standard normal, $P(z \geq \delta) \leq \frac{1}{\sqrt{2\pi}} \frac{1}{\delta} e^{-\delta^{2} / 2}$:
\begin{center}
  $P\Big(\underset{i\in\{1,\ldots,p\}}{\max}|y_i|> a \Big) 
     \leq  p\; 2 \, P\left(z>\sqrt{n}a\right) \leq \frac{1}{\sqrt{\pi}}\frac{\sqrt{2}}{\sqrt{n}a}p e^{-\frac{n}{2} a^{2}}.$
\end{center}
Since $a \geq \sqrt{\frac{2\ln(p)}{n}}$, $\frac{\sqrt{2}}{\sqrt{n}a}\leq \frac{1}{\sqrt{\ln(p)}}$. Taking $a^2=a^2-\frac{2\ln(p)}{n}+\frac{2\ln(p)}{n}$, we get \begin{center}
$ P\Big(\underset{i\in\{1,\ldots,p\}}{\max}|y_i|> a \Big) \leq \frac{1}{\sqrt{\pi\ln(p)}} e^{-\frac{n}{2} \left(a^{2}-\frac{2\ln(p)}{n}\right)}.$ \\\end{center}
b) $ \underset{i \in \{1,\ldots,s\}}{\min}|\beta_i|- \underset{i \in \{1,\ldots,s\}}{\max}|y_i-\beta_i| > a$ implies $\forall i \; \underset{i \in \{1,\ldots,s\}}{\min}|\beta_i| - |y_i-\beta_i| > a $. Since $\forall i \; |\beta_i|< |\beta_i-y_i|+|y_i|$ we have that $ \forall i \; \min_{i \in \{1,\ldots,s\}}|\beta_i| - |y_i-\beta_i| > a$ implies $\forall i \; |y_i| > a $ and $ \underset{i \in \{1,\ldots,s\}}{\min}|y_i| > a$. By the properties of the probability $P$: 
\begin{center}$ P\Big(\underset{i \in \{1,\ldots,s\}}{\min}|\beta_i|-\underset{i \in \{1,\ldots,s\}}{\max}|y_i-\beta_i| > a\Big) \leq P\Big(\underset{i \in \{1,\ldots,s\}}{\min}|y_i| > a\Big) $. \end{center}
c)\,\, We start by bounding the probability $P(|y_i|)>a$ for any $i$. We have that:
$$P(|y_i|>a) = P(y_i>a) + P(y_i<-a) = P(z>\frac{a-\mu}{\sigma}) + P(z<-\frac{a+\mu}{\sigma}) \quad \text{where }z\sim N(0,1).$$
If $\mu>0$, then $P(|y_i|>a) =  P(z>\frac{a-|\mu|}{\sigma}) + P(z<-\frac{a+|\mu|}{\sigma})$. If $\mu<0$, by symmetry, $P(z>\frac{a-\mu}{\sigma})=P(z>\frac{a+|\mu|}{\sigma})=P(z<-\frac{a+|\mu|}{\sigma})$ and $P(z<-\frac{a+\mu}{\sigma})=P(z<\frac{|\mu|-a}{\sigma})=P(z>\frac{a-|\mu|}{\sigma})$. Then for any $\mu$, $P(|y_i|>a)  = P(z>\frac{a-|\mu|}{\sigma}) + P(z<-\frac{a+|\mu|}{\sigma}) $ and using the symmetry of the standard Gaussian for any $a<|\mu|$: 
\begin{equation}\label{eq:random}
    P(|y_i|>a) = 1+P(0<z<\frac{|\mu|-a}{\sigma})-P(0<z<\frac{a+|\mu|}{\sigma}).
\end{equation}
From \cite{chu}, for any $\delta>0$, $\frac{1}{2}\left(1-e^{-\delta^2/2}\right)^{\frac{1}{2}}\leq P(0 <z < \delta)\leq \frac{1}{2}\left(1-e^{-2\delta^2/\pi}\right)^{\frac{1}{2}}$. Thus, for any $\gamma,\delta>0$:
$$ P(0<z \leq \gamma)-P(0<z \leq \delta)\leq
\frac{1}{2}\left(1-e^{-2\gamma^2/\pi}\right)^{\frac{1}{2}}-\frac{1}{2}\left(1-e^{-\delta^2/2}\right)^{\frac{1}{2}}.$$
Applying the above to \eqref{eq:random}, we get:
$$P(|y_i|>a) \leq 
1+ \frac{1}{2}\left(\left(1-e^{-2(|\mu|-a)^2/\pi\sigma^2}\right)^{\frac{1}{2}}-\left(1-e^{-(a+|\mu|)^2/2\sigma^2}\right)^{\frac{1}{2}}\right).\\$$

Using that $x^{\frac{1}{2}}-y^{\frac{1}{2}}=\frac{x-y}{x^{\frac{1}{2}}+y^{\frac{1}{2}}}$ for $x,y>0$, we get:
$$P(|y_i|>a) \leq 1- \frac{e^{-2(|\mu|-a)^2/\pi\sigma^2}-e^{-(a+|\mu|)^2/2\sigma^2}}{2\left(\left(1-e^{-2(|\mu|-a)^2/\pi\sigma^2}\right)^{\frac{1}{2}}+\left(1-e^{-(a+|\mu|)^2/2\sigma^2}\right)^{\frac{1}{2}}\right)}.$$

Finally, $P(\min_{i\in\{1,\ldots,s\}}|y_i|>a)=\exp\left(s\ln{P(|y_i|>a)}\right)$, using that $\ln\left(1+ x \right)< x$ for $x>-1$ we get the desired inequality.

\subsection{Proof of Lemma~\ref{lem:L0toL1convergencegen}}
Suppose $NC(M_1)+\ldots+NC(M_k)>\frac{k}{k+1}$ then for any $M\neq \{M_1,\ldots,M_k\}$, $NC(M)<\frac{1}{k+1}$. In addition, if $NC(M_1)+\ldots+NC(M_k)>\frac{k}{k+1}$ then necessarily $\max_{l=1,\ldots,k}NC(M_l)> \frac{1}{k+1}$ and $\hat{S}^{b}\in \{M_1,\ldots,M_k\}$. Consequently, 
$$
P(\hat{S}^b\not\in \{M_1,\ldots,M_k\})\; \leq\; P\Big(NC(M_1)+\ldots+NC(M_k)\leq\frac{k}{k+1}\Big).$$
Moreover, we have:
$$P\Big(NC(M_1)+\ldots+NC(M_k)\leq\frac{k}{k+1}\Big)\;=\;P\Big(\sum_{M \in \M\setminus \{M_1,\ldots,M_k\}} NC(M)\geq \frac{1}{k+1}\Big).$$
The result follows from the the Markov's inequality applied to the right-hand side above.

\subsection{Proof of Lemma~\ref{lem:boundrhogen}} 
The non-centrality parameter $\mu_{Q_TM}$, as defined in \eqref{eq:muQM}, satisfies
\begin{eqnarray*}
    \mu_{Q_TM}&=&\|\big(I_n-P_M\big) \bs X_{Q_T\setminus M} \bs \beta^{\star}_{Q_T\setminus M}\|^2\\
    &=&\|\big(I_n-P_M\big) \bs X_{T\setminus M} \bs \beta^{\star}_{T\setminus M}\|^2\\
    &= & n {\beta^{\star}_{T\setminus M}}^\top \left(\tfrac1n \bs X_{T\setminus M}^\top(I_n-P_M)\bs X_{T\setminus M}\right)\beta^{\star}_{T\setminus M}\\
    &\geq & n\lambda_{\min}\left(\tfrac1n \bs X_{T\setminus M}^\top(I_n-P_M)\bs X_{T\setminus M}\right)\|\beta^{\star}_{T\setminus M}\|^2,
\end{eqnarray*}
where the second equality follows from observing that $Q_T\setminus M=T\setminus M$. Since $T$ is a susbet of $S$, by reordering columns, $X_{S\setminus M}=[X_{T\setminus M},X_{S\setminus (M\cup T)}]$ and $\tfrac1n \bs X_{T\setminus M}^\top(I_n-P_M)\bs X_{T\setminus M}$ is a diagonal block of $\tfrac1n \bs X_{S\setminus M}^\top(I_n-P_M)\bs X_{S\setminus M}$. By interlacing, $\lambda_{\min}\left(\tfrac1n \bs X_{T\setminus M}^\top(I_n-P_M)\bs X_{T\setminus M}\right) \geq \lambda_{\min}\left(\tfrac1n \bs X_{S\setminus M}^\top(I_n-P_M)\bs X_{S\setminus M}\right)$. Finally observing that 
$\lambda_{\min}\left(\tfrac1n \bs X_{T\setminus M}^\top(I_n-P_M)\bs X_{T\setminus M}\right) \geq \rho(\bs X)$ and $\|\beta^{\star}_{T\setminus M}\|^2\geq \sum_{j=1}^b|T_j\setminus M_j|\,\min_{i\in T_j\setminus M_j}{\beta_{i}^{\star}}^2$, gives the desired result.

\subsection{Proof of Lemma~\ref{lemma:s1s3}}
By \cite{birge}, Lemma 8.1 we have for any $x>0$
$$P(W> (\nu+\mu)+2 \sqrt{(\nu+2 \mu) x}+2 x) \leq e^{-x}.$$
The function $f:x\mapsto (\nu+\mu)+2 \sqrt{(\nu+2 \mu) x}+2 x$ is  one-to-one between $\R^+$ and $( \nu+
\mu,+\infty)$. 
We have for any $w>\mu+\nu$:
$$P(W > w) \leq e^{-f^{-1}(w)} \:=\; e^{-\big(\frac{w+\mu}{2}  - \sqrt{2w(2 \mu+\nu)-2 \mu \nu - \nu^2 }\big) } .$$
Observe that:
$$\frac{w+\mu}{2}  - \sqrt{2w( 2 \mu+ \nu)-2 \mu \nu - \nu^2 }= \frac{w}{2}\bigg(1+\frac{\mu}{w} -\sqrt{\frac{8(2 \mu+\nu)}{w}\bigg(1-\frac{\nu\mu}{2w\mu+w\nu}-\frac{\nu^2}{2(2w\mu+w\nu)} \bigg)}\bigg).$$
 If $\nu=o(w)$ and $\mu=o(w)$, we then have $\frac{w+\mu}{2}  - \sqrt{2w( 2 \mu+ \nu)-2 \mu \nu - \nu^2 }= \frac{w}{2}(1+o(1))$ and for any $\phi\in(0,1)$ and $n$ large enough:
$$P(W > w) \leq  e^{-\phi\frac{w}{2}}$$

\subsection{Proofs of Lemma~\ref{lemma:s2}}
The result follows directly from \cite{rossell2022concentration}, Lemma S2.

\subsection{Proof of Lemma~\ref{lemma:s18}}
For any $u\in [\underline{u},\bar{u}]$, we have $2\ln \left(\frac{g}{1 / u-1}\right) \geq 2\ln \left(\frac{g}{1 / \underline{u}-1}\right)$. Since $\underline{u}\geq (1+g^{\phi}e^{-(\nu+\mu)/2})^{-1}$,  we also have for any $u\in (\underline{u},\bar{u})$:
$$2\ln \left(\frac{g}{1 / u-1}\right) \geq 2(1-\phi)\ln(g)+\nu+\mu .$$
By assumption $\nu=o\left(\ln (g)\right)$ and $\mu=\left(\ln (g)\right)$, then for any $u\in (\underline{u},\bar{u})$, $\nu=o\left(2\ln \left(\frac{g}{1 / u-1}\right)\right)$ and $\mu=o\left(2\ln \left(\frac{g}{1 / u-1}\right)\right)$. By Lemma~\ref{lemma:s1s3}, for any $\phi\in(0,1)$ and $n$ large enough:
\begin{equation}\label{eq:randomm}
    \int_{\underline{u}}^{\bar{u}} P\left(W>2\ln \left(\frac{g}{1 / u-1}\right)\right) d u< \frac{1}{g^{\phi}} \int_{\underline{u}}^{\bar{u}}(1 / u-1)^{\phi} d u .
\end{equation}
Applying the change of variables $v=1 / u-1$ to the integral above gives
$$
\int_{\underline{u}}^{\bar{u}}(1 / u-1)^{\phi} d u=\int_{1 / \bar{u}-1}^{1 / \underline{u}-1} \frac{v^{\phi}}{(v+1)^2} d v.
$$
Note that $(1+x)^r \leq 1+r x$  for $x \geq-1$ and $r\in[0,1]$ and that $v^{\phi}=(v-1 + 1)^{\phi}$ with $v-1>-1$ since $v<1 / \bar{u}-1$ and $\bar{u}<1$. Then $v^{\phi}=(v-1 + 1)^{\phi}\leq 1 + \phi (v-1)\leq 1 + \phi(v+1)$. As a result:
\begin{equation}\label{eq:random2}
\int_{\underline{u}}^{\bar{u}}(1 / u-1)^{\phi} d u<\int_{1 / \bar{u}-1}^{1 / \underline{u}-1} \frac{1}{(v+1)^2} +\frac{\phi}{v+1} d v = \bar{u}-\underline{u}+\phi\ln\Big(\frac{\bar{u}}{\underline{u}}\Big)
\end{equation}
The results follows inputing the bound from \eqref{eq:random2} in \eqref{eq:randomm} and using that $\phi<1$.


\subsection{Proof of Lemma~\ref{lemma:s20}}
Since a probability is bounded by $1$, for any $a\in(0,1)$.
\begin{equation}\label{eq:basebound}
    \int_0^1 P\left(W>2 \ln \left(\frac{g}{1 / u-1}\right)\right) d u \;\leq\; 2 a+\int_{a}^{1-a} P\left(W>2\ln\left(\frac{g}{1 / u-1}\right)\right) d u.
\end{equation}

Take $a=\left(1+g^{\phi}e^{-(\nu+\mu)/ 2}\right)^{-1}$ for some $\phi\in(\alpha,1)$. By Lemma~\ref{lemma:s18} with $\underline{u}=a$ and $\bar{u}=1-a$:
$$
\int_0^1 P\left(W>2 \ln \left(\frac{g}{1 / u-1}\right)\right) d u \;\leq\; \frac{2}{1+g^{\phi}e^{-(\nu+\mu)/2}}+ 
\frac{1-2a+\ln(g^{\phi}e^{-(\nu+\mu)/2})}{g^{\phi}} $$
Since $1+g^{\phi}e^{-(\nu+\mu)/ 2}>g^{\phi}e^{-(\nu+\mu)/ 2}$ and $1-2a-\frac{\nu+\mu}{2}\leq \ln(g^\phi)$ for $n$ large enough, we have:
$$
\int_0^1 P\left(W>2 \ln \left(\frac{g}{1 / u-1}\right)\right) d u \;\leq\; g^{-\phi}\big(2e^{(\nu+\mu)/ 2}+2\ln(g^{\phi})\big) $$
We have:
\begin{equation}\label{eq:random24}
    \frac{g^{-\phi}2e^{(\nu+\mu)/ 2}}{g^{-\alpha}} \;=\;e^{-(\phi-\alpha)\ln(g)\big(1-\frac{\nu+\mu}{2(\phi-\alpha)\ln(g)}- \frac{\ln(2)}{(\phi-\alpha)\ln(g)}\big)}
\end{equation}
and 
\begin{equation}\label{eq:random25}
    \frac{g^{-\phi}2\ln(g^{\phi})}{g^{-\alpha}} \;=\;e^{-(\phi-\alpha)\ln(g)\big(1-\frac{\ln(\phi\ln(g))}{(\phi-\alpha)\ln(g)}- \frac{\ln(2)}{(\phi-\alpha)\ln(g)}\big)}
\end{equation}
Since $\alpha<\phi$, $g$ is increasing and $\nu=o(\ln(g))$ and $\mu=o(\ln(g))$, both expressions in \eqref{eq:random24} and \eqref{eq:random25} vanish as $n$ grows and for $n$ large enough:
$$
\int_0^1 P\left(W>2 \ln \left(\frac{g}{1 / u-1}\right)\right) d u \;\leq\; g^{-\alpha}$$

\subsection{Proof of Lemma~\ref{lem:boundexpnc}}
Since $M\in \mathcal M\setminus \{T\}$, by Lemma~\ref{lem:L0toL1convergence} (ii),  $NC(M)\leq (1+e^{C(T)-C(M)})^{-1}$.  By the standard formula for expectation of a nonnegative random variable in terms of its CDF, we get:
\begin{eqnarray*}
  \E\left(NC(M)\right)  &\leq& \int_0^1 P\left((1+e^{C(T)-C(M)})^{-1}\geq u\right){\rm d} u\\
  &=&\int_0^1 P\left(C(T)-C(M) \leq \ln{(\tfrac1u-1)}\right){\rm d}u\\
  &=&\int_0^1 P\Big(-\tfrac12 L_{TM} \geq \Delta_{MT}-\ln{(\tfrac1u-1)}\Big) {\rm d}u.
\end{eqnarray*}
Let $Q_T=M \cup T$ and recall that $L_{TM}=L_{Q_TM}-L_{Q_TT}$. For any $\gamma\in (\frac{1}{2},1]$
$$
-\tfrac12 L_{TM} -(\Delta_{MT}-\ln{(\tfrac1u-1)})\;=\;\left(\tfrac12 L_{Q_TT}-\ln{\Bigl(\frac{e^{\gamma\Delta_{MT}}}{\tfrac1u-1}\Bigl)}\right)-\Bigl(\tfrac12 L_{Q_TM}+(1-\gamma)\Delta_{MT}\Bigl).
$$
Observe that for any random variables $U$, $V$, and any $\epsilon,\gamma' \geq 0$, the event $
\{U-V \geq 0\}$ implies $\{U \geq \gamma' \epsilon\}\cup
\{V  < \gamma' \epsilon\}$. Let $U=\tfrac12 L_{Q_TT}-\ln{\Bigl(\frac{e^{\gamma\Delta_{MT}}}{\tfrac1u-1}\Bigl)}$ and $V=\tfrac12 L_{Q_TM}+(1-\gamma)\Delta_{MT}$. Take $\epsilon=\mu_{Q_TM}$ and $\gamma'=\tfrac16({1-\gamma})$. Recall that $A_T=\gamma\Delta_{MT}+\gamma'\mu_{Q_TM}$ and
$$
\{U\geq \gamma'\epsilon\}\;=\;\{\tfrac12 L_{Q_TT}\geq \ln{\Bigl(\frac{e^{\gamma\Delta_{MT}}}{\tfrac1u-1}}\Bigl)+\gamma'\mu_{Q_TM}\}\;=\;\{\tfrac12 L_{Q_TT}\geq \ln{\Bigl(\frac{e^{A_T}}{\tfrac1u-1}}\Bigl)\},
$$
$$
\{V<\gamma'\epsilon\}\;=\;\{\tfrac12 L_{Q_TM}<-(1-\gamma)\Delta_{MT}+\gamma'\mu_{Q_TM}\}\;=\;\{\tfrac12 L_{Q_TM}<\gamma'(\mu_{Q_TM}-6\Delta_{MT})\}.
$$
By the union bound we have:
\begin{equation}\label{eq:prooflemmas20nonspu}
    \E\left(NC(M)\right)  \leq \int_0^1P\Bigl(\tfrac12 L_{Q_TT}\geq \ln{\Bigl(\frac{e^{A_T}}{\tfrac1u-1}\Bigl)}\Bigl){\rm d}u+ P\Bigl(\tfrac12 L_{Q_TM}<\gamma'(\mu_{Q_TM}-6\Delta_{MT})\Bigl).
\end{equation}
We first bound the first term in \eqref{eq:prooflemmas20nonspu}. If $M \subset T$ then $Q_T=T$, $L_{Q_TT}=0$, and this term is zero. Suppose now that $T$ is a strict subset of $Q_T$. By Lemma~\ref{lemma:s7}, $L_{Q_TT}$ has a $\chi^2$-distribution. More precisely, $L_{Q_TT} \sim \mathcal{\chi}_{|Q_T\setminus T|}^2(\mu_{Q_TT})$ with $|Q_T\setminus T|=|M\setminus T|$. By assumption, $A_T>0$, $|M\setminus T|=o(\ln(e^{A_T}))$ and $\mu_{Q_TT}=o(\ln(e^{A_T}))$, then  by Lemma~\ref{lemma:s20}, for $\alpha\in(\psi,1)$, and $n$ large enough:
\begin{equation}\label{eq:boundLQS}
\int_0^1P\Bigl(L_{Q_TT}>2\ln{\Bigl(\frac{e^{A_T}}{1 / u-1}\Bigl)}\Bigl){\rm d}u \;<\; 
e^{-\alpha A_T}.
\end{equation}
We now bound the second term in \eqref{eq:prooflemmas20nonspu}. If $T \subset M$, then $Q_T=M$, $L_{Q_TM}=0$, $\mu_{Q_TM}=0$, and this term is zero. If $M$ is a strict subset of $Q_T$ then, by Lemma~\ref{lemma:s7}, $L_{Q_TM} \sim \chi_{|Q_T\setminus M|}^2(\mu_{Q_TM})$ with $|Q_T\setminus M|=|T\setminus M|$.
Clearly, when $\mu_{Q_TM}\leq 6\Delta_{MT}$, this term is also zero so suppose  $\mu_{Q_TM}> 6\Delta_{MT}$.  
We have  
$$P\left(L_{Q_TM}<2\gamma'(\mu_{Q_TM}-6\Delta_{MT})\right)\; \leq \; P\left(L_{Q_TM}<2\gamma'\mu_{Q_TM}\right)$$ 
and, by Lemma~\ref{lemma:s2}, as $\gamma'\in (0,\tfrac{1}{12})$
$$P\left(L_{Q_TM}<2\gamma'\mu_{Q_TM}\right) \;\;\leq\;\; (2\gamma')^{\frac{|T\setminus M|}{4}} e^{-\frac{1}{2}(1-\sqrt{2\gamma'})^2\mu_{Q_TM}}\;\;\leq\;\;  (\tfrac12)^{\tfrac34|T\setminus M|}e^{-\frac{1}{2}(1-\sqrt{2\gamma'})^2\mu_{Q_TM}}.$$
Since $\mu_{Q_TM}>6\Delta_{MT}$, we get
$$
A_T\;=\;\gamma'\mu_{Q_TM}+\gamma\Delta_{MT}\;<\;\frac{1-\gamma}{6}\mu_{Q_TM}+\frac{\gamma}{6}\mu_{Q_TM}\;=\;\frac16\mu_{Q_TM}\;\leq\;\frac12(1-\sqrt{2\gamma'})^2\mu_{Q_TM},
$$
where the last inequality follows from the fact that $\gamma'\in(0, \tfrac{1}{12})$. We conclude:
\begin{equation}\label{eq:boundLQM}
    P\left(L_{Q_TM}<2\gamma'(\mu_{Q_TM}-8\Delta_{MT})\right) \;\leq\; \left(\frac12\right)^{\tfrac34|T\setminus M|}e^{-A_T} \leq e^{-\alpha A_T}
\end{equation}
Summing the bounds in \eqref{eq:boundLQS} and \eqref{eq:boundLQM} gives that for $n$ large enough $\E\left(NC(M)\right) < 2e^{-\alpha A_T}$. Since $\psi<\alpha$, for $n$ large enough $\E\left(NC(M)\right) < e^{-\psi A_T}$.

\subsection{Proof of Lemma~\ref{lem:upperboundnoncentralgen}}
By Lemma~\ref{lemma:s7}, we have
\begin{eqnarray*}
    \mu_{Q_TT}&=&{\bs \beta^{\star}_{Q_T\setminus T}}^{\top}\bs X_{Q_T\setminus T}^{\top}  \big(I_n-P_T\big) \bs X_{Q_T\setminus T} \bs \beta^{\star}_{Q_T\setminus T}\\
    &=&{\bs \beta^{\star}_{M\setminus T}}^{\top}\bs X_{M\setminus T}^{\top}  \big(I_n-P_T\big) \bs X_{M\setminus T} \bs \beta^{\star}_{M\setminus T}\\
    &=&{\bs \beta^{\star}_{(S\cap M)\setminus T}}^{\top}\bs X_{(S\cap M)\setminus T}^{\top}  \big(I_n-P_T\big) \bs X_{(S\cap M)\setminus T} \bs \beta^{\star}_{(S\cap M)\setminus T}
\end{eqnarray*}
where the third equality follows from $\bs \beta^*_{M\setminus S}=0$.
We start with the upper bound. Denote for any square matrix $A$, its largest eigenvalue 
$\lambda_{\max}(A)$. We have:
$$\mu_{Q_TT}\leq n\lambda_{\max}\Big(\frac{1}{n}\bs X_{(S\cap M)\setminus T}^{\top}  \big(I_n-P_T\big) \bs X_{(S\cap M)\setminus T}\Big) \|\bs \beta^{\star}_{(S\cap M)\setminus T}\|^2$$
Let $B=\tfrac{1}{n}\bs X_{(S\cap M)\setminus T}^{\top}  \big(I_n-P_T\big) \bs X_{(S\cap M)\setminus T}$ and $C=\tfrac{1}{n}\bs X_{(S\cap M)\cup T}^{\top}\bs X_{(S\cap M)\cup T}$. $B$ is the Schur complement $\big(C \,  / \, \tfrac{1}{n}\bs X_T^{\top} \bs X_T\big)$. The inverse $B^{-1}$ is then a diagonal block of $C^{-1}$. By interlacing. $\lambda_{min}(B^{-1})\geq \lambda_{min}(C^{-1})$ and then $\lambda_{\max}(B)\leq \lambda_{\max}(C)$. We also have $(S\cap M)\cup T \subseteq S$, then by interlacing again $\lambda_{\max}(C)\leq \lambda_{\max}\Big(\tfrac{1}{n}\bs X_S^{\top}\bs X_S\Big)=\bar{\lambda}$. The upper bound follows from also observing that $\|\bs \beta^{\star}_{(S\cap M)\setminus T}\|_2^2\leq \sum_{j=1}^b |(S_j\cap M_j)\setminus T_j|\max_{i\in (S_j\cap M_j)\setminus T_j}{\beta_i^*}^2$.

A lower bound on $\mu_{Q_TT}$ is: 
$$\mu_{Q_TT}\geq n\lambda_{\min}\Big(\frac{1}{n}\bs X_{(S\cap M)\setminus T}^{\top}  \big(I_n-P_T\big) \bs X_{(S\cap M)\setminus T}\Big) \|\bs \beta^{\star}_{(S\cap M)\setminus T}\|_2^2$$
The final bound is obtained by similar interlacing argument as for the upper bound and the bound $\|\bs \beta^{\star}_{(S\cap M)\setminus T}\|_2^2\geq \sum_{j=1}^b |(S_j\cap M_j)\setminus T_j|\min_{i\in (S_j\cap M_j)\setminus T_j}{\beta_i^*}^2$

\section{Properties in the Gaussian sequence model with $s$ finite}\label{sec:appendixGSMsfinite}

We derive here selection properties of the block $\ell_0$ penalties in the Gaussian Sequence model dropping assumption~A3 and assuming instead:\\
(A11) \quad for all $j$, \;$s_j \,< \, \infty$.

Changing A3 for A11 implies redeveloping results relative to the probability of false negatives and consequently betamin conditions. Observe that the Lemmas and Properties of Section~\ref{sec:GSMgen} as well as Property \ref{prop:suffnecIblockthresh} in Section~\ref{sec:properties_legression} do not depend on A3, the probability of false negatives or a betamin condition. They hold equally in the regime where $s_j$ is finite for all $j=1,\ldots,b$.

Consider the assumption:\\

(A12) \quad for all $j$, \; $\sqrt{n}(\beta_{\min,j}^{\star}-\tau_j) \to \infty$.\\

Under A11, A12 can be shown to be sufficient for $\hat{S}^{b}$ to hold the screening property. Moreover in Proposition~\ref{prop:necIIblockthresh} a) we showed A12 was a general necessary condition for the screening property without assuming A3. A12 is then necessary and sufficient for $\hat{S}^{b}$ to hold the property under A11.
\begin{prop}\label{prop:suffnecIIblockthreshbis}
 In the Gaussian sequence model \eqref{eq:GSM}, under assumptions A1, A2 and A12, $\hat{S}^{b}$ satisfies:
$$P(S\subseteq \hat{S}^{b} ) \to 1$$
\end{prop}
\begin{proof}
    By the proof of Proposition~\ref{prop:suffnecIblockthresh}, we have: \begin{equation*}
P\left(\hat{S}^{b} \not\subseteq S\right)\;\leq\; \sum_{j=1}^b\frac{e^{-\tfrac{n}{2}\left((\beta^\star_{\min,j}-\tau_j)^2-\tfrac{2\ln(s_j)}{n}\right)^2}}{\sqrt{\pi\ln(s_j)}}
 \end{equation*}
 where the right-hand side goes to 0 by assumption A12.
\end{proof}
By Proposition~\ref{prop:suffnecIblockthresh}, to guarantee model recovery with high probability the block thresholds $\tau_j$ must be greater or equal to $\sqrt{2\ln(p_j-s_j)/n}$. Under A11, we then get the minimal betamin condition, for all $j=1,\ldots,b$:
\begin{equation}\label{eq:nectypeIIGSMfinite}
    \sqrt{n}\beta_{\min,j}^{\star}-\sqrt{2\ln(p_j-s_j)} \to \infty
\end{equation}
Condition \eqref{eq:nectypeIIGSMfinite} is necessary for asymptotic support recovery with $\hat{S}^{b}$:
\begin{lem}\label{lem:necbetaminblockthresh2}
    Suppose there exists $j\in\{1,\ldots,b\}$ such that
    \begin{equation}\label{eq:necbetaminblockthreshbis}
      \lim_{n\to \infty}\sqrt{n}\beta_{\min,j}^*-\sqrt{2\ln(p_j-s_j)} < +\infty  
    \end{equation}
then under assumptions A1, A2 and A11,\;  $P(\hat{S}^{b} = S)\not \to 1$
\end{lem}

\begin{proof}

  We have $\sqrt{n}\beta_{\min,j}^*-\sqrt{2\ln(p_j-s_j)} \,=\, \sqrt{n}\beta_{\min,j}^*-\sqrt{n}\tau_j+\sqrt{2\ln(p_j-s_j)}\Big(\tfrac{\sqrt{n}\tau_j}{\sqrt{2\ln(p_j-s_j)}}-1\Big)$ and \eqref{eq:necbetaminblockthreshbis} implies that there exists $c\in\R^+$ such that:
  \begin{equation}\label{eq:proofnecbetaminbis}
      \lim_{n\to \infty}\sqrt{n}\beta_{\min,j}^*-\sqrt{n}\tau_j+\sqrt{2\ln(p_j-s_j)}\Big(\tfrac{\sqrt{n}\tau_j}{\sqrt{2\ln(p_j-s_j)}}-1\Big)\leq c
  \end{equation}

Consider the case $\lim_{n\to +\infty}\tfrac{\sqrt{n}\tau_j}{\sqrt{2\ln(p_j-s_j)}}-1<0$. Then by Proposition~\ref{prop:suffnecIblockthresh},
$P(\hat{S}^{b} = S)\not \to 1$. 
Now consider the case $\lim_{n\to +\infty}\tfrac{\sqrt{n}\tau_j}{\sqrt{2\ln(p_j-s_j)}}-1 \geq0$. Then by \eqref{eq:proofnecbetaminbis},  $\lim_{n\to \infty}\sqrt{n}\beta_{\min,j}^*-\sqrt{n}\tau_j \leq c$ and by Proposition~\ref{prop:necIIblockthresh} a), $P(\hat{S}^{b} = S)\not \to 1$.
\end{proof}

Theorem~\ref{theo:fullrecovGS} does not assume A3 but requires the earlier betamin condition~A5. Now under A11, A12 implies A5. As a result, when $s_j$ is finite for all $j=1,\ldots,b$, Theorem~\ref{theo:fullrecovGS} holds equally assuming A11 and A12 instead of A5.

It follows that the gains in terms of convergence rate with  informed block thresholding under A11 are the same as under A3. To discuss gains in terms of betamin conditions, we assume without loss of generality that $\beta_{\min}^{\star}$ is located in block $b$. By \eqref{eq:nectypeIIGSMfinite}, the minimal betamin conditions with block informed thresholding $\hat{S}^{b}$ and standard thresholding $\hat{S}$ are respectively:
\begin{eqnarray}
    \sqrt{n}\beta_{\min,b}^{\star}-\sqrt{2\ln(p_b-s_b)} \to \infty \quad\text{and}\quad
    \sqrt{n}\beta_{\min}^{\star}-\sqrt{2\ln(p-s)} \to \infty
\end{eqnarray}
The relevant quantities to compare the minimal requirement on the growth rate of betamin are thus $\sqrt{\tfrac{\ln(p_b-s_b)}{n}}$ and $\sqrt{\tfrac{\ln(p-s)}{n}}$. Since $p_b-s_b \leq p - s$, the minimal betamin condition with $\hat{S}^{b}$ is always milder than with $\hat{S}$ under A11 again. Moreover, observe that under Examples 1 to 3 in Section~\ref{sec:examples}, the relevant ratio to compare minimal betamin condition was also $\Theta\Big(\sqrt{\tfrac{\ln(p_b-s_b)}{\ln(p-s)}}\Big)$. As a result, those examples also illustrate the magnitude of the gains in terms of minimal betamin conditions under A11.


\section{Non linear block $\ell_0$ penalty in high-dimensional linear regression}\label{sec:appendixnonlinear}
\sdescription{Short description of Supplement C.}

In this section we derive selection properties of the block $\ell_0$ penalties in the linear regression setting without assuming linearity. 

Let $\eta_j$ for $j=1,\ldots,b$ be any non-negative and increasing penalty functions on the natural numbers. The selector based on those $\ell_0$ block penalties is:
\begin{equation}
  \widehat S^b\;\in\;\arg\max_{M\in \M} \left\{\max_{\beta\in \mathcal{L}_M} \ell(\bs y;\bs \beta)-\sum_{j=1}^b\eta_j (|M_j|)\right\},\qquad M=M_1\cup \ldots\cup M_b.	
\end{equation}
where $\ell$ is the log-likelihood function.

Rewrite the difference in penalty between any model $M$ and the true model $S$:
\begin{equation}\label{eq:DeltaMbis}
\Delta_{MS} \; := \; \sum_{j=1}^b \eta_j(|M_j|)-\eta_j(|S_j|)
\end{equation}
%The proof of Lemma~\ref{lem:boundexpnc} does not rely on the linearity of block penalties and the Lemma holds identically with the redefined $\Delta_{M}$. 
We define the normalized difference in block penalty between any model $M$ and the true support $S$ as:
 $$\tilde{\kappa}_{M, j}\; := \; \begin{cases}
     \frac{\eta_j(|M_j|)-\eta_j(|S_j|)}{|M_j|\,-\,|S_j|} \;\;\; \text{if}\;\;\; |M_j|\,\neq\,|S_j| \\
     0 \;\;\;\text{if}\;\;\; |M_j|\,=\,|S_j|
 \end{cases}  $$
The quantity $\tilde{\kappa}_{M, j}$ plays a similar role as $\kappa_j$, representing the average penalty incurred for adding a variable from block $B_j$ to the model $M$. However, unlike in the case of linear penalties, $\tilde{\kappa}_{M, j}$ is not constant but a function of $|M_j|$ and $|S_j|$. For any $M\in\M$, $\tilde{\kappa}_{M, j}$ is exactly $\kappa_j$ if the penalty is linear as defined in Section \ref{sec:intro_blockpen}.

As in the linear case, consistency is achieved by setting conditions on the $\tilde{\kappa}_{M, j}$ and $\beta^*_{\min, j}$:\\
(AXXX) for each $j$, there exist $f_j\to \infty$ (as $n\to \infty$) such that for all $M \in \M$ such that $|M_j| \neq |S_j|$ and $|M_j \setminus S_j| >0$, for $n$ sufficiently large, and  all $j=1,\ldots,b$:
\begin{equation}\label{cond:suffIregnonlinpen}
    \tilde{\kappa}_{M, j} =\;\ln\bigg(\frac{p_j-s_j}{|M_j \setminus S_j| }\bigg) + f_j
\end{equation}

(AXXX) for each $j$, there exist $h_j\to +\infty$ such that for $n$ sufficiently large, for all $j=1,\ldots,b$
    \begin{equation}\label{cond:suffIIregnonlinpen}
        \sqrt{\frac{(1-\gamma) n\rho(\bs X)}{2}}{{\beta_{\min}^{\star}}_j} \,- \,\sqrt{\ln(p_j-s_j) + f_j} \;=\; \sqrt{\ln(s_j)} + h_j .
    \end{equation}
where $\alpha_j:=\frac{\ln(p_j-s_j)}{\ln(p_j-s_j) + f_j}$ and $\gamma:=\frac12(1+\max_j \alpha_j)\in (\tfrac12,1)$ for $n$ large enough.

\begin{thm}\label{theo:suffcondlinearmodelnonlinpen}
Suppose that A1, A2, A7 and A10 hold, then 
    % \begin{eqnarray}
    %  &\qquad \tilde{\kappa}_{M, j} > \ln(p_j-s_j)  \label{cond:suffIregnonlinpen}\\
    %    &\qquad \sqrt{\frac{2(1-\gamma) n\rho(\bs X)}{8}}{{\beta_{\min}^{\star}}_j}  \geq \sqrt{\tilde{\kappa}_{M, j}}+{\sqrt{\ln(s_j)}}  .\label{cond:suffIIregnonlinpen}
    % \end{eqnarray}
Then $$\sum_{M \in \M\setminus \{S\}}\E\left(NC(M)\right)\to 0 \qquad \mbox{and} \qquad P(\widehat{S}^b = S) \to 1$$
\end{thm}

\begin{proof}
The proof is essentially the same as proof Theorem~\ref{theo:suffcondlinearmodel} replacing $\kappa_j$ by $\tilde{\kappa}_{M, j}$ for every $j=1,\ldots,b$. We rewrite here the main intermediate results of the proof Theorem~\ref{theo:suffcondlinearmodel} and redevelop only the steps altered by the non-linearity of the penalty. We refer the reader to the proof of Theorem~\ref{theo:suffcondlinearmodel} for more details.

Observe that $\Delta_{MS}$, defined in \eqref{eq:DeltaMbis}, can be rewritten as $\Delta_{MS}=\sum_{j=1}^b(|M_j\setminus S_j|-|S_j\setminus M_j|)\tilde{\kappa}_{M}$. By  Lemma~\ref{lem:boundrho}, for every $n\in \N$ we have 
\begin{equation}\label{eq:nonoverfittingcondgen}
   A_S\;=\; \gamma \Delta_{MS}+\frac{1-\gamma}{2}\mu_{Q_SM}\;\geq\; \gamma\sum_{j=1}^b|M_j \setminus S_j|\tilde{\kappa}_{M} + \sum_{j=1}^b |S_j\setminus M_j |\left(\tfrac{1-\gamma}{2} n \rho(\bs X) {\beta^{\star}_{\min,j}}^{2} - \gamma\tilde{\kappa}_{M, j}\right) 
\end{equation} 
As for Theorem~\ref{theo:suffcondlinearmodel}, by assumption A10 we have $\gamma\sum_{j}|M_j\setminus S_j|\tilde{\kappa}_{M, j}\geq \tfrac{3}{4}|M\setminus S|$ as long as $\min_j \ln(p_j-s_j)\geq \tfrac32$, which always holds in our asymptotic regime for $n$ sufficiently large. Observe that the proof of Lemma~\ref{lem:boundexpnc} does not rely on the linearity of block penalties and the Lemma holds identically with the redefined $\Delta_{MS}$.  By Lemma~\ref{lem:boundexpnc} we have:
$$
\E(NC(M))\;\leq\;16 e^{-A_S}A_S^{|M\setminus S|/2+1}+e^{-A_S},
$$
where $A_S=\gamma\Delta_{MS}+\tfrac{1-\gamma}{2}\mu_{Q_SM}$. Similarly to Theorem~\ref{theo:suffcondlinearmodel}, defining:
$$A_S^*\;:=\;\gamma\sum_{j=1}^b|M_j \setminus S_j|\tilde{\kappa}_{M, j} + \sum_{j=1}^b |S_j\setminus M_j |\left(\tfrac{1-\gamma}{2} n \rho(\bs X) {\beta^{\star}_{\min,j}}^{2} - \gamma\tilde{\kappa}_{M, j}\right) $$
we get 
$$
\E(NC(M))\;\leq\;16 e^{-A_S^*}{A_S^*}^{|M\setminus S|/2+1}+e^{-A_S^*}.
$$
Note that 
$$A_S^*\;=\;\gamma\sum_{j \in \{1,\ldots, b\} \text{s.t.} |M_j \setminus S_j| \neq 0}^b|M_j \setminus S_j|\tilde{\kappa}_{M, j} + \sum_{j=1}^b |S_j\setminus M_j |\left(\tfrac{1-\gamma}{2} n \rho(\bs X) {\beta^{\star}_{\min,j}}^{2} - \gamma\tilde{\kappa}_{M, j}\right) $$
Furthermore, condition \eqref{cond:suffIIregnonlinpen} implies there exist $g'_j \to +\infty$ such that:
$$\frac{(1-\gamma) n\rho(\bs X)}{2}{{\beta_{\min}^{\star}}_j}^2 \,- \,\tilde{\kappa}_{M, j} \;=\; \ln(s_j) + g'_{j}.$$
We also have that for $x$ sufficiently large and any $\psi \in (0,1)$
\begin{equation}\label{eq:boundexp3}
16e^{-x}x^{\tfrac12|M\setminus S|+1}+e^{-x}\leq e^{-\psi x}.    
\end{equation}
To choose $\psi$ we proceed as in the proof of Theorem~\ref{theo:suffcondlinearmodel}. Let $\delta \in (0,1)$ and denote:
$$\bar{m}_j=\begin{cases}
    \max\bigg\{\tfrac{2\ln\big(\tfrac{p_j-s_j}{|M_j\setminus S_j|}\big)}{f_j},\tfrac{2\ln(s_j)}{g'_{j}}\bigg\},\;\;\; \text{if}\;\;|M_j \setminus S_j| \neq 0 \\
    \frac{2\ln(s_j)}{g'_{j}},\;\;\; \text{if}\;\;|M_j \setminus S_j| \neq 0
\end{cases}$$
Take $\psi=\max_{j=1, \ldots,b}\frac{a + \delta + \bar{m}}{1+\bar{m}}$ for some $a \in(0, 1- \delta)$ then $\psi \in (0,1)$ and for every $j=1,\ldots,b$:
$$\delta \leq \psi-(1-\psi)\bar{m}_j.$$ 
By definition of $\gamma$ and $\alpha_j$, we have for every $j=1,\ldots,b$ such that $|M_j \setminus S_j|\neq0$:
$$
\gamma\tilde{\kappa}_{M, j}\;\geq\;\tfrac12 (1+\alpha_j)\tilde{\kappa}_{M, j}\geq \tfrac12 \bigg(1+\tfrac{\ln\big(\tfrac{p_j-s_j}{|M_j\setminus S_j|}\big)}{\tilde{\kappa}_{M, j}}\bigg)\tilde{\kappa}_{M, j} \;=\; \ln\big(\tfrac{p_j-s_j}{|M_j\setminus S_j|}\big)+ \tfrac12f_j.
$$
From our choice of $\psi,\delta$, we further get:
$$\psi\gamma\tilde{\kappa}_{M, j} \;\geq\; \psi\bigg(\ln\big(\tfrac{p_j-s_j}{|M_j\setminus S_j|}\big)+ \frac12f_j\bigg) \;\geq\; \ln\big(\tfrac{p_j-s_j}{|M_j\setminus S_j|}\big)+ \delta\frac12f_j.$$
Observe that, for any $M$, $\tilde{\kappa}_{M, j} \leq \ln(p_j-s_j)+f_j$. Then, by \eqref{cond:suffIIregnonlinpen} and our choice of $\psi,\delta$:
$$\psi\left(\tfrac{1-\gamma}{2} n \rho(\bs X) {\beta^{\star}_{\min,j}}^{2} - \gamma\tilde{\kappa}_{M, j}\right) \;\geq\; \psi\left(\ln(s_j)+g_j'\right) \;\geq\; \ln(s_j)+\delta \frac{g'_j}{2}.
$$
Using the expression for $A_S^*$, the bound in \eqref{eq:boundexp3}, and the above inequalities, we get
$$
\E(NC(M))\;\leq\;\exp\{- \sum_{j \in \{1,\ldots, b\} \text{s.t.} |M_j \setminus S_j| > 0} |M_j\setminus S_j|\bigg(\ln\bigg(\frac{p_j-s_j}{|M_j\setminus S_j|}\bigg)+\delta \tfrac{f_j}{2}\bigg)-\sum_{j=1}^b |S_j\setminus M_j|(\ln(s_j)+\delta \tfrac{g'_j}{2})\}.
$$
Following the same subsequent steps of the proof of Theorem~\ref{theo:suffcondlinearmodel} shows selection consistency.
\end{proof}

The EBIC penalty is equivalent to considering that all the variables belong to a single block, that is $b=1$, and setting $\eta_1(|M|)=\psi\ln{\binom{p}{|M|}} + \tfrac{1}{2}\ln(n)$. The EBIC penalty can be shown to satisfy assumption~A10 whenever $s=o(\ln(n))$ and the following corollary holds:
\begin{cor}
    Suppose that $b=1$, $\eta_1(|M|)=\psi \ln {\binom{p}{|M|}} + \tfrac{1}{2}\ln(n)$ with $\psi\geq1$, $s=o(\sqrt{n})$, $\max|M|_{M\in\M} \leq \tfrac{p}{2}$ and that there exists $c\to +\infty$ such that for $n$ sufficiently large
    \begin{equation}\label{eq:betaminebic}
        \sqrt{\frac{\ln(n)-2(s-2)}{2 \ln(p)+\ln(n)-2(s-2)}\frac{n\rho(\bs X)}{16}}{{\beta_{\min}^{\star}}} \,- \,\sqrt{\psi\ln(p) + \frac{1}{2}\ln(n)} \;=\; \sqrt{\ln(s)} + c .
    \end{equation}
    Then A10 and A11 hold and under A1 and A2: $$\sum_{M \in \M\setminus \{S\}}\E\left(NC(M)\right)\to 0 \qquad \mbox{and} \qquad P(\widehat{S}^b = S) \to 1$$    
\end{cor}

\begin{proof}
    We start by showing $\eta_1$ satisfies assumption~A10. Let $M\in\M$ and denote $|M\setminus S|=u$ and $|S\setminus M|=w$. We have $|M|=u+s-w$ and $|M|-|S|=u-w$. As in A10 consider the case $u-w>0$ and $u>0$, then:
    \begin{equation*}
        \tilde{\kappa}_{M,1}=\frac{\psi}{u-w}\ln\bigg[\frac{\binom{p}{|M|}}{\binom{p}{|S|}}\bigg]+ \frac{1}{2}\ln n
    \end{equation*}
    A well-known property of binomial coefficients is that for any positive integers $n$, $h$, $k$ we have:
    \begin{equation}\label{eq:idbinom}
        \binom{n}{h}\binom{n-h}{k}=\binom{n}{h+k}\binom{h+k}{h}
    \end{equation}
    Consider first the case in which $u-w>0$, take $n=p$, $k=u-w$ and $h=s$ in \eqref{eq:idbinom}, we have:
    \begin{equation}\label{eq:identitybinom}
        \frac{\binom{p}{|M|}}{\binom{p}{s}}=
        %\frac{{n \choose h+k}}{{n \choose h}}= \frac{{n -h\choose k}}{{h+k \choose h}} =
        \frac{\binom{p -s}{u-w}}{\binom{|M|}{s}}
    \end{equation}
    % \begin{equation}\label{eq:identitybinom}
    %     \frac{{p \choose |M|}}{{p \choose s}}=\frac{{p \choose p-s+w-u}}{{p \choose s}} =
    %     \frac{{n \choose h+k}}{{n \choose n-h}}=\frac{{n \choose h+k}}{{n \choose h}}= \frac{{n -h\choose k}}{{h+k \choose h}} =
    %     \frac{{p -s \choose u-w}}{{|M| \choose s}}
    % \end{equation}
A standard bounds on binomial coefficient for $1\leq k\leq n$ are:
\begin{equation}\label{eq:boundbinom}
  \left(\frac{n}{k}\right)^k \leq  \binom{n}{k} \leq \frac{n^n}{k^k(n-k)^{n-k}}
\end{equation}
Combining \eqref{eq:identitybinom} and \eqref{eq:boundbinom}, we get:
$$ \frac{\psi}{u-w}\ln\bigg[\frac{\binom{p}{|M|}}{\binom{p}{s}}\bigg] \geq \psi\ln\bigg(\frac{p-s}{u-w}\bigg)-\psi\bigg[\frac{s}{u-w}\ln\bigg(\frac{u-w}{s}+1\bigg) +\ln\bigg(1+\frac{s}{u-w}\bigg) \bigg]$$
We have:
$$\psi\ln\bigg(\frac{p-s}{u-w}\bigg) \geq \psi\ln\bigg(\frac{p-s}{u}\bigg) \geq \ln\bigg(\frac{p-s}{u}\bigg)$$
Observe also that $f:x \mapsto \frac{\ln(x+1)}{x}+\ln(1+{x}^{-1})$ for $x>0$ is decreasing. We have $\frac{u-w}{s} \geq s^{-1}$ and then: 
\begin{equation}\label{eq:lowerboundb1ebic2}
        \tilde{\kappa}_{M,1} \geq \ln\big(\frac{p-s}{u}\big)-s\ln(s^{-1}+1)+ \ln \big(\tfrac{\sqrt{n}}{1+s}\big)
\end{equation}
 We have that $\lim_{n\to+\infty}s\ln(s^{-1}+1)< +\infty$ and since $s=o(\sqrt{n})$, $\ln(\tfrac{\sqrt{n}}{s+1}) \to +\infty$, then $\eta_1$ satisfies assumption~A10 in the case $u-w>0$.\\
 Consider now the case where $u-w<0$. We have:
     \begin{equation*}
        \tilde{\kappa}_{M,1}=\frac{\psi}{w-u}\ln\bigg[\frac{\binom{p}{s}}{\binom{p}{|M|}}\bigg]+ \frac{1}{2}\ln n
    \end{equation*}
 Taking $n=p$, $k=w-u$ and $h=|M|$ in \eqref{eq:idbinom}, we obtain:
     \begin{equation}
        \frac{\binom{p}{s}}{\binom{p}{|M|}}=
        %\frac{{n \choose h+k}}{{n \choose h}}= \frac{{n -h\choose k}}{{h+k \choose h}} =
        \frac{\binom{p - |M|}{w-u}}{\binom{s}{|M|}}
    \end{equation}
By bounds in \eqref{eq:boundbinom}, we get:
$$ \frac{\psi}{w-u}\ln\bigg[\frac{\binom{p}{s}}{\binom{p}{|M|}}\bigg] \geq \psi\ln\bigg(\frac{p-|M|}{w-u}\bigg)+\psi\bigg[\frac{s}{w-u}\ln\bigg(1-\frac{w-u}{s}\bigg) -\ln\bigg(\frac{s}{w-u}-1\bigg) \bigg]$$
Since $|M|<s\leq p/2$ and $\psi\geq 1$ we have:
$$\psi\ln\bigg(\frac{p-|M|}{w-u}\bigg) \geq \psi\ln\bigg(\frac{p-s}{w-u}\bigg) \geq \psi\ln\bigg(\frac{p-s}{w}\bigg)\geq \ln\bigg(\frac{p-s}{u}\bigg)$$
Observe that $g:x \mapsto \frac{\ln(1-x)}{x}-\ln({x}^{-1}-1)$ for $x\in(0,1)$ is increasing. We have $\frac{w-u}{s} \geq s^{-1}$ and then: 
\begin{equation}
        \tilde{\kappa}_{M,1} \geq \ln\big(\frac{p-s}{u}\big)-s\ln(1-s^{-1})+ \ln \big(\tfrac{\sqrt{n}}{s-1}\big)
\end{equation}
Note that the righ-handside above is well defined. Since $u>0$ and $u-w<0$, we have $w\geq2$ and $|M|<s$. The later is not possible unless $s-1>0$. We have that $\lim_{n\to+\infty}s\ln(1-s^{-1})< +\infty$ and since $s=o(\sqrt{n})$, $\ln(\tfrac{\sqrt{n}}{s-1}) \to +\infty$, then $\eta_1$ satisfies assumption~A10 in the case $u-w<0$ too.

 We now show that if XXXX holds then A11 holds. Since A10 holds there exists $b_1$ such that for all $M\in\M$ such that $|M|\neq s$ and $|M\setminus S|>0$, $\tilde{\kappa}_{M, 1} =\;\ln\bigg(\frac{p-s}{|M \setminus S| }\bigg) + b_{1}$. By inequalities XXX and \eqref{eq:lowerboundb1ebic2}, for any of those $M$, $b_1 \geq \tfrac{1}{2}\ln(n)-(s+2)$. By assumption~A11 we have: $1-\gamma=1-\frac{1}{2}\frac{\ln(p-s)}{})$ XXX -MISSING XXX
\end{proof}


\end{document}
